\begin{filecontents*}[overwrite]{note_refs.bib}
@inproceedings{lambert2022variational,
  author    = {Lambert, Marc and Chewi, Sinho and Bach, Francis and Bonnabel, Silv\`ere and Rigollet, Philippe},
  title     = {Variational inference via {W}asserstein gradient flows},
  booktitle = {Advances in Neural Information Processing Systems},
  volume    = {35},
  year      = {2022}
}
@inproceedings{diao2023forward,
  author    = {Diao, Michael and Balasubramanian, Krishnakumar and Chewi, Sinho and Salim, Adil},
  title     = {Forward-backward {G}aussian variational inference via {JKO} in the {B}ures--{W}asserstein space},
  booktitle = {Proceedings of the 40th International Conference on Machine Learning},
  year      = {2023}
}
\end{filecontents*}

\documentclass[11pt]{article}

% ============================================================
% Packages
% ============================================================
\usepackage[margin=1in]{geometry}
\usepackage{amsmath,amssymb,amsthm,mathtools}
\usepackage{bm}
\usepackage{xcolor}
\usepackage[numbers,sort&compress]{natbib}
\usepackage{hyperref}
\usepackage{aliascnt}
\usepackage[nameinlink,capitalize]{cleveref}

\usepackage{comment}

\hypersetup{
    colorlinks=true,
    linkcolor=blue!60!black,
    citecolor=blue!60!black,
    urlcolor=blue!60!black
}

% ============================================================
% Theorem environments
% ============================================================
\theoremstyle{plain}
\newtheorem{theorem}{Theorem}[section]

\newaliascnt{proposition}{theorem}
\newtheorem{proposition}[proposition]{Proposition}
\aliascntresetthe{proposition}

\newaliascnt{lemma}{theorem}
\newtheorem{lemma}[lemma]{Lemma}
\aliascntresetthe{lemma}

\newaliascnt{corollary}{theorem}
\newtheorem{corollary}[corollary]{Corollary}
\aliascntresetthe{corollary}

\theoremstyle{definition}
\newaliascnt{assumption}{theorem}
\newtheorem{assumption}[assumption]{Assumption}
\aliascntresetthe{assumption}

\newaliascnt{definition}{theorem}
\newtheorem{definition}[definition]{Definition}
\aliascntresetthe{definition}

\newaliascnt{example}{theorem}
\newtheorem{example}[example]{Example}
\aliascntresetthe{example}

\theoremstyle{remark}
\newaliascnt{remark}{theorem}
\newtheorem{remark}[remark]{Remark}
\aliascntresetthe{remark}

\crefname{theorem}{Theorem}{Theorems}
\Crefname{theorem}{Theorem}{Theorems}
\crefname{proposition}{Proposition}{Propositions}
\Crefname{proposition}{Proposition}{Propositions}
\crefname{lemma}{Lemma}{Lemmas}
\Crefname{lemma}{Lemma}{Lemmas}
\crefname{corollary}{Corollary}{Corollaries}
\Crefname{corollary}{Corollary}{Corollaries}
\crefname{assumption}{Assumption}{Assumptions}
\Crefname{assumption}{Assumption}{Assumptions}
\crefname{definition}{Definition}{Definitions}
\Crefname{definition}{Definition}{Definitions}
\crefname{example}{Example}{Examples}
\Crefname{example}{Example}{Examples}
\crefname{remark}{Remark}{Remarks}
\Crefname{remark}{Remark}{Remarks}

% ============================================================
% Macros
% ============================================================
\newcommand{\R}{\mathbb{R}}
\newcommand{\N}{\mathcal{N}}
\newcommand{\E}{\mathbb{E}}
\newcommand{\Prob}{\mathbb{P}}
\newcommand{\KL}{\mathrm{KL}}
\newcommand{\Tr}{\operatorname{Tr}}
\newcommand{\SPD}{\mathbb{S}_{++}}
\newcommand{\Sym}{\mathbb{S}}
\newcommand{\Id}{I}
\newcommand{\dd}{\mathrm{d}}
\newcommand{\grad}{\operatorname{grad}}
\newcommand{\op}{\mathrm{op}}
\newcommand{\Frob}{\mathrm{F}}
\newcommand{\diag}{\operatorname{diag}}
\newcommand{\Kdisc}{K^{\mathrm{disc}}}
\newcommand{\post}{\mathrm{post}}
\newcommand{\FR}{\mathrm{FR}}
\newcommand{\Riem}{\mathrm{R}}
\newcommand{\calE}{\mathcal{E}}
\newcommand{\PG}{\mathcal{P}_{\mathrm G}}
\newcommand{\Aspace}{\mathcal{A}}
\newcommand{\Fil}{\mathcal{F}}
\newcommand{\STL}{\mathrm{STL}}
\newcommand{\norm}[1]{\left\lVert #1\right\rVert}
\newcommand{\inner}[2]{\left\langle #1,#2\right\rangle}
\newcommand{\logp}{\log_+}
\newcommand{\glob}{\mathrm{glob}}
\newcommand{\cov}{\mathrm{cov}}
\newcommand{\loc}{\mathrm{loc}}
\newcommand{\Gradnorm}{\mathsf{G}}
\newcommand{\Ureg}{\mathcal{U}}
\newcommand{\Hreg}{\mathcal{H}}
\newcommand{\gstar}{\gamma_\star}
\newcommand{\gund}{\underline{\gamma}}
\newcommand{\Lstar}{L_\star}
\newcommand{\sharpthr}{\Delta^{\sharp}}
\newcommand{\alphastar}{\alpha_\star}
\newcommand{\betastar}{\beta_\star}
\newcommand{\kappastar}{\kappa_\star}

% ============================================================
% Title
% ============================================================
\title{Sharp Global Rates, Covariance Bootstrap, and Gaussian-Core Local Convergence
for Gaussian Natural Gradient Variational Inference}
\author{}
\date{17 July 2026}

\begin{document}
\maketitle

\begin{abstract}
We sharpen the deterministic convergence theory of Gaussian natural-gradient variational
inference in both its global and local phases.  The continuous-time analysis is formulated in
coordinates whitened by the optimizing Gaussian covariance.  If
$\alphastar\Id\preceq\nabla^2\widehat V\preceq\betastar\Id$,
$\kappastar:=\betastar/\alphastar$, and
$\lambda_{0,\star}:=\lambda_{\min}(C_\star^{-1/2}C_0C_\star^{-1/2})$, then the exact energy estimate
\[
    \Delta(a_t)
    \leq
    \Delta_0\bigl[1+\betastar\lambda_{0,\star}(e^t-1)\bigr]^{-1/\kappastar}
\]
holds globally.  Consequently, the global-localization time is
$O(\logp(1/(\betastar\lambda_{0,\star}))+\kappastar\logp(\Delta_0/\delta_{\mathrm{loc}}))$.
Both $\kappastar$ and $\betastar\lambda_{0,\star}$ are invariant under invertible affine
reparametrizations.  We also give an explicit two-dimensional logarithmic-spiral family with an
admissible anisotropic spectral initialization for which
$K\leq\kappastar\leq1.03K$ and the exact Gaussian Fisher--Rao trajectory satisfies
$T_{\mathrm{loc}}\geq K/100\geq\kappastar/103$.  Thus the linear dependence on the intrinsic
affine-invariant condition number is sharp, up to logarithmic input-size factors, for the general
spectral initialization class.

For the deterministic discretizations, an explicit covariance bootstrap improves the dependence
on an underdispersed initial covariance from inverse to logarithmic and organizes the complexity
into covariance burn-in, global localization, and local convergence.  A short
Bures--Wasserstein warm-start can remove the covariance burn-in before a permanent switch to
Fisher--Rao dynamics.  We complement the upper bounds with a one-dimensional smooth bump-train
construction.  For every fixed $h=\gamma/\kappa$, $0<\gamma\leq1/2$, it has
$1\leq V_\kappa''\leq\kappa$, a dimension-free bound $\|V_\kappa'''\|_\infty=O(1)$, and an
initial covariance $c_0=1/\kappa$ already on the curvature scale, yet both the Riemannian and
KL/Bregman schemes require $\Omega(\kappa^2/\gamma)=\Omega(\kappa/h)$ iterations for a
constant-factor energy reduction.  Thus the polynomial fixed-step global-localization factor is
sharp for both schemes, even under Hessian Lipschitzness.  The local phase is analyzed through a
Gaussian-core decomposition $F=F_{\mathrm G}+H$, retaining the exact nonlinear Gaussian field
$F_{\mathrm G}(m,C)=(-Cm,C-C^2)$ and Taylor-bounding only the genuinely non-Gaussian remainder.
The resulting continuous-time energy-sublevel region is sharp for Gaussian targets, while the two
deterministic schemes admit explicit sufficient local regions and optimizer-free residual entry
tests.
\end{abstract}

\tableofcontents

\section{Purpose of the note}

This note records five complementary refinements of the convergence analysis for the Gaussian
natural-gradient flow and its two deterministic discretizations: the Riemannian retraction scheme
and the KL/Bregman scheme.

\paragraph{Affine-invariant continuous-time rate and sharpness.}
The natural intrinsic conditioning is obtained after whitening by the optimizing Gaussian
covariance.  In these coordinates the optimizer is $(0,\Id)$ and the curvature condition number is
$\kappastar=\betastar/\alphastar$.  We prove the global estimate
\[
    \Delta(a_t)
    \leq
    \Delta_0\bigl[1+\betastar\lambda_{0,\star}(e^t-1)\bigr]^{-1/\kappastar},
\]
which yields a global-localization time
\[
    \logp\frac1{\betastar\lambda_{0,\star}}
    +\kappastar\logp\frac{\Delta_0}{\delta_{\mathrm{loc}}}.
\]
The quantities $\kappastar$ and $\betastar\lambda_{0,\star}$ are invariant under every invertible
affine change of variables.  An explicit logarithmic-spiral potential gives a matching
$\Omega(\kappastar)$ delay in dimension two under the manuscript's general spectral
initialization.  The lower construction is deliberately stated with its limitations: it uses an
anisotropic initial covariance and does not settle the stricter initialization
$C_0=\beta^{-1}\Id$.

\paragraph{Deterministic refinement: covariance bootstrap.}
The previous global proofs use the fixed lower covariance bound
\[
    C_t\succeq \lambda_{\min}\Id,
    \qquad
    \lambda_{\min}=\min\left\{\lambda_{0,\min},\frac1\beta\right\},
\]
for the entire global phase.  This gives a valid but pessimistic warm-up time proportional to
$1/\lambda_{0,\min}$.  When the covariance is too small it grows; tracking this growth through an
explicit covariance bootstrap improves the initialization dependence from inverse to logarithmic,
\[
    \frac1{\lambda_{0,\min}}
    \quad\leadsto\quad
    \logp\frac1{\lambda_{0,\min}}.
\]
The deterministic complexity is thereby organized into
\[
    \boxed{
    \text{covariance burn-in}
    \; + \;
    \text{global localization}
    \; + \;
    \text{local convergence}.}
\]

\paragraph{Sharpness of the fixed-step discretization rates.}
The global upper bounds contain the factor $\kappa/\Delta t$.  We prove that this factor is sharp
for both explicit discretizations at constant-factor accuracy.  More precisely, for every fixed
$\Delta t=\gamma/\kappa$, $0<\gamma\leq1/2$, a one-dimensional even $C^\infty$ bump-train
potential with $1\leq V''\leq\kappa$ and $\|V'''\|_\infty=O(1)$ keeps the covariance matched at
$c_n\simeq1/\kappa$ while moving the mean by only a relative
$\Theta(\gamma/\kappa^2)$ per step.  Both the Riemannian and KL/Bregman schemes therefore need
$\Omega(\kappa^2/\gamma)=\Omega(\kappa/\Delta t)$ iterations for a constant-factor energy
reduction.  Since the example starts at $c_0=1/\kappa$, the lower bound is a genuine global
localization obstruction and not a covariance burn-in effect.  The claim is deliberately scoped
to one globally fixed step; adaptive relative-curvature, line-search, and implicit schemes are not
covered.

\paragraph{Transport-bootstrap refinement: Wasserstein warm-start.}
The logarithmic covariance burn-in is intrinsic to pure Fisher--Rao dynamics when the initial
covariance is severely underdispersed.  It can nevertheless be removed algorithmically by using
Bures--Wasserstein transport only as a short warm-start.  In the small-covariance regime,
Fisher--Rao covariance growth is multiplicative, whereas Bures--Wasserstein covariance growth is
additive and scale-independent.  We therefore first lift the covariance to a fixed curvature-scale
floor and then switch permanently to Fisher--Rao for global localization and local convergence.

\paragraph{Local refinement: a Gaussian-sharp region via the Gaussian core.}
The local convergence theory (\Cref{sec:local}) is built on the forward-invariant energy sublevel
$\Ureg_\delta=\{a:\Delta(a)\leq\delta\}$.  On this region the field is split as
$F=F_{\mathrm G}+H$, where the exact Gaussian core
$F_{\mathrm G}(m,C)=(-Cm,C-C^2)$ is dissipated without approximation and only the
non-Gaussian remainder $H$ is Taylor-bounded.  Since $H\equiv0$ for Gaussian targets, the
certified region attains the exact Gaussian threshold in every dimension.  For general strongly
log-concave targets it degrades continuously through a target-adaptive non-Gaussian modulus.
The deterministic schemes are treated through their exact Gaussian maps, and entry into the local
region is certified by an optimizer-free residual test in the original coordinates.

\paragraph{Status of the constants and lower bound.}
Closed-form universal constants and exact Gaussian thresholds are displayed explicitly.
Target-adaptive quantities such as $\alphastar,\betastar$, the linearized-gap constant $\Gamma$,
and the Hessian-dispersion modulus are defined by explicit variational or scalar conditions.
The spiral lower bound is constructive and its shadowing argument is isolated in a technical lemma
with explicit numerical certificates.  The fixed-step discrete lower bound is likewise constructive;
its flat-top bump train is even, strongly convex, globally smooth, and uniformly Hessian-Lipschitz.
The continuous lower bound proves sharpness of the intrinsic $\kappa_\star$ dependence, while the
discrete lower bound proves sharpness of the factor $\kappa/\Delta t$ for constant-factor global
localization.  A finite worst-case localization statement must charge the initial energy, initial
radius, or an equivalent input-size parameter: without such a restriction, the worst-case
localization time is infinite even for a quadratic target.

\section{Setup and notation}

We consider the Gaussian variational family
\[
    \PG
    :=
    \left\{
        \rho_a=\N(m,C):
        a=(m,C)\in\Aspace:=\R^{N_\theta}\times\SPD(N_\theta)
    \right\},
\]
and the objective
\[
    \calE(a)
    :=
    \KL(\rho_a\,\|\,\rho_{\post}),
    \qquad
    \rho_{\post}(\theta)\propto \exp(-V(\theta)),
    \qquad
    a_\star=(m_\star,C_\star)
    \in
    \operatorname*{argmin}_{a\in\Aspace}\calE(a).
\]
We write the energy gap as
\[
    \Delta(a):=\calE(a)-\calE(a_\star),
    \qquad
    \Delta_0:=\Delta(a_0),
    \qquad
    \Delta_n:=\Delta(a_n).
\]
The stepsize is denoted $\Delta t>0$, and $N_\theta$ is the parameter dimension. Throughout,
$\Delta t$ always appears with the differential $t$ attached, so there is no clash with the
energy gap $\Delta(\cdot)$.

\begin{assumption}[Strong log-concavity and smoothness]
\label{assump:logconcave-smooth}
There exist constants $0<\alpha\leq \beta<\infty$ such that
\[
    \alpha\Id
    \preceq
    \nabla^2V(\theta)
    \preceq
    \beta\Id,
    \qquad
    \forall \theta\in\R^{N_\theta}.
\]
Equivalently, $\alpha\Id\preceq -\nabla_\theta\nabla_\theta\log\rho_{\post}(\theta)\preceq\beta\Id$.
We write the condition number
\[
    \kappa:=\frac{\beta}{\alpha}\geq 1.
\]
\end{assumption}

For a Gaussian parameter $a=(m,C)$ we abbreviate
\[
    g(a):=\E_{\rho_a}[\nabla_\theta\log\rho_{\post}(X)]
    =-\E_{\rho_a}[\nabla V(X)],
    \qquad
    A(a):=\E_{\rho_a}[\nabla^2V(X)],
\]
and the \emph{precision residual}
\[
    R(a):=C^{-1}-A(a).
\]
\Cref{assump:logconcave-smooth} gives $\alpha\Id\preceq A(a)\preceq\beta\Id$. At the
minimizer, the first-order conditions read
\begin{equation}
\label{eq:first-order}
    g(a_\star)=0,
    \qquad
    A(a_\star)=C_\star^{-1},
    \qquad\text{i.e.}\qquad
    R(a_\star)=0 .
\end{equation}

\subsection{Continuous flow and deterministic schemes}

The Gaussian natural gradient flow is
\begin{equation}
\label{eq:NG-flow}
    \frac{\dd m_t}{\dd t}
    =
    C_tg(a_t),
    \qquad
    \frac{\dd C_t}{\dd t}
    =
    C_t-C_tA(a_t)C_t.
\end{equation}
The deterministic Riemannian retraction scheme is
\begin{equation}
\label{eq:Riem-det}
    m_{n+1}=m_n+\Delta t\,C_ng(a_n),
    \qquad
    C_{n+1}
    =
    C_n^{1/2}
    \exp\bigl(\Delta t(\Id-C_n^{1/2}A(a_n)C_n^{1/2})\bigr)
    C_n^{1/2}.
\end{equation}
The deterministic KL/Bregman scheme is
\begin{equation}
\label{eq:KL-det}
    m_{n+1}=m_n+\Delta t\,C_ng(a_n),
    \qquad
    C_{n+1}
    =(1+\Delta t)\bigl(C_n^{-1}+\Delta t\,A(a_n)\bigr)^{-1},
\end{equation}
or equivalently $C_{n+1}^{-1}=\tfrac1{1+\Delta t}\bigl(C_n^{-1}+\Delta t\,A(a_n)\bigr)$.
We use the energy splitting $\calE=\calE_1+\calE_2$ with
\[
    \calE_1(a)=\E_{\rho_a}[\log\rho_a]=-\tfrac12\log\det C+\text{const.},
    \qquad
    \calE_2(a)=-\E_{\rho_a}[\log\rho_{\post}].
\]

\subsection{Single-sample stochastic schemes}

Let
\[
    \Fil_n:=\sigma(a_0,\ldots,a_n),
    \qquad
    X_n\mid\Fil_n\sim\N(m_n,C_n),
    \qquad
    H_n:=\nabla^2V(X_n).
\]
The \emph{raw} one-sample estimators are
$b_n=\nabla\log\rho_{\post}(X_n)=-\nabla V(X_n)$ and $S_n=-H_n$, and the raw stochastic
schemes replace $g(a_n)$ and $A(a_n)$ by $b_n$ and $H_n$ in \eqref{eq:Riem-det} and
\eqref{eq:KL-det}. Alternatively, one uses the sticking-the-landing (STL) estimators. The
STL score estimator is
\begin{equation}
\label{eq:STL-score}
    \widehat b_n
    :=
    \nabla\log\rho_{\post}(X_n)-\nabla\log\rho_{a_n}(X_n)
    =
    -\nabla V(X_n)+C_n^{-1}(X_n-m_n),
\end{equation}
and the STL residual-Hessian estimator is
\begin{equation}
\label{eq:STL-residual}
    \widehat K_n
    :=
    C_n^{-1}+\nabla^2\log\rho_{\post}(X_n)
    =
    C_n^{-1}-H_n.
\end{equation}
They are conditionally unbiased for the deterministic gradient and residual:
\begin{equation}
\label{eq:STL-unbiased}
    \E[\widehat b_n\mid\Fil_n]=g(a_n),
    \qquad
    \E[\widehat K_n\mid\Fil_n]=R(a_n).
\end{equation}
The first identity uses $\E_{\rho_{a_n}}[C_n^{-1}(X_n-m_n)]=0$; the second uses the
definition of $A(a_n)$. The \emph{STL Riemannian scheme} is
\begin{equation}
\label{eq:Riem-STL}
    m_{n+1}=m_n+\Delta t\,C_n\widehat b_n,
    \qquad
    C_{n+1}
    =
    C_n^{1/2}\exp\bigl(\Delta t\,C_n^{1/2}\widehat K_nC_n^{1/2}\bigr)C_n^{1/2},
\end{equation}
and the \emph{STL KL scheme} is
\begin{equation}
\label{eq:KL-STL}
    m_{n+1}=m_n+\Delta t\,C_n\widehat b_n,
    \qquad
    C_{n+1}
    =
    (1+\Delta t)\bigl((1+\Delta t)C_n^{-1}-\Delta t\,\widehat K_n\bigr)^{-1}.
\end{equation}
Because $\widehat K_n=C_n^{-1}-H_n$, the covariance updates in \eqref{eq:Riem-STL} and
\eqref{eq:KL-STL} simplify to
\begin{equation}
\label{eq:STL-cov-simplify}
    C_{n+1}
    =
    C_n^{1/2}\exp\bigl(\Delta t(\Id-C_n^{1/2}H_nC_n^{1/2})\bigr)C_n^{1/2},
    \qquad
    C_{n+1}
    =
    (1+\Delta t)\bigl(C_n^{-1}+\Delta t\,H_n\bigr)^{-1},
\end{equation}
respectively. Thus the STL covariance updates coincide with the raw Hessian-sampling
updates. In particular, since $\alpha\Id\preceq H_n\preceq\beta\Id$ almost surely, the
covariance bounds proved below hold pathwise for the stochastic schemes.

\subsection{Initial covariance and shorthands}

We assume the initial covariance satisfies
\[
    \lambda_{0,\min}\Id\preceq C_0\preceq \lambda_{0,\max}\Id,
\]
and we set, as in the global theory,
\[
    \lambda_0:=\lambda_{0,\min},
    \qquad
    \lambda_{\max}:=
    \max\left\{
        \lambda_{0,\max},\frac1\alpha
    \right\},
    \qquad
    \logp x:=\max\{\log x,0\}.
\]

\subsection{Fisher--Rao metric, gradient norm, and dissipation identity}

The Fisher--Rao metric of the Gaussian family, at $a=(m,C)$ and for tangent vectors
$(u,U),(v,W)\in\R^{N_\theta}\times\Sym(N_\theta)$, is
\begin{equation}
\label{eq:tangent-norm}
    \inner{(u,U)}{(v,W)}_a
    :=
    u^\top C^{-1}v
    +
    \tfrac12\Tr\bigl(C^{-1}UC^{-1}W\bigr),
    \qquad
    \norm{(u,U)}_a^2
    :=
    u^\top C^{-1}u+\tfrac12\norm{C^{-1/2}UC^{-1/2}}_{\Frob}^2 .
\end{equation}
This is the Fisher information inner product of the Gaussian family. The induced squared
gradient norm of the objective is
\begin{equation}
\label{eq:G-norm}
    \Gradnorm(a)^2
    :=
    g(a)^\top Cg(a)
    +
    \tfrac12\norm{C^{1/2}R(a)C^{1/2}}_{\Frob}^2 .
\end{equation}

\begin{lemma}[Metric identification and exact dissipation identity]
\label{lem:metric-dissipation}
Under \Cref{assump:logconcave-smooth}, the differential of $\calE$ at $a=(m,C)$ in the
direction $(u,U)$ is
\[
    D\calE(a)[(u,U)]
    =
    -g(a)^\top u
    +
    \tfrac12\Tr\bigl((A(a)-C^{-1})U\bigr).
\]
Consequently, the Riemannian gradient of $\calE$ with respect to
\eqref{eq:tangent-norm} is
\[
    \grad\calE(a)
    =
    \bigl(-Cg(a),\;CA(a)C-C\bigr),
\]
the flow \eqref{eq:NG-flow} is exactly the Fisher--Rao gradient flow
$\dot a_t=-\grad\calE(a_t)$, its squared gradient norm is
$\norm{\grad\calE(a)}_a^2=\Gradnorm(a)^2$ as in \eqref{eq:G-norm}, and along
\eqref{eq:NG-flow},
\[
    \frac{\dd}{\dd t}\Delta(a_t)
    =
    -\Gradnorm(a_t)^2 .
\]
\end{lemma}

\begin{proof}
\emph{Differential.} Write $\calE=\calE_1+\calE_2$ with
$\calE_1(a)=-\tfrac12\log\det C+\mathrm{const}$ and
$\calE_2(a)=\E_{\rho_a}[V]+\mathrm{const}$. For the entropy part,
$\partial_m\calE_1=0$ and
$\partial_C\bigl(-\tfrac12\log\det C\bigr)=-\tfrac12C^{-1}$. For the potential part,
writing $X=m+C^{1/2}Z$ with $Z\sim\N(0,\Id)$, differentiation under the integral sign
(justified by dominated convergence, since $\norm{\nabla^2V}_{\op}\leq\beta$ implies
$\norm{\nabla V(x)}\leq\norm{\nabla V(0)}+\beta\norm{x}$, so all integrands are bounded
by polynomials times Gaussian densities) gives
\[
    \partial_m\E_{\rho_a}[V]=\E_{\rho_a}[\nabla V]=-g(a),
    \qquad
    \partial_C\E_{\rho_a}[V]=\tfrac12\E_{\rho_a}[\nabla^2V]=\tfrac12A(a),
\]
the second identity being Price's theorem (Gaussian integration by parts applied twice).
Adding the two parts,
\[
    D\calE(a)[(u,U)]
    =
    -g(a)^\top u
    +
    \Tr\Bigl(\bigl(\tfrac12A(a)-\tfrac12C^{-1}\bigr)U\Bigr),
\]
which is the stated formula.

\emph{Gradient.} The gradient $(\widehat u,\widehat U):=\grad\calE(a)$ is defined by
$\inner{(\widehat u,\widehat U)}{(u,U)}_a=D\calE(a)[(u,U)]$ for all $(u,U)$. Matching the
mean blocks gives $C^{-1}\widehat u=-g(a)$, i.e.\ $\widehat u=-Cg(a)$. Matching the
covariance blocks gives
$\tfrac12\Tr(C^{-1}\widehat UC^{-1}U)=\tfrac12\Tr((A(a)-C^{-1})U)$ for all symmetric $U$,
i.e.\ $C^{-1}\widehat UC^{-1}=A(a)-C^{-1}$, i.e.\ $\widehat U=CA(a)C-C$. Therefore
$-\grad\calE(a)=(Cg(a),\,C-CA(a)C)$, which is the right-hand side of \eqref{eq:NG-flow}.

\emph{Gradient norm.} By \eqref{eq:tangent-norm},
\[
    \norm{\grad\calE(a)}_a^2
    =
    g(a)^\top CC^{-1}Cg(a)
    +
    \tfrac12\Tr\bigl(C^{-1}(CA(a)C-C)C^{-1}(CA(a)C-C)\bigr).
\]
The first term is $g(a)^\top Cg(a)$. For the second, $C^{-1}(CA(a)C-C)=A(a)C-\Id$, so
\[
    \Tr\bigl((A(a)C-\Id)(A(a)C-\Id)\bigr)
    =
    \Tr\bigl(C(A(a)-C^{-1})C(A(a)-C^{-1})\bigr)
    =
    \norm{C^{1/2}R(a)C^{1/2}}_{\Frob}^2,
\]
where the last equality uses cyclicity of the trace and $R(a)=C^{-1}-A(a)$ (the sign
cancels in the square). This gives \eqref{eq:G-norm}.

\emph{Dissipation.} By the chain rule and the gradient identification,
\[
    \frac{\dd}{\dd t}\Delta(a_t)
    =
    D\calE(a_t)[\dot a_t]
    =
    D\calE(a_t)\bigl[-\grad\calE(a_t)\bigr]
    =
    -\inner{\grad\calE(a_t)}{\grad\calE(a_t)}_{a_t}
    =
    -\Gradnorm(a_t)^2 .
    \qedhere
\]
\end{proof}

\section{Two basic inequalities}

\begin{lemma}[Wasserstein PL inequality on Gaussians]
\label{lem:W2-PL}
Under \Cref{assump:logconcave-smooth}, for every Gaussian parameter $a\in\Aspace$,
\[
    \norm{\grad^{W_2}\calE(a)}_{W_2}^2
    \geq
    2\alpha\,\Delta(a).
\]
\end{lemma}

This is the Wasserstein Polyak--{\L}ojasiewicz inequality induced by the
$\alpha$-strong convexity of $V$, restricted to the Gaussian family; see the base
manuscript and \citet{lambert2022variational}.

\begin{lemma}[Fisher--Rao/Wasserstein comparison]
\label{lem:FR-W2}
Suppose $a=(m,C)$ with $C\succeq\ell\Id$ for some $\ell>0$. Then, with \eqref{eq:G-norm},
\eqref{eq:G-norm},
\[
    \Gradnorm(a)^2
    \geq
    \frac{\ell}{2}\,\norm{\grad^{W_2}\calE(a)}_{W_2}^2
    \geq
    \alpha\ell\,\Delta(a).
\]
Moreover,
\[
    \Gradnorm(a)^2
    \leq
    \norm{C}_{\op}\,\norm{\grad^{W_2}\calE(a)}_{W_2}^2 .
\]
\end{lemma}

\begin{proof}
The Wasserstein Gaussian gradient norm is
\[
    \norm{\grad^{W_2}\calE(a)}_{W_2}^2
    =
    \norm{g(a)}^2+\Tr\bigl(R(a)CR(a)\bigr),
\]
and \eqref{eq:G-norm} reads
\[
    \Gradnorm(a)^2
    =
    g(a)^\top Cg(a)+\tfrac12\Tr\bigl(CR(a)CR(a)\bigr),
\]
where we used $\norm{C^{1/2}R(a)C^{1/2}}_{\Frob}^2=\Tr(CR(a)CR(a))$. If $C\succeq\ell\Id$,
then $g(a)^\top Cg(a)\geq\ell\norm{g(a)}^2\geq\tfrac{\ell}{2}\norm{g(a)}^2$, and since
$R(a)CR(a)\succeq0$,
\[
    \tfrac12\Tr\bigl(CR(a)CR(a)\bigr)
    =
    \tfrac12\Tr\bigl(C\,[R(a)CR(a)]\bigr)
    \geq
    \frac{\ell}{2}\,\Tr\bigl(R(a)CR(a)\bigr).
\]
Adding the two estimates gives
$\Gradnorm(a)^2\geq\tfrac{\ell}{2}\norm{\grad^{W_2}\calE(a)}_{W_2}^2$, and combining with
\Cref{lem:W2-PL} gives
$\Gradnorm(a)^2\geq\tfrac{\ell}{2}\cdot2\alpha\Delta(a)=\alpha\ell\Delta(a)$. For the
upper bound, with $L_C:=\norm{C}_{\op}$,
\[
    \begin{aligned}
    g(a)^\top Cg(a)&\leq L_C\norm{g(a)}^2,\\
    \tfrac12\Tr\bigl(CR(a)CR(a)\bigr)
    &=\tfrac12\Tr\bigl(C\,[R(a)CR(a)]\bigr)
    \leq \tfrac{L_C}2\Tr\bigl(R(a)CR(a)\bigr)
    \leq L_C\Tr\bigl(R(a)CR(a)\bigr),
    \end{aligned}
\]
and adding gives $\Gradnorm(a)^2\leq L_C\norm{\grad^{W_2}\calE(a)}_{W_2}^2$.
\end{proof}

\section{Affine-invariant continuous-time global convergence}
\label{sec:cont}

The global continuous-time estimate is most naturally stated after whitening by the optimizing
Gaussian covariance.  This removes purely affine anisotropy and exposes the intrinsic curvature
parameter that remains after the optimal Gaussian geometry has been normalized.

\subsection{Optimizer whitening and intrinsic conditioning}

Let $a_\star=(m_\star,C_\star)$ be the unique optimizing Gaussian and set
\begin{equation}
\label{eq:optimizer-whitening}
    \widehat m:=C_\star^{-1/2}(m-m_\star),
    \qquad
    \widehat C:=C_\star^{-1/2}CC_\star^{-1/2},
    \qquad
    \widehat V(z):=V(m_\star+C_\star^{1/2}z).
\end{equation}
The optimizer becomes $(\widehat m_\star,\widehat C_\star)=(0,\Id)$.  Define the optimal
whitened curvature constants
\begin{equation}
\label{eq:star-curvature}
    \alphastar
    :=\inf_{z\in\R^{N_\theta}}\lambda_{\min}(\nabla^2\widehat V(z)),
    \qquad
    \betastar
    :=\sup_{z\in\R^{N_\theta}}\lambda_{\max}(\nabla^2\widehat V(z)),
    \qquad
    \kappastar:=\frac{\betastar}{\alphastar},
\end{equation}
and the whitened initial covariance floor
\begin{equation}
\label{eq:lambda0-star}
    \lambda_{0,\star}
    :=\lambda_{\min}(\widehat C_0)
    =\lambda_{\min}(C_\star^{-1/2}C_0C_\star^{-1/2}).
\end{equation}

\begin{lemma}[Optimizer whitening, flow covariance, and affine invariance]
\label{lem:optimizer-whitening}
Under \Cref{assump:logconcave-smooth}, the transformed parameters solve the same Gaussian
Fisher--Rao system,
\begin{equation}
\label{eq:whitened-flow}
    \dot{\widehat m}_t=\widehat C_t\widehat g(\widehat a_t),
    \qquad
    \dot{\widehat C}_t
    =\widehat C_t-\widehat C_t\widehat A(\widehat a_t)\widehat C_t,
\end{equation}
where $\widehat g=-\E[\nabla\widehat V]$ and
$\widehat A=\E[\nabla^2\widehat V]$.  Moreover,
\[
    \E_{Z\sim\N(0,\Id)}[\nabla\widehat V(Z)]=0,
    \qquad
    \E_{Z\sim\N(0,\Id)}[\nabla^2\widehat V(Z)]=\Id,
\]
so $\alphastar\leq1\leq\betastar$.  The quantities $\alphastar,\betastar,\kappastar$ and the
spectrum of $\widehat C_0$ are invariant under arbitrary invertible affine changes of the original
coordinates.
\end{lemma}

\begin{proof}
Write $S=C_\star^{1/2}$.  Since
$\nabla\widehat V(z)=S\nabla V(m_\star+Sz)$ and
$\nabla^2\widehat V(z)=S\nabla^2V(m_\star+Sz)S$, one has
$\widehat g=Sg$ and $\widehat A=SAS$.  Hence
\[
    \dot{\widehat m}
    =S^{-1}\dot m
    =S^{-1}Cg
    =\widehat C\widehat g,
\]
and
\[
    \dot{\widehat C}
    =S^{-1}\dot CS^{-1}
    =\widehat C-\widehat C\widehat A\widehat C.
\]
The first-order conditions \eqref{eq:first-order} give the two expectation identities at the
whitened optimizer.  In particular, the average whitened Hessian is $\Id$, so its pointwise lower
and upper spectral bounds satisfy $\alphastar\leq1\leq\betastar$.

Now apply an invertible affine map $y=Bx+b$.  The transformed optimizer covariance is
$\widetilde C_\star=BC_\star B^\top$.  The matrix
\[
    O:=\widetilde C_\star^{-1/2}BC_\star^{1/2}
\]
is orthogonal, because $OO^\top=\Id$.  Thus optimizer whitening before and after the affine
change differs only by the orthogonal map $O$.  The whitened Hessians are orthogonally congruent,
and the whitened initial covariances are orthogonally similar.  Their spectral data are therefore
unchanged.
\end{proof}

\begin{lemma}[Relation between intrinsic and original condition numbers]
\label{lem:kappa-star-kappa}
Under \Cref{assump:logconcave-smooth},
\begin{equation}
\label{eq:Cstar-band}
    \frac1\beta\Id\preceq C_\star\preceq\frac1\alpha\Id,
\end{equation}
and
\begin{equation}
\label{eq:kappa-star-kappa}
    \frac1\kappa\Id\preceq\nabla^2\widehat V(z)\preceq\kappa\Id,
    \qquad
    1\leq\kappastar\leq\kappa^2.
\end{equation}
The upper comparison can be very loose: for an anisotropic quadratic target, $\kappastar=1$ even
when $\kappa$ is arbitrarily large.
\end{lemma}

\begin{proof}
The optimizer condition $C_\star^{-1}=A(a_\star)$ and
$\alpha\Id\preceq A(a_\star)\preceq\beta\Id$ imply \eqref{eq:Cstar-band}.  Therefore
\[
    \nabla^2\widehat V(z)
    =C_\star^{1/2}\nabla^2V(m_\star+C_\star^{1/2}z)C_\star^{1/2}
    \succeq\alpha C_\star\succeq\frac1\kappa\Id,
\]
and similarly $\nabla^2\widehat V(z)\preceq\kappa\Id$.  Thus
$\alphastar\geq1/\kappa$, $\betastar\leq\kappa$, and
$\kappastar\leq\kappa^2$; the lower bound $\kappastar\geq1$ follows from
$\alphastar\leq1\leq\betastar$.  If $V(x)=\tfrac12x^\top Hx$, then
$C_\star=H^{-1}$ and $\widehat V(z)=\tfrac12\norm{z}^2$, so $\kappastar=1$.
\end{proof}

\subsection{Precision bootstrap and exact energy estimate}

\begin{lemma}[Precision Duhamel formula]
\label{lem:precision-duhamel}
Let $P_t:=C_t^{-1}$.  Along \eqref{eq:NG-flow},
\begin{equation}
\label{eq:precision-duhamel}
    \dot P_t=A(a_t)-P_t,
    \qquad
    P_t=e^{-t}P_0+\int_0^te^{-(t-s)}A(a_s)\,\dd s.
\end{equation}
The identical formula holds for
$\widehat P_t:=\widehat C_t^{-1}$ with $A$ replaced by $\widehat A$.
\end{lemma}

\begin{proof}
Since $P_t=C_t^{-1}$,
\[
    \dot P_t=-P_t\dot C_tP_t
    =-P_t(C_t-C_tA(a_t)C_t)P_t
    =A(a_t)-P_t.
\]
Multiplication by the integrating factor $e^t$ gives \eqref{eq:precision-duhamel}.  The whitened
formula follows from \eqref{eq:whitened-flow}.
\end{proof}

\begin{lemma}[Continuous covariance bootstrap in both coordinate systems]
\label{lem:cont-cov}
Along \eqref{eq:NG-flow},
\begin{equation}
\label{eq:ell-original}
    C_t\succeq\ell(t)\Id,
    \qquad
    \ell(t):=\frac{\lambda_0e^t}{1+\beta\lambda_0(e^t-1)},
    \qquad
    C_t\preceq\lambda_{\max}\Id.
\end{equation}
In optimizer-whitened coordinates,
\begin{equation}
\label{eq:ell-star}
    \widehat C_t\succeq\ell_\star(t)\Id,
    \qquad
    \ell_\star(t)
    :=\frac{\lambda_{0,\star}e^t}{1+\betastar\lambda_{0,\star}(e^t-1)}.
\end{equation}
In particular,
$C_t\succeq(2\beta)^{-1}\Id$ for
$t\geq\logp(1/(\beta\lambda_0))$, and
$\widehat C_t\succeq(2\betastar)^{-1}\Id$ for
$t\geq\logp(1/(\betastar\lambda_{0,\star}))$.
\end{lemma}

\begin{proof}
By \Cref{lem:precision-duhamel},
\[
    P_t\preceq
    \bigl(e^{-t}\lambda_0^{-1}+\beta(1-e^{-t})\bigr)\Id,
\]
which gives the lower bound in \eqref{eq:ell-original} after inversion.  The upper bound follows
from $A(a_t)\succeq\alpha\Id$: no eigenvalue of $C_t$ can cross $1/\alpha$ from below, and any
eigenvalue above $1/\alpha$ is pushed downward.  The whitened statement follows by the same
argument with $\lambda_{0,\star}$ and $\betastar$.  The two curvature-scale consequences follow by
requiring the decaying initial-precision term to be no larger than the corresponding curvature
upper bound.
\end{proof}

\begin{theorem}[Affine-invariant continuous-time global convergence]
\label{thm:cont-global}
Under \Cref{assump:logconcave-smooth}, for every $t\geq0$,
\begin{equation}
\label{eq:cont-star-decay}
    \boxed{
    \Delta(a_t)
    \leq
    \Delta_0
    \bigl[1+\betastar\lambda_{0,\star}(e^t-1)\bigr]^{-1/\kappastar}.}
\end{equation}
Consequently, for every $\delta\in(0,\Delta_0)$,
\begin{equation}
\label{eq:cont-star-hitting-exact}
    T_{\mathrm c}^{\glob}(\delta)
    :=\inf\{t\geq0:\Delta(a_t)\leq\delta\}
    \leq
    \log\left(
        1+
        \frac{(\Delta_0/\delta)^{\kappastar}-1}
             {\betastar\lambda_{0,\star}}
    \right),
\end{equation}
and, in particular,
\begin{equation}
\label{eq:cont-star-hitting-simple}
    \boxed{
    T_{\mathrm c}^{\glob}(\delta)
    \leq
    \logp\frac1{\betastar\lambda_{0,\star}}
    +\kappastar\log\frac{\Delta_0}{\delta}.}
\end{equation}
The original-coordinate estimate remains valid as well; hence one may take the minimum of
\eqref{eq:cont-star-hitting-exact} and
$\log(1+((\Delta_0/\delta)^\kappa-1)/(\beta\lambda_0))$.
\end{theorem}

\begin{proof}
The energy and the Fisher--Rao dissipation are invariant under the optimizer whitening of
\eqref{eq:optimizer-whitening}.  We therefore work in the whitened coordinates and omit hats.
By \Cref{lem:metric-dissipation},
$\dot\Delta_t=-\Gradnorm(a_t)^2$.  The covariance lower bound
$C_t\succeq\ell_\star(t)\Id$ and \Cref{lem:FR-W2} give
\[
    \Gradnorm(a_t)^2
    \geq
    \alphastar\ell_\star(t)\Delta(a_t),
\]
so
\[
    \dot\Delta_t
    \leq
    -\alphastar\ell_\star(t)\Delta_t.
\]
Gr\"onwall's inequality yields
\[
    \Delta_t
    \leq
    \Delta_0
    \exp\left(-\alphastar\int_0^t\ell_\star(s)\,\dd s\right).
\]
The scalar integral is exact:
\[
    \int_0^t\ell_\star(s)\,\dd s
    =\frac1{\betastar}
      \log\bigl[1+\betastar\lambda_{0,\star}(e^t-1)\bigr].
\]
Since $\alphastar/\betastar=1/\kappastar$, this proves
\eqref{eq:cont-star-decay}.  Solving the equality in \eqref{eq:cont-star-decay} gives
\eqref{eq:cont-star-hitting-exact}.

For \eqref{eq:cont-star-hitting-simple}, set
$R=(\Delta_0/\delta)^{\kappastar}\geq1$ and
$b=\betastar\lambda_{0,\star}$.  If $b\leq1$, then
$1+(R-1)/b\leq R/b$; if $b\geq1$, then
$1+(R-1)/b\leq R$.  Taking logarithms in the two cases gives the displayed bound.  Applying the
same argument without whitening gives the original-coordinate estimate.
\end{proof}

\begin{remark}[What is and is not encoded by $\kappastar$]
The intrinsic condition number removes anisotropy that can be eliminated by an affine
reparametrization.  It does not, by itself, encode the initialization size.  The factor
$\log(\Delta_0/\delta)$ or an equivalent radius/input-size term is indispensable: if arbitrary
initial means are allowed, the worst-case time to reach a fixed energy level is infinite even for a
quadratic target.  The sharpness result below therefore compares the polynomial dependence on
$\kappastar$ while retaining logarithmic initialization factors.
\end{remark}

\section{Sharpness of the continuous-time global-localization rate}
\label{sec:cont-sharp}

We now construct an explicit logarithmic-spiral family showing that the linear dependence on
$\kappastar$ in \Cref{thm:cont-global} cannot be improved, up to logarithmic initialization and
localization factors, under the manuscript's general spectral initialization.

\subsection{The spiral potential}

Fix $K\geq16$ and set
\begin{equation}
\label{eq:spiral-parameters}
    h:=2\alpha,
    \qquad
    c:=\frac{\sqrt{K-2}}3,
    \qquad
    \delta_K:=\frac{10^{-4}}K,
    \qquad
    D_K:=\frac{18h(K-1)}{K+7},
    \qquad
    \beta:=\alpha(2K+1).
\end{equation}
Thus the original-coordinate condition number is $\kappa=2K+1$.

Define
\[
    s(u):=35u^4-84u^5+70u^6-20u^7,
    \qquad 0\leq u\leq1,
\]
and the even $C^3$ cutoff
\begin{equation}
\label{eq:phase-bump}
    \rho(u):=
    \begin{cases}
        1, & |u|\leq\tfrac12,\\[1mm]
        s(2-2|u|), & \tfrac12<|u|<1,\\[1mm]
        0, & |u|\geq1.
    \end{cases}
\end{equation}
Extend $\rho_{\delta_K}(\phi):=\rho(\phi/\delta_K)$ periodically with period $\pi$, and set
\[
    \bar\rho_{\delta_K}
    :=\frac1\pi\int_{-\pi/2}^{\pi/2}\rho_{\delta_K}(\phi)\,\dd\phi,
    \qquad
    g_K(\phi)
    :=\frac{\rho_{\delta_K}(\phi)-\bar\rho_{\delta_K}}
            {1-\bar\rho_{\delta_K}}.
\]
Let $f_K$ be the even $\pi$-periodic solution of
\begin{equation}
\label{eq:fK}
    f_K''=g_K,
    \qquad
    f_K(0)=f_K'(0)=0.
\end{equation}
Then
\begin{equation}
\label{eq:fK-core-bounds}
    f_K(\phi)=\frac12\phi^2
    \quad\text{for }|\phi|\leq\frac{\delta_K}{2},
    \qquad
    \norm{f_K'}_\infty\leq3\delta_K,
    \qquad
    \norm{f_K}_\infty\leq2\pi\delta_K.
\end{equation}

Define the radial cutoff
\begin{equation}
\label{eq:radial-cutoff}
    \chi(\ell):=
    \begin{cases}
        0, & \ell\leq0,\\
        35\ell^4-84\ell^5+70\ell^6-20\ell^7,
            & 0<\ell<1,\\
        1, & \ell\geq1.
    \end{cases}
\end{equation}
For a scale $r_0>0$, use polar coordinates $(r,\theta)$ and set
\[
    \ell:=\log\frac r{r_0},
    \qquad
    \phi:=\theta+c\ell.
\]
The spiral potential is
\begin{equation}
\label{eq:spiral-potential}
    \boxed{
    V_K(r,\theta)
    :=\frac12r^2\bigl[h+D_K\chi(\ell)f_K(\phi)\bigr].}
\end{equation}
For $r\leq r_0$, $V_K(x)=\tfrac h2\norm{x}^2$.  The matching of the first three derivatives of
$\rho$ and $\chi$ at their endpoints gives $V_K\in C^3(\R^2)$; the $\pi$-periodicity of $f_K$
makes $V_K$ even.

\begin{lemma}[Global Hessian certificate]
\label{lem:spiral-Hessian}
The potential \eqref{eq:spiral-potential} satisfies
\begin{equation}
\label{eq:spiral-Hessian-bounds}
    0.98h\Id\preceq\nabla^2V_K(x)\preceq h(K+0.02)\Id
    \qquad\text{for every }x\in\R^2.
\end{equation}
Consequently,
$\alpha\Id\preceq\nabla^2V_K\preceq\beta\Id$ with $\beta=\alpha(2K+1)$.
On the flat part of a spiral valley, the two Hessian eigenvalues are exactly $h$ and $hK$.
\end{lemma}

\begin{proof}
Put $q(\ell,\theta):=\chi(\ell)f_K(\theta+c\ell)$.  For a general profile
$A=A(\ell,\theta)$, differentiation in the orthonormal polar frame gives
\[
    \nabla^2\left(\frac12r^2A\right)
    =
    \begin{pmatrix}
      A+\frac32A_\ell+\frac12A_{\ell\ell}
      &\frac12A_\theta+\frac12A_{\ell\theta}\\[1mm]
      \frac12A_\theta+\frac12A_{\ell\theta}
      &A+\frac12A_\ell+\frac12A_{\theta\theta}
    \end{pmatrix}.
\]
Consequently,
\begin{equation}
\label{eq:spiral-Hessian-decomposition}
 \nabla^2V_K-h\Id
 =\frac{D_K}{2}\chi g_K
     \begin{pmatrix}c^2&c\\c&1\end{pmatrix}
   +D_K\{\chi B_K+E_K\},
\end{equation}
where all functions are evaluated at $(\ell,\phi)$ and
\[
 B_K=
 \begin{pmatrix}
 f_K+\frac32cf_K'&\frac12f_K'\\
 \frac12f_K'&f_K+\frac12cf_K'
 \end{pmatrix},
 \qquad
 E_K=
 \begin{pmatrix}
 (\frac32\chi'+\frac12\chi'')f_K+c\chi'f_K'&\frac12\chi'f_K'\\
 \frac12\chi'f_K'&\frac12\chi'f_K
 \end{pmatrix}.
\]
This identity covers both the uncut region and the transition annulus.

The normalization of the phase cutoff is explicit.  Since
$\int_{-1}^1\rho(u)\,\dd u=3/2$,
\[
 \bar\rho_{\delta_K}=\frac{3\delta_K}{2\pi},
 \qquad
 -\frac{\bar\rho_{\delta_K}}{1-\bar\rho_{\delta_K}}
 \leq g_K\leq1.
\]
For $K\geq16$,
\begin{equation}
\label{eq:spiral-negative-background}
 (K-1)\frac{\bar\rho_{\delta_K}}{1-\bar\rho_{\delta_K}}
 <4.8\cdot10^{-5}.
\end{equation}
Moreover, the polynomial cutoff satisfies
\[
 0\leq\chi\leq1,
 \qquad
 \norm{\chi'}_\infty=\frac{35}{16},
 \qquad
 \norm{\chi''}_\infty<7.514.
\]
Using \eqref{eq:fK-core-bounds}, $D_K/h\leq18$, and
$c\leq\sqrt K/3$, the maximum-row-sum bound in
\eqref{eq:spiral-Hessian-decomposition} gives
\begin{align*}
 \frac{D_K}{h}\norm{\chi B_K+E_K}_{\op}
 &\leq18\Bigl[
   \left(1+\frac32\frac{35}{16}+\frac{7.514}{2}\right)2\pi\delta_K\\
 &\hspace{27mm}
  +\left(\frac32c+\frac12+c\frac{35}{16}+\frac12\frac{35}{16}\right)3\delta_K
  \Bigr]
 <0.0078.
\end{align*}
Every term on the right decreases when $K$ increases, so it is enough to insert $K=16$.

The rank-one matrix in \eqref{eq:spiral-Hessian-decomposition} has nonzero eigenvalue
$1+c^2=(K+7)/9$, and
$\tfrac{D_K}{2}(1+c^2)=h(K-1)$.  Its positive part is therefore at most
$h(K-1)\Id$, while \eqref{eq:spiral-negative-background} bounds its negative part by
$4.8\cdot10^{-5}h\Id$.  Thus, in fact,
\[
 0.992h\Id\preceq\nabla^2V_K
 \preceq h(K+0.0078)\Id,
\]
which implies \eqref{eq:spiral-Hessian-bounds}.  On an uncut valley core,
$g_K=1$ and $f_K(0)=f_K'(0)=0$, so the two eigenvalues are exactly $h$ and $hK$.
Finally, $0.98h=1.96\alpha>\alpha$ and
$h(K+0.02)<\alpha(2K+1)=\beta$.
\end{proof}

\subsection{Reference orbit and Gaussian shadowing}

Let $R_\theta$ denote the planar rotation through angle $\theta$, set
$s_K:=\sqrt{K-2}$, and define
\begin{equation}
\label{eq:QK}
    Q_K
    :=\frac1{K+7}
      \begin{pmatrix}
        (K+1)^2 & 2s_K(K+1)\\
        2s_K(K+1) & 6(K+1)
      \end{pmatrix}.
\end{equation}
A principal-minor calculation gives
\begin{equation}
\label{eq:QK-band}
    \frac32\Id\prec Q_K\prec K\Id,
    \qquad
    Q_K^{-1}e_1
    =\frac1{K+1}\begin{pmatrix}3\\-\sqrt{K-2}\end{pmatrix}.
\end{equation}
Indeed, $\Tr Q_K=K+1$, $\det Q_K=2(K+1)^2/(K+7)$, and
\[
 \det\left(Q_K-\frac32\Id\right)
 =\frac{2K^2-23K+29}{4(K+7)}>0
 \qquad(K\geq16),
\]
while the first diagonal entry of $Q_K-3\Id/2$ is positive.  The upper bound follows in the
same way from the principal minors of $K\Id-Q_K$.
Choose $m_0$ with $\norm{m_0}=e^2r_0$ and angle $\theta_0$ satisfying
\begin{equation}
\label{eq:valley-phase-init}
    \theta_0+c\log\frac{\norm{m_0}}{r_0}\equiv0\pmod\pi,
\end{equation}
and initialize the precision by
\begin{equation}
\label{eq:spiral-P0}
    P_0:=C_0^{-1}:=hR_{\theta_0}Q_KR_{\theta_0}^\top.
\end{equation}
Then
\begin{equation}
\label{eq:spiral-C0-band}
    \frac1\beta\Id\preceq C_0\preceq\frac1\alpha\Id.
\end{equation}

The point-system reference orbit is
\begin{equation}
\label{eq:reference-orbit}
    r_t^{\mathrm{ref}}
    :=e^2r_0\exp\left(-\frac{3t}{K+1}\right),
    \qquad
    \theta_t^{\mathrm{ref}}
    :=\theta_0+\frac{\sqrt{K-2}}{K+1}t,
\end{equation}
with
\[
    P_t^{\mathrm{ref}}
    :=hR_{\theta_t^{\mathrm{ref}}}Q_KR_{\theta_t^{\mathrm{ref}}}^\top.
\]
Indeed, on the valley $\nabla V_K(m)=hm$, and \eqref{eq:QK-band} gives
\[
    \dot r=-\frac3{K+1}r,
    \qquad
    r\dot\theta=\frac{\sqrt{K-2}}{K+1}r.
\]
A direct substitution also verifies
$\dot P^{\mathrm{ref}}=\nabla^2V_K(m^{\mathrm{ref}})-P^{\mathrm{ref}}$.

\begin{lemma}[Gaussian shadowing certificate]
\label{lem:spiral-shadowing}
Assume
\begin{equation}
\label{eq:r0-shadowing}
    r_0\sqrt\alpha
    \geq
    10^{25}K^{3/2}\sqrt{\log(eK)}.
\end{equation}
Let $(m_t,C_t)$ be the exact Gaussian Fisher--Rao flow for $V_K$ initialized by
\eqref{eq:valley-phase-init}--\eqref{eq:spiral-P0}.  Then, for
$0\leq t\leq K/100$,
\begin{equation}
\label{eq:radial-shadowing}
    -\frac4K
    \leq
    \frac{\dd}{\dd t}\log\norm{m_t}
    \leq
    -\frac2K.
\end{equation}
In particular,
\begin{equation}
\label{eq:mean-lower-shadowing}
    \norm{m_t}\geq e^{-0.04}\norm{m_0}
    \qquad(0\leq t\leq K/100).
\end{equation}
\end{lemma}

\begin{proof}
Write $P_t=C_t^{-1}$ and choose the continuous lift of the polar angle starting at $\theta_0$.
Set
\begin{equation}
\label{eq:spiral-transverse-coordinates}
 z_t:=\theta_t+c\log\frac{r_t}{r_0},
 \qquad
 S_t:=\frac1hR_{\theta_t}^{\top}P_tR_{\theta_t},
 \qquad
 E_t:=S_t-Q_K,
\end{equation}
where $z_t$ is represented near $0$ modulo $\pi$.  Put
\[
 w:=\begin{pmatrix}c\\1\end{pmatrix},
 \qquad
 J:=\begin{pmatrix}0&-1\\1&0\end{pmatrix},
 \qquad
 d_K:=\frac{D_K}{h}=\frac{18(K-1)}{K+7}.
\]
On the uncut flat phase core, $f_K(z)=z^2/2$, and the normalized point gradient and Hessian are
\begin{align}
\label{eq:spiral-core-point-fields}
 b_K(z)
 &:=\frac1{hr}R_\theta^\top\nabla V_K(m)
 =\begin{pmatrix}
 1+\frac{d_K}{2}z^2+\frac{d_Kc}{2}z\\[1mm]
 \frac{d_K}{2}z
 \end{pmatrix},\\
 \mathsf H_K(z)
 &:=\frac1hR_\theta^\top\nabla^2V_K(m)R_\theta\notag\\
 &=\Id+\frac{d_K}{2}ww^\top
 +d_K\begin{pmatrix}
 \frac12z^2+\frac32cz&\frac12z\\[1mm]
 \frac12z&\frac12z^2+\frac12cz
 \end{pmatrix}.
\end{align}
For the exact Gaussian flow define
\[
 \bar b_t:=\frac1{hr_t}R_{\theta_t}^\top\E[\nabla V_K(X_t)],
 \qquad
 \bar H_t:=\frac1hR_{\theta_t}^\top\E[\nabla^2V_K(X_t)]R_{\theta_t},
 \qquad
 q_t:=S_t^{-1}\bar b_t.
\]
The mean and precision equations then become
\begin{equation}
\label{eq:spiral-rotated-system}
 \frac{\dot r_t}{r_t}=-(q_t)_1,
 \qquad
 \dot\theta_t=-(q_t)_2,
 \qquad
 \dot z_t=-w^\top q_t,
 \qquad
 \dot S_t=\bar H_t-S_t-(q_t)_2(S_tJ-JS_t).
\end{equation}
In particular, if the Gaussian averages in \eqref{eq:spiral-rotated-system} are replaced by the
point fields \eqref{eq:spiral-core-point-fields}, then $(z,S)=(0,Q_K)$ is an equilibrium.  Indeed,
with $\omega_K=s_K/(K+1)$,
\[
 Q_K^{-1}e_1=\begin{pmatrix}3/(K+1)\\-\omega_K\end{pmatrix},
 \qquad
 w^\top Q_K^{-1}e_1=0,
 \qquad
 \mathsf H_K(0)-Q_K+\omega_K(Q_KJ-JQ_K)=0.
\]

We next give the linear and nonlinear certificates explicitly.  Let
\begin{equation}
\label{eq:spiral-y-scaling}
 \mathsf T_K:=\diag(K,\sqrt K,1,\sqrt K),
 \qquad
 y:=\mathsf T_K(z,E_{11},E_{12},E_{22})^\top.
\end{equation}
For $S=Q_K+E$, let $\mathcal F_K(y)$ be the scaled point vector field obtained from
\eqref{eq:spiral-rotated-system} with $\bar b=b_K(z)$ and $\bar H=\mathsf H_K(z)$, and set
\[
 A_K:=D\mathcal F_K(0),
 \qquad
 \mathcal N_K(y):=\mathcal F_K(y)-A_Ky.
\]
If $s=s_K=\sqrt{K-2}$, direct differentiation gives
\begin{equation}
\label{eq:spiral-linear-matrix}
 A_K=\mathsf T_K\mathsf B_K\mathsf T_K^{-1},
\end{equation}
where, in the order $(z,E_{11},E_{12},E_{22})$,
{\scriptsize
\[
\mathsf B_K=
\begin{pmatrix}
-\frac{3(K-1)}{2(K+1)}&0&\frac{K+7}{2(K+1)^2}&-\frac{s(K+7)}{6(K+1)^2}\\[1mm]
\frac{3s(K-1)}{K+7}&-\frac{K^2+20K-17}{(K+1)(K+7)}&\frac{12s}{K+7}&-\frac{2(K-2)}{K+7}\\[1mm]
\frac{3(K-1)(K+1)}{2(K+7)}&\frac{2(K-11)s}{(K+1)(K+7)}&-\frac{7K^2-10K+19}{2(K+1)(K+7)}&\frac{s(K^2-2K+9)}{2(K+1)(K+7)}\\[1mm]
\frac{9s(K-1)}{K+7}&\frac{12(K-2)}{(K+1)(K+7)}&-\frac{12s}{K+7}&\frac{K-11}{K+7}
\end{pmatrix}.
\]
}
Its characteristic polynomial has the factorization
\begin{equation}
\label{eq:spiral-characteristic}
 \det(\lambda\Id-A_K)
 =\frac{\lambda+1}{2(K+1)^3}\,p_K(\lambda+1),
\end{equation}
where
\[
 p_K(\mu)
 =2(K+1)^3\mu^3
 +2(K-5)(K+1)^2\mu^2
 +4(K-2)(K+1)^2\mu
 +(K-2)(K^2-10K-23).
\]
For $K\geq16$, all coefficients are positive and the cubic Routh--Hurwitz determinant is
\[
 6(K-2)(K-1)^2(K+1)^3>0.
\]
Thus every eigenvalue of $A_K$ has real part at most $-1$.  The entries of the scaled matrix obey
\[
 \bigl(|(A_K)_{ij}|\bigr)
 \preceq
 \begin{pmatrix}
 3/2&0&2/3&1/4\\
 3&3/2&12&2\\
 3/2&1/8&7/2&1/2\\
 9&3/4&12&1
 \end{pmatrix},
\]
so $\norm{A_K}_{\op}<21$.  In a complex Schur decomposition $A_K=U(D+N)U^*$,
$N$ is strictly upper triangular, $\norm{N}_{\op}\leq42$, and
$\norm{e^{Dt}}_{\op}\leq e^{-t}$.  Iterating Duhamel's formula terminates after three factors of
$N$ and yields
\begin{equation}
\label{eq:spiral-semigroup}
 \norm{e^{A_Kt}}_{\op}
 \leq e^{-t}\sum_{j=0}^3\frac{(42t)^j}{j!}
 <2\cdot10^6e^{-t/2}.
\end{equation}

The nonlinear estimate is also an elementary finite-dimensional calculation.  Inserting the
explicit inverse of the $2\times2$ matrix $S$ into \eqref{eq:spiral-core-point-fields} gives, for
$\norm{y}\leq10^{-12}$ and $K\geq16$,
\begin{equation}
\label{eq:spiral-second-derivative-table}
 \sup\sum_{j,\ell=1}^4
 \left|\partial_{j\ell}^2(\mathcal F_K)_i\right|
 \leq(100,2000,500,2000)_i.
\end{equation}
For completeness, the bounds follow term by term from $d_K\leq18$, $c/\sqrt K\leq1/3$,
$Q_K\succeq3\Id/2$, and
\[
 E=\begin{pmatrix}y_2/\sqrt K&y_3\\y_3&y_4/\sqrt K\end{pmatrix};
\]
the largest denominator is controlled by
$\norm{(Q_K+E)^{-1}}_{\op}\leq1$.  Taylor's theorem and
\eqref{eq:spiral-second-derivative-table} imply
\begin{equation}
\label{eq:spiral-nonlinear}
 \norm{\mathcal N_K(y)}
 \leq\frac12\sqrt{100^2+2000^2+500^2+2000^2}\,\norm{y}^2
 <4000\norm{y}^2.
\end{equation}

It remains to bound the Gaussian-versus-point error at the actual state.  Define the stopping time
\begin{equation}
\label{eq:spiral-stopping-time}
 \tau:=\frac K{100}\wedge
 \inf\left\{t\geq0:\norm{y_t}=10^{-12}
 \ \text{or}\ r_t\notin[e^{3/2}r_0,e^{5/2}r_0]\right\}.
\end{equation}
On $[0,\tau]$, \eqref{eq:QK-band} and \eqref{eq:spiral-y-scaling} give
$S_t\succeq\Id$, hence $C_t\preceq h^{-1}\Id$.  Put $Z_t:=X_t-m_t$ and
\[
 \mathcal G_t:=\left\{\norm{Z_t}\leq\frac{\delta_Kr_t}{8\sqrt K}\right\},
 \qquad
 p_K:=\sup_{t\leq\tau}\Prob(\mathcal G_t^c),
 \qquad
 M_K:=\sup_{t\leq\tau}
 \E\left[\sqrt h\norm{Z_t}\mathbf1_{\mathcal G_t^c}\right].
\]
For every point on the segment $m_t+uZ_t$, $0\leq u\leq1$, the logarithmic radius and angle
obey, on $\mathcal G_t$,
\[
 |\Delta\theta|+c|\Delta\ell|
 \leq2(1+c)\frac{\norm{Z_t}}{r_t}<\frac{\delta_K}{4}.
\]
Since $|z_t|\leq10^{-12}/K<\delta_K/4$ and $r_t\geq e^{3/2}r_0$, the whole segment lies in the
uncut flat phase core.  If $G\sim\N(0,\Id_2)$, the covariance bound implies
$\sqrt h\norm{Z_t}\leq\norm{G}$ under a common Gaussian representation.  The exact radial tail in
dimension two therefore gives
\begin{equation}
\label{eq:spiral-tail-certificate}
 p_K
 \leq\exp\left(-\frac{\delta_K^2hr_0^2}{128K}\right)
 \leq10^{-50}K^{-30},
 \qquad
 M_K\leq\sqrt{2p_K}\leq\sqrt2\,10^{-25}K^{-15}.
\end{equation}

On the flat core, differentiating \eqref{eq:spiral-core-point-fields} in Cartesian coordinates gives
\begin{equation}
\label{eq:spiral-core-Hessian-Lipschitz}
 \operatorname{Lip}_{\op}(\nabla^2V_K)
 \leq\frac{15hK}{r_t}
\end{equation}
along these segments.  One way to see the constant is to note that
$\norm{\partial_z\mathsf H_K}_{\op}\leq12\sqrt K$,
$\norm{\nabla z}\leq2\sqrt K/(5r)$, and rotation of the polar frame contributes at most
$2h(K+0.02)/r$; allowing a factor two for the radius along the segment gives
\eqref{eq:spiral-core-Hessian-Lipschitz}.

Let
\begin{align*}
 \varepsilon_b(t)
 &:=\left\|\frac1{hr_t}R_{\theta_t}^\top
   \bigl(\E[\nabla V_K(X_t)]-\nabla V_K(m_t)\bigr)\right\|,\\
 \varepsilon_H(t)
 &:=\left\|\frac1hR_{\theta_t}^\top
   \bigl(\E[\nabla^2V_K(X_t)]-\nabla^2V_K(m_t)\bigr)
   R_{\theta_t}\right\|_{\op}.
\end{align*}
Put $A_0:=\sqrt\alpha\,r_0$.  Splitting at $\mathcal G_t$, using
\eqref{eq:spiral-core-Hessian-Lipschitz} for the Taylor remainder on $\mathcal G_t$, the global
Hessian bound on its complement, $\E Z_t=0$, and
$\E\norm{Z_t}^2\leq2/h$, gives
\begin{equation}
\label{eq:spiral-average-errors}
 \varepsilon_b(t)\leq8\frac K{A_0^2}+3\frac K{A_0}M_K,
 \qquad
 \varepsilon_H(t)\leq15\frac K{A_0}+3Kp_K.
\end{equation}
The forcing $\xi_K(t)$ in the scaled transverse system is, by definition, the contribution of
these actual-state Gaussian-versus-point errors.  Since
$\norm{S_t^{-1}}_{\op}\leq1$, $\norm{S_t}_{\op}\leq2K$, and
$\norm{w}\leq2\sqrt K/5$, equations \eqref{eq:spiral-rotated-system} and
\eqref{eq:spiral-y-scaling} imply
\begin{align}
\label{eq:smoothing-certificate}
 \norm{\xi_K(t)}
 &\leq\frac25K^{3/2}\varepsilon_b(t)
   +\sqrt{2K}\bigl(\varepsilon_H(t)+4K\varepsilon_b(t)\bigr)\notag\\
 &\leq23\frac{K^{3/2}}{A_0}
   +5K^{3/2}p_K
   +52\frac{K^{5/2}}{A_0^2}
   +20\frac{K^{5/2}}{A_0}M_K.
\end{align}
The radius condition \eqref{eq:r0-shadowing} and \eqref{eq:spiral-tail-certificate} make the last
line smaller than $3\cdot10^{-24}$, and in particular smaller than $10^{-22}$, on $[0,\tau]$.

The exact transverse equation is now
\begin{equation}
\label{eq:transverse-system}
 \dot y_t=A_Ky_t+\mathcal N_K(y_t)+\xi_K(t),
 \qquad y_0=0.
\end{equation}
Variation of constants, \eqref{eq:spiral-semigroup}, \eqref{eq:spiral-nonlinear}, and
\eqref{eq:smoothing-certificate} yield, for $t\leq\tau$,
\[
 \norm{y_t}
 \leq2\cdot10^6\int_0^te^{-(t-s)/2}
 \bigl(4000\cdot10^{-24}+10^{-22}\bigr)\,\dd s
 <2\cdot10^{-14}.
\]
Thus the transverse component cannot cause the exit in \eqref{eq:spiral-stopping-time}.

Finally, define the point radial field
$\mathcal R_K(y):=-e_1^\top S^{-1}b_K(z)$.  Differentiating the same explicit $2\times2$ inverse
gives
\[
 \sup_{\norm{y}\leq10^{-12}}\norm{\nabla\mathcal R_K(y)}
 \leq\frac{20}{K},
 \qquad
 \mathcal R_K(0)=-\frac3{K+1}.
\]
The Gaussian correction is at most $\varepsilon_b(t)$, so on $[0,\tau]$,
\[
 \left|\frac{\dd}{\dd t}\log r_t+\frac3{K+1}\right|
 \leq\frac{20}{K}\norm{y_t}+\varepsilon_b(t)
 <\frac1{4K}.
\]
This already implies $-4/K\leq\dot r_t/r_t\leq-2/K$.  Hence the radius cannot reach the upper
barrier in \eqref{eq:spiral-stopping-time}, and before time $K/100$ it remains above
$e^{2-0.04}r_0>e^{3/2}r_0$.  Therefore $\tau=K/100$, which proves
\eqref{eq:radial-shadowing}; integration gives \eqref{eq:mean-lower-shadowing}.
\end{proof}

\subsection{Optimizer whitening and the lower bound}

\begin{lemma}[Optimizer covariance and intrinsic condition number of the spiral]
\label{lem:spiral-kappa-star}
Let $C_\star$ be the optimizing covariance for $V_K$ and define
\[
    \eta_K:=K\exp\left(-\frac{\alpha r_0^2}{2}\right).
\]
Under \eqref{eq:r0-shadowing}, $\eta_K\leq10^{-3}$ and
\begin{equation}
\label{eq:spiral-Pstar}
    h(1-\eta_K)\Id
    \preceq
    C_\star^{-1}
    \preceq
    h(1+\eta_K)\Id.
\end{equation}
Moreover,
\begin{equation}
\label{eq:spiral-kappa-star}
    K\leq\kappastar(V_K)\leq1.03K,
\end{equation}
and the whitened initial covariance satisfies
\begin{equation}
\label{eq:spiral-initial-product}
    \betastar\lambda_{0,\star}
    \geq\frac{1-\eta_K}{1+\eta_K}.
\end{equation}
Thus the covariance-initialization term in \eqref{eq:cont-star-hitting-simple} is smaller than
$0.003$.
\end{lemma}

\begin{proof}
The potential is exactly quadratic with Hessian $h\Id$ on $\{\norm{x}\leq r_0\}$.  Since
$C_\star\preceq\alpha^{-1}\Id$, write $X_\star=C_\star^{1/2}Z$ with
$Z\sim\N(0,\Id_2)$.  Then
$\norm{X_\star}>r_0$ implies $\norm{Z}>\sqrt\alpha r_0$, and for a standard two-dimensional
Gaussian,
$\Prob(\norm{Z}>s)=e^{-s^2/2}$.  Using
$\norm{\nabla^2V_K-h\Id}_\op\leq hK$ and
$C_\star^{-1}=\E[\nabla^2V_K(X_\star)]$ gives \eqref{eq:spiral-Pstar}.

Combining \eqref{eq:spiral-Hessian-bounds} with \eqref{eq:spiral-Pstar} gives
\[
    \alphastar\geq\frac{0.98}{1+\eta_K},
    \qquad
    \betastar\leq\frac{K+0.02}{1-\eta_K},
\]
and hence $\kappastar\leq1.03K$.  At the origin the Hessian is $h\Id$, whereas at a valley
point it has an eigenvalue $hK$.  Therefore
\[
    \alphastar\leq h\lambda_{\min}(C_\star),
    \qquad
    \betastar\geq hK\lambda_{\min}(C_\star),
\]
which gives $\kappastar\geq K$.

Finally, \eqref{eq:spiral-C0-band} gives
$C_0\succeq(hK)^{-1}\Id$, while
$C_\star^{-1}\succeq h(1-\eta_K)\Id$.  Thus
$\lambda_{0,\star}\geq(1-\eta_K)/K$.  The valley Hessian and
\eqref{eq:spiral-Pstar} also give
$\betastar\geq K/(1+\eta_K)$, proving
\eqref{eq:spiral-initial-product}.
\end{proof}

\begin{theorem}[Continuous-time logarithmic-spiral lower bound]
\label{thm:spiral-sharp}
For every $K\geq16$, the potential $V_K$ of \eqref{eq:spiral-potential}, with $r_0$ satisfying
\eqref{eq:r0-shadowing}, is globally $C^3$ and satisfies
$\alpha\Id\preceq\nabla^2V_K\preceq\beta\Id$ with $\kappa=2K+1$.  The initialization
\eqref{eq:valley-phase-init}--\eqref{eq:spiral-P0} is admissible under the general spectral
condition $(1/\beta)\Id\preceq C_0\preceq(1/\alpha)\Id$.  For any local level satisfying
\begin{equation}
\label{eq:spiral-local-level}
    \delta_{\mathrm{loc}}
    <\frac\alpha2e^{-0.08}\norm{m_0}^2,
\end{equation}
the exact Gaussian Fisher--Rao flow obeys
\begin{equation}
\label{eq:spiral-lower-time}
    \boxed{
    T_{\mathrm{loc}}
    :=\inf\{t\geq0:\Delta(a_t)\leq\delta_{\mathrm{loc}}\}
    \geq\frac K{100}
    \geq\frac{\kappastar(V_K)}{103}.}
\end{equation}
\end{theorem}

\begin{proof}
The potential is even, so the optimizing mean is $m_\star=0$.  Strong convexity in the mean and
global optimality of $(0,C_\star)$ imply, for every $(m,C)$,
\[
    \Delta(m,C)
    \geq\frac\alpha2\norm{m}^2.
\]
By \Cref{lem:spiral-shadowing}, for $0\leq t\leq K/100$,
\[
    \Delta(a_t)
    \geq
    \frac\alpha2e^{-0.08}\norm{m_0}^2
    >\delta_{\mathrm{loc}}.
\]
Thus $T_{\mathrm{loc}}\geq K/100$.  The second inequality in
\eqref{eq:spiral-lower-time} follows from
$\kappastar\leq1.03K$ in \Cref{lem:spiral-kappa-star}.  The admissibility and Hessian claims follow
from \Cref{lem:spiral-Hessian} and \eqref{eq:spiral-C0-band}.
\end{proof}

\begin{corollary}[Sharp polynomial dependence in dimension at least two]
\label{cor:cont-sharp}
For $N_\theta\geq2$, the worst-case polynomial dependence of the continuous-time
global-localization phase on the affine-invariant condition number is
\begin{equation}
\label{eq:cont-sharp-rate}
    \boxed{T_{\mathrm{loc}}=\widetilde\Theta(\kappastar)}
\end{equation}
under the general spectral initialization class, where the tilde suppresses logarithmic dependence
on initialization size and localization accuracy.  For the spiral family,
$\kappastar=\Theta(\kappa)$ and
$T_{\mathrm{loc}}\geq(\kappa-1)/200$.
\end{corollary}

\begin{proof}
The upper bound is \Cref{thm:cont-global}; the lower bound is
\Cref{thm:spiral-sharp}.  For dimensions larger than two, add independent quadratic coordinates,
\[
    V_K^{(d)}(x)
    :=V_K(x_1,x_2)+\frac h2\sum_{j=3}^d x_j^2,
\]
and initialize those coordinates at equilibrium.  The block-diagonal subsystem is dynamically
invariant.  Moreover, Hadamard's determinant inequality shows that a cross-covariance block can
only decrease the Gaussian determinant while leaving the separated potential expectation
unchanged; hence the optimizing covariance is block diagonal.  The two-dimensional slow subsystem
and its condition-number comparison are therefore unchanged.
\end{proof}

\begin{remark}[Hessian-Lipschitz version]
Differentiating \eqref{eq:spiral-Hessian-decomposition} once more and using
$\norm{\chi'''}_\infty\leq52.5$ and
$\norm{g_K'}_\infty\leq4.376/\delta_K$ gives the explicit upper bound
\[
    \operatorname{Lip}_{\op}(\nabla^2V_K)
    \leq
    10^8\frac{hK^{5/2}}{r_0}.
\]
Hence, for every prescribed $L_H>0$, choosing
\[
    r_0\geq
    \max\left\{
        \frac{10^8hK^{5/2}}{L_H},
        \frac{10^{25}K^{3/2}\sqrt{\log(eK)}}{\sqrt\alpha}
    \right\}
\]
preserves the lower bound while ensuring
$\operatorname{Lip}_{\op}(\nabla^2V_K)\leq L_H$.  At the endpoint $L_H=0$ the target is
quadratic, so the linear lower bound is not asserted under a radius budget fixed independently of
$L_H$.
\end{remark}

\begin{remark}[Scope of the sharpness theorem]
The matching construction uses an admissible anisotropic initial covariance.  It does not establish
an $\Omega(\kappastar)$ lower bound under the stricter initialization
$C_0=\beta^{-1}\Id$.  It also allows the initial radius to grow with $K$ and, in the
Hessian-Lipschitz version, with $1/L_H$; the logarithm of that radius is part of the input size
charged by the upper theorem.  These restrictions should be retained whenever the sharpness result
is quoted.
\end{remark}

\section{Riemannian scheme: spectral bounds and two-sided covariance bootstrap}

\begin{lemma}[Riemannian covariance bootstrap]
\label{lem:Riem-cov}
Assume \Cref{assump:logconcave-smooth}, $\lambda_{0,\min}\Id\preceq C_0\preceq\lambda_{0,\max}\Id$,
and
\[
    0<\Delta t\leq\min\left\{1,\frac1{\beta\lambda_{\max}}\right\}.
\]
Then $C_n\preceq\lambda_{\max}\Id$ for all $n\geq0$. Define
\[
    L_0:=\min\left\{\lambda_0,\frac1{2\beta}\right\},
    \qquad
    L_{n+1}:=\frac{e^{\Delta t}L_n}{1+(e-1)\Delta t\,\beta L_n}.
\]
Then $C_n\succeq L_n\Id$ for all $n\geq0$, and with
\[
    y_0:=\frac1{\beta L_0},
    \qquad
    y_\infty:=\frac{(e-1)\Delta t}{e^{\Delta t}-1}\in[1,e-1)\subset[1,2),
\]
one has $y_n:=(\beta L_n)^{-1}=y_\infty+(y_0-y_\infty)e^{-n\Delta t}$. Hence
$C_n\succeq\tfrac1{2\beta}\Id$ for all
\[
    n\geq
    N_{\cov,-}^{\Riem}
    :=
    \left\lceil
        \frac1{\Delta t}
        \logp\frac{y_0-y_\infty}{2-y_\infty}
    \right\rceil
    =
    O\!\left(\frac1{\Delta t}\logp\frac1{\beta\lambda_0}\right).
\]
Moreover, with $z_0:=\tfrac1{\alpha\lambda_{0,\max}}$ and
$z_\infty:=\tfrac{\Delta t}{e^{\Delta t}-1}\in[\tfrac1{e-1},1)$, one has
$z_n:=\lambda_{\min}(C_n^{-1})/\alpha\geq z_\infty+(z_0-z_\infty)e^{-n\Delta t}$, hence
$C_n\preceq\tfrac2\alpha\Id$ for all
\[
    n\geq
    N_{\cov,+}^{\Riem}
    :=
    \begin{cases}
        0, & z_0\geq\tfrac12,\\[4pt]
        \left\lceil
            \dfrac1{\Delta t}
            \log\dfrac{z_\infty-z_0}{z_\infty-1/2}
        \right\rceil, & z_0<\tfrac12,
    \end{cases}
    \qquad
    =
    O\!\left(\frac1{\Delta t}\logp(\alpha\lambda_{0,\max})\right).
\]
The case split is genuine: when $z_0\geq\tfrac12$ the bound holds already at $n=0$, whereas
when $z_0<\tfrac12$ one necessarily has $z_0<\tfrac12<z_\infty$, so the logarithm has a
well-defined positive argument $(z_\infty-z_0)/(z_\infty-\tfrac12)>1$ (the earlier one-line
form with $\logp$ is undefined for $z_0>z_\infty$).
\end{lemma}

\begin{proof}
Write $B_n:=C_n^{1/2}A(a_n)C_n^{1/2}$, so that the covariance update reads
$C_{n+1}=C_n^{1/2}\exp(\Delta t(\Id-B_n))C_n^{1/2}$ and
$C_{n+1}^{-1}=C_n^{-1/2}\exp(\Delta t(B_n-\Id))C_n^{-1/2}$.

\emph{Upper bound.} Using $e^X\succeq\Id+X$ for symmetric $X$,
\[
    C_{n+1}^{-1}
    \succeq
    C_n^{-1/2}\bigl[\Id+\Delta t(B_n-\Id)\bigr]C_n^{-1/2}
    =
    (1-\Delta t)C_n^{-1}+\Delta t\,A(a_n),
\]
using $C_n^{-1/2}B_nC_n^{-1/2}=A(a_n)$. Inductively assume $C_n\preceq\lambda_{\max}\Id$;
then $C_n^{-1}\succeq\lambda_{\max}^{-1}\Id$ and $A(a_n)\succeq\alpha\Id\succeq\lambda_{\max}^{-1}\Id$
(since $\lambda_{\max}\geq1/\alpha$). As $\Delta t\leq1$,
\[
    C_{n+1}^{-1}
    \succeq
    \bigl[(1-\Delta t)\lambda_{\max}^{-1}+\Delta t\,\alpha\bigr]\Id
    \succeq
    \lambda_{\max}^{-1}\Id,
\]
hence $C_{n+1}\preceq\lambda_{\max}\Id$; the base case is the hypothesis on $C_0$.

\emph{Lower envelope.} Since $A(a_n)\preceq\beta\Id$ and $C_n\preceq\lambda_{\max}\Id$,
$0\preceq B_n\preceq\beta\lambda_{\max}\Id$, so $0\preceq\Delta t\,B_n\preceq\Id$ (because
$\Delta t\leq1/(\beta\lambda_{\max})$). The scalar inequality $e^x\leq1+(e-1)x$ on $[0,1]$
gives, by functional calculus, $\exp(\Delta t\,B_n)\preceq\Id+(e-1)\Delta t\,B_n$, hence
\[
    C_{n+1}^{-1}
    =
    e^{-\Delta t}C_n^{-1/2}\exp(\Delta t\,B_n)C_n^{-1/2}
    \preceq
    e^{-\Delta t}\bigl[C_n^{-1}+(e-1)\Delta t\,A(a_n)\bigr].
\]
If $C_n\succeq L_n\Id$ then $C_n^{-1}\preceq L_n^{-1}\Id$ and $A(a_n)\preceq\beta\Id$, so
\[
    C_{n+1}^{-1}
    \preceq
    e^{-\Delta t}\bigl(L_n^{-1}+(e-1)\Delta t\,\beta\bigr)\Id,
\]
i.e.\ $C_{n+1}\succeq L_{n+1}\Id$. With $y_n=(\beta L_n)^{-1}$ the recurrence becomes the
affine recursion $y_{n+1}=e^{-\Delta t}(y_n+(e-1)\Delta t)$, whose solution is
$y_n=y_\infty+(y_0-y_\infty)e^{-n\Delta t}$ with $y_\infty=(e-1)\Delta t/(e^{\Delta t}-1)$.
For $\Delta t\in(0,1]$ the map $\Delta t\mapsto y_\infty$ is continuous with limit $e-1$ as
$\Delta t\to0$ and value $1$ at $\Delta t=1$, and is monotone, so $y_\infty\in[1,e-1)\subset[1,2)$.
The condition $C_n\succeq\tfrac1{2\beta}\Id$ is $y_n\leq2$, i.e.\
$(y_0-y_\infty)e^{-n\Delta t}\leq2-y_\infty$, giving $N_{\cov,-}^{\Riem}$. Since $L_0=\min\{\lambda_0,1/(2\beta)\}$, this is
$O(\Delta t^{-1}\logp(1/(\beta\lambda_0)))$.

\emph{Upper envelope.} Using $e^X\succeq\Id+X$ again,
$C_{n+1}^{-1}\succeq e^{-\Delta t}(C_n^{-1}+\Delta t\,A(a_n))\succeq e^{-\Delta t}(C_n^{-1}+\Delta t\,\alpha\Id)$.
Hence $p_n:=\lambda_{\min}(C_n^{-1})$ obeys $p_{n+1}\geq e^{-\Delta t}(p_n+\Delta t\,\alpha)$, and with
$z_n=p_n/\alpha$ the equality recursion has solution
$z_n=z_\infty+(z_0-z_\infty)e^{-n\Delta t}$ with $z_\infty=\Delta t/(e^{\Delta t}-1)$. For
$\Delta t\in(0,1]$, $z_\infty\in[\tfrac1{e-1},1)$, in particular $z_\infty>1/2$. The
condition $C_n\preceq\tfrac2\alpha\Id$ is $z_n\geq1/2$. If $z_0\geq1/2$ then, since $z_n$ is
a convex combination of $z_0$ and $z_\infty$ and both are $\geq1/2$, it holds for all $n$
(so $N_{\cov,+}^{\Riem}=0$). If $z_0<1/2$ then $z_0<1/2<z_\infty$, the sequence $z_n$
increases to $z_\infty$, and $z_n\geq1/2$ first holds when
$e^{-n\Delta t}\leq(z_\infty-\tfrac12)/(z_\infty-z_0)$, giving the stated case-split
$N_{\cov,+}^{\Riem}$, which is $O(\Delta t^{-1}\logp(\alpha\lambda_{0,\max}))$.
\end{proof}


\subsection{A proved Riemannian one-step descent estimate}
\label{subsec:Riem-descent}

The Riemannian covariance update in \eqref{eq:Riem-det} is an exponential retraction in the
covariance variable.  We prove its descent property directly, rather than importing a
geodesic-smoothness statement.  The proof uses only a Gaussian coupling, smoothness of the
potential, Gaussian integration by parts, and a scalar exponential inequality.

\begin{lemma}[A scalar exponential inequality]
\label{lem:scalar-retraction}
For every $w\in[-1,1]$,
\[
    \left(e^{w/2}-1-\frac w2\right)
    +\frac12\left(e^{w/2}-1\right)^2
    \leq \frac{w^2}{2}.
\]
\end{lemma}

\begin{proof}
Put $v=w/2\in[-1/2,1/2]$.  Since
\[
    (e^v-1)+\frac12(e^v-1)^2=\frac12(e^{2v}-1),
\]
the desired inequality is equivalent, with $x=2v=w\in[-1,1]$, to
\[
    e^x\leq 1+x+x^2.
\]
Indeed,
\[
    e^x-1-x
    =\sum_{k\geq2}\frac{x^k}{k!}
    \leq x^2\sum_{k\geq2}\frac{|x|^{k-2}}{k!}
    \leq (e-2)x^2<x^2,
\]
which proves the claim.
\end{proof}

\begin{lemma}[Retraction smoothness and deterministic Riemannian descent]
\label{lem:Riem-descent}
Assume $0\preceq\nabla^2V\preceq\beta\Id$ and $C\preceq\Lambda\Id$.  Set
$K:=\beta\Lambda$ and suppose $K\geq1$.  For a tangent vector $\zeta=(u,U)$ at
$a=(m,C)$ define
\[
    W:=C^{-1/2}UC^{-1/2},
    \qquad
    \operatorname{Retr}_a(\zeta)
    :=\left(m+u,\;C^{1/2}e^W C^{1/2}\right).
\]
If $\|W\|_{\op}\leq1$, then
\begin{equation}
\label{eq:Riem-retraction-smooth}
    \calE(\operatorname{Retr}_a(\zeta))
    \leq
    \calE(a)+\inner{\grad\calE(a)}{\zeta}_a
    +K\|\zeta\|_a^2.
\end{equation}
Consequently, for the deterministic update \eqref{eq:Riem-det}, if
$0<h\leq1/K$, then
\begin{equation}
\label{eq:Riem-onestep}
    \boxed{
    \calE(a^+)\leq
    \calE(a)-h(1-Kh)\Gradnorm(a)^2 .
    }
\end{equation}
In particular, if $h\leq1/(2K)$, then
\[
    \calE(a^+)\leq\calE(a)-\frac h2\Gradnorm(a)^2.
\]
\end{lemma}

\begin{proof}
Let $Z\sim\N(0,\Id)$ and introduce the coupling
\[
    X:=m+C^{1/2}Z\sim\rho_a,
    \qquad
    Y:=m+u+C^{1/2}e^{W/2}Z
    \sim\rho_{\operatorname{Retr}_a(\zeta)}.
\]
Write $E_W:=e^{W/2}-\Id$ and $M:=C^{1/2}A(a)C^{1/2}$.  The entropy change is exact:
\[
    \calE_1(\operatorname{Retr}_a(\zeta))-\calE_1(a)
    =-\frac12\Tr W.
\]
By $\beta$-smoothness of $V$,
\[
    \calE_2(\operatorname{Retr}_a(\zeta))-\calE_2(a)
    \leq
    \E\!\left[\nabla V(X)^\top(Y-X)\right]
    +\frac\beta2\E\|Y-X\|^2.
\]
Because $Y-X=u+C^{1/2}E_WZ$, $\E[\nabla V(X)]=-g(a)$, and Gaussian integration by
parts gives $\E[Z\nabla V(X)^\top]=C^{1/2}A(a)$, the first-order contribution equals
\[
    -g(a)^\top u+\Tr(E_WM).
\]
The cross term in the quadratic contribution vanishes, and $C\preceq\Lambda\Id$ yields
\[
    \frac\beta2\E\|Y-X\|^2
    \leq
    \frac K2\,u^\top C^{-1}u
    +\frac K2\|E_W\|_{\Frob}^2.
\]
On the other hand, since
$\grad\calE(a)=(-Cg(a),-C R(a)C)$ and
$C^{1/2}R(a)C^{1/2}=\Id-M$,
\[
    \inner{\grad\calE(a)}{\zeta}_a
    =-g(a)^\top u-\frac12\Tr((\Id-M)W).
\]
Subtracting this linear term from the entropy and covariance contributions leaves
\[
    \Theta
    :=\Tr\!\left[\left(e^{W/2}-\Id-\frac W2\right)M\right]
      +\frac K2\|e^{W/2}-\Id\|_{\Frob}^2.
\]
The matrix $N_W:=e^{W/2}-\Id-W/2$ is positive semidefinite, while
$0\preceq M\preceq K\Id$.  Hence
\[
    \Tr(N_WM)\leq K\Tr N_W.
\]
Diagonalizing $W$ and applying \Cref{lem:scalar-retraction} eigenvalue by eigenvalue gives
\[
    \Theta
    \leq K\sum_i\left[
       e^{w_i/2}-1-\frac{w_i}{2}
       +\frac12(e^{w_i/2}-1)^2
    \right]
    \leq \frac K2\|W\|_{\Frob}^2.
\]
Combining the mean and covariance estimates, and using
$\|\zeta\|_a^2=u^\top C^{-1}u+\tfrac12\|W\|_{\Frob}^2$, proves
\eqref{eq:Riem-retraction-smooth}.

For the deterministic step take
$\zeta=-h\grad\calE(a)=(hCg(a),hC R(a)C)$.  Then
\[
    W=hC^{1/2}R(a)C^{1/2}
      =h\bigl(\Id-C^{1/2}A(a)C^{1/2}\bigr).
\]
The eigenvalues of the matrix in parentheses lie in $[1-K,1]$, so
$\|W\|_{\op}\leq hK\leq1$.  Moreover
\[
    \inner{\grad\calE(a)}{\zeta}_a=-h\Gradnorm(a)^2,
    \qquad
    \|\zeta\|_a^2=h^2\Gradnorm(a)^2.
\]
Substitution into \eqref{eq:Riem-retraction-smooth} gives
\eqref{eq:Riem-onestep}.
\end{proof}

\begin{theorem}[Improved global hitting time, Riemannian scheme]
\label{thm:Riem-global}
Under \Cref{assump:logconcave-smooth}, with
\[
    0<\Delta t\leq\frac1{2\beta\lambda_{\max}},
\]
the hitting time
$N_{\Riem}^{\glob}(\delta):=\inf\{n\geq0:\Delta(a_n)\leq\delta\}$ of \eqref{eq:Riem-det}
satisfies, for any $\delta\in(0,\Delta_0)$,
\[
    \boxed{
    N_{\Riem}^{\glob}(\delta)
    =
    O\!\left(
        \frac1{\Delta t}\logp\frac1{\beta\lambda_0}
        +
        \frac\kappa{\Delta t}\logp\frac{\Delta_0}{\delta}
    \right).
    }
\]
\end{theorem}

\begin{proof}
Apply \Cref{lem:Riem-descent} with $K=\beta\lambda_{\max}$.  Since
$\Delta t\leq1/(2K)$,
\[
    \Delta(a_{n+1})
    \leq
    \Delta(a_n)-\frac{\Delta t}{2}\Gradnorm(a_n)^2.
\]
If $C_n\succeq L_n\Id$, \Cref{lem:FR-W2} gives
$\Gradnorm(a_n)^2\geq\alpha L_n\Delta(a_n)$, and therefore
\[
    \Delta(a_{n+1})
    \leq
    \left(1-\frac{\alpha L_n\Delta t}{2}\right)\Delta(a_n).
\]
After $N_{\cov,-}^{\Riem}$ iterations, \Cref{lem:Riem-cov} gives
$L_n\geq1/(2\beta)$, so the per-step factor is at most
$1-\Delta t/(4\kappa)$.  The covariance burn-in costs
$O(\Delta t^{-1}\logp(1/(\beta\lambda_0)))$ iterations, and reducing the energy gap from
$\Delta_0$ to $\delta$ afterward costs
$O((\kappa/\Delta t)\logp(\Delta_0/\delta))$, using
$-\log(1-x)\geq x$.
\end{proof}

\section{KL/Bregman scheme: spectral bounds and exact two-sided covariance bootstrap}

\begin{lemma}[KL covariance bootstrap]
\label{lem:KL-cov}
Assume \Cref{assump:logconcave-smooth} and $\lambda_{0,\min}\Id\preceq C_0\preceq\lambda_{0,\max}\Id$.
Along \eqref{eq:KL-det} (and, pathwise, along the stochastic KL update with $A(a_n)$
replaced by $H_n$), $C_n\preceq\lambda_{\max}\Id$ for all $n\geq0$. Define
\[
    L_0:=\min\left\{\lambda_0,\frac1{2\beta}\right\},
    \qquad
    L_{n+1}=\frac{(1+\Delta t)L_n}{1+\Delta t\,\beta L_n},
\]
so that $C_n\succeq L_n\Id$, with the explicit solution
\[
    L_n
    =
    \frac1{\beta\bigl[1+(\tfrac1{\beta L_0}-1)(1+\Delta t)^{-n}\bigr]} .
\]
Hence $C_n\succeq\tfrac1{2\beta}\Id$ for all
\[
    n\geq
    N_{\cov,-}^{\KL}
    :=
    \left\lceil
        \frac{\logp\bigl(\tfrac1{\beta L_0}-1\bigr)}{\log(1+\Delta t)}
    \right\rceil
    =
    O\!\left(\frac1{\Delta t}\logp\frac1{\beta\lambda_0}\right),
\]
and $C_n\preceq\tfrac2\alpha\Id$ for all
\[
    n\geq
    N_{\cov,+}^{\KL}
    :=
    \left\lceil\frac{\log2}{\log(1+\Delta t)}\right\rceil
    =
    O\!\left(\frac1{\Delta t}\right).
\]
\end{lemma}

\begin{proof}
The update is $C_{n+1}^{-1}=\tfrac1{1+\Delta t}(C_n^{-1}+\Delta t\,A(a_n))$ with
$\alpha\Id\preceq A(a_n)\preceq\beta\Id$ (the same bound holds pathwise for $H_n$).

\emph{Upper bound.} If $C_n\preceq\lambda_{\max}\Id$ then $C_n^{-1}\succeq\lambda_{\max}^{-1}\Id$
and $A(a_n)\succeq\alpha\Id\succeq\lambda_{\max}^{-1}\Id$, so
$C_{n+1}^{-1}\succeq\tfrac1{1+\Delta t}(\lambda_{\max}^{-1}+\Delta t\lambda_{\max}^{-1})\Id=\lambda_{\max}^{-1}\Id$,
i.e.\ $C_{n+1}\preceq\lambda_{\max}\Id$.

\emph{Lower envelope.} If $C_n\succeq L_n\Id$ then $C_n^{-1}\preceq L_n^{-1}\Id$,
$A(a_n)\preceq\beta\Id$, so $C_{n+1}^{-1}\preceq\tfrac1{1+\Delta t}(L_n^{-1}+\Delta t\beta)\Id$,
i.e.\ $C_{n+1}\succeq L_{n+1}\Id$. With $y_n=(\beta L_n)^{-1}$ the recursion is
$y_{n+1}=\tfrac{y_n+\Delta t}{1+\Delta t}=1+\tfrac{y_n-1}{1+\Delta t}$, so
$y_n=1+(y_0-1)(1+\Delta t)^{-n}$, which is the stated closed form. The condition
$C_n\succeq\tfrac1{2\beta}\Id$ is $y_n\leq2$, i.e.\
$(y_0-1)(1+\Delta t)^{-n}\leq1$, giving $N_{\cov,-}^{\KL}$ (here $y_0=1/(\beta L_0)$,
and $\logp(y_0-1)$ covers the case $y_0\le2$).

\emph{Upper envelope.} With $A(a_n)\succeq\alpha\Id$, $p_n:=\lambda_{\min}(C_n^{-1})$ obeys
$p_{n+1}\geq\tfrac{p_n+\Delta t\,\alpha}{1+\Delta t}$, and with $z_n=p_n/\alpha$ the equality
recursion is $z_{n+1}=1+\tfrac{z_n-1}{1+\Delta t}$, so $z_n=1+(z_0-1)(1+\Delta t)^{-n}$. If
$z_0\geq1$ then $z_n\geq1\geq\tfrac12$ for all $n$. If $0\leq z_0<1$ then $z_n\geq\tfrac12$
once $(1-z_0)(1+\Delta t)^{-n}\leq\tfrac12$, for which $(1+\Delta t)^{-n}\leq\tfrac12$ (i.e.\
$n\geq N_{\cov,+}^{\KL}$) suffices since $1-z_0\leq1$. In either case $C_n\preceq\tfrac2\alpha\Id$
for $n\geq N_{\cov,+}^{\KL}$.
\end{proof}


\subsection{A proved KL/Bregman one-step estimate}
\label{subsec:KL-descent}

For the displayed KL update, the useful descent divergence is the \emph{forward} divergence
$\KL(\rho_{a^+}\|\rho_a)$.  We prove the estimate directly.  This orientation is important:
the reverse divergence $\KL(\rho_a\|\rho_{a^+})$ does not satisfy the corresponding descent
inequality for arbitrary mean displacement.

\begin{lemma}[Scalar inequality for the KL covariance update]
\label{lem:scalar-KL}
Let $K\geq1$, $0<h\leq1/K$, $\mu\in[0,K]$, and
\[
    b:=\sqrt{\frac{1+h}{1+h\mu}}.
\]
Then
\begin{equation}
\label{eq:scalar-KL}
    \mu(b-1)-\log b+\frac K2(b-1)^2
    \leq
    -\frac1{2h}\bigl(b^2-1-2\log b\bigr).
\end{equation}
\end{lemma}

\begin{proof}
Because $K\leq1/h$ and $(b-1)^2\geq0$, it suffices to replace $K$ on the left by $1/h$.
The relation between $b$ and $\mu$ is
\[
    \mu=\frac{(1+h)b^{-2}-1}{h}.
\]
After multiplying the resulting inequality by $h$ and moving all terms to the left, define
\[
    G_h(b)
    :=h\mu(b-1)-h\log b+\frac12(b-1)^2
      +\frac12(b^2-1-2\log b).
\]
A direct differentiation gives
\[
    G_h(1)=0,
    \qquad
    G_h'(b)=\frac{(b-1)\bigl(2b^3-(1+h)(b+2)\bigr)}{b^3}.
\]
Since $0\leq h\mu\leq hK\leq1$,
\[
    b\in\left[\sqrt{\frac{1+h}{2}},\sqrt{1+h}\right].
\]
The polynomial $p_h(b):=2b^3-(1+h)(b+2)$ is convex on $(0,\infty)$, and at the two
endpoints of this interval,
\[
    p_h\!\left(\sqrt{\frac{1+h}{2}}\right)=-2(1+h)<0,
    \qquad
    p_h(\sqrt{1+h})=(1+h)(\sqrt{1+h}-2)<0,
\]
where $h\leq1$.  Convexity therefore implies $p_h(b)<0$ throughout the interval.  Hence
$G_h'(b)\geq0$ for $b\leq1$ and $G_h'(b)\leq0$ for $b\geq1$.  Thus $G_h$ attains its maximum
at $b=1$, where it vanishes, proving \eqref{eq:scalar-KL}.
\end{proof}

\begin{proposition}[Forward-KL descent and Wasserstein-gradient control]
\label{prop:KL-onestep}
Assume $0\preceq\nabla^2V\preceq\beta\Id$ and
$L\Id\preceq C\preceq\Lambda\Id$.  Put $K:=\beta\Lambda\geq1$ and let
$0<h\leq1/K$.  For the KL update
\[
    m^+=m+hCg(a),
    \qquad
    C^+=(1+h)(C^{-1}+hA(a))^{-1},
\]
one has
\begin{equation}
\label{eq:KL-forward-descent}
    \boxed{
    \calE(a^+)\leq\calE(a)-\frac1h\KL(\rho_{a^+}\|\rho_a).
    }
\end{equation}
Moreover,
\begin{equation}
\label{eq:KL-forward-gradient-lower}
    \KL(\rho_{a^+}\|\rho_a)
    \geq
    \frac{Lh^2}{4(1+hK)^2}
    \norm{\grad^{W_2}\calE(a)}_{W_2}^2.
\end{equation}
\end{proposition}

\begin{proof}
Set
\[
    M:=C^{1/2}A(a)C^{1/2},
    \qquad
    B:=\sqrt{1+h}\,(\Id+hM)^{-1/2}.
\]
Then $0\preceq M\preceq K\Id$, $B$ is a function of $M$, and
$C^+=C^{1/2}B^2C^{1/2}$.  With $Z\sim\N(0,\Id)$, couple
\[
    X:=m+C^{1/2}Z\sim\rho_a,
    \qquad
    Y:=m+hCg(a)+C^{1/2}BZ\sim\rho_{a^+}.
\]
The entropy change is
\[
    \calE_1(a^+)-\calE_1(a)=-\log\det B.
\]
Using $\beta$-smoothness of $V$, Gaussian integration by parts, and
$C\preceq\Lambda\Id$ exactly as in \Cref{lem:Riem-descent},
\[
\begin{aligned}
    \calE_2(a^+)-\calE_2(a)
    &\leq
    -h\,g(a)^\top Cg(a)
    +\Tr\bigl((B-\Id)M\bigr)\\
    &\qquad
    +\frac{Kh^2}{2}g(a)^\top Cg(a)
    +\frac K2\|B-\Id\|_{\Frob}^2.
\end{aligned}
\]
Because $hK\leq1$, the mean contribution is at most
$-(h/2)g(a)^\top Cg(a)$, which equals minus $h^{-1}$ times the mean part of
$\KL(\rho_{a^+}\|\rho_a)$.

Let $\mu_i$ and $b_i=\sqrt{(1+h)/(1+h\mu_i)}$ be the aligned eigenvalues of $M$ and $B$.
The remaining covariance contribution is the sum of
\[
    \mu_i(b_i-1)-\log b_i+\frac K2(b_i-1)^2.
\]
By \Cref{lem:scalar-KL}, this is bounded by
\[
    -\frac1{2h}\sum_i(b_i^2-1-2\log b_i),
\]
which is exactly minus $h^{-1}$ times the covariance part of
$\KL(\rho_{a^+}\|\rho_a)$.  This proves \eqref{eq:KL-forward-descent}.

For \eqref{eq:KL-forward-gradient-lower}, write
$R:=C^{-1}-A(a)$ and $S:=C^{1/2}RC^{1/2}=\Id-M$.  The mean part of the forward KL is
\[
    \frac12(m^+-m)^\top C^{-1}(m^+-m)
    =\frac{h^2}{2}g(a)^\top Cg(a)
    \geq\frac{Lh^2}{2}\|g(a)\|^2.
\]
For the covariance part, let
\[
    Q:=C^{1/2}(C^+)^{-1}C^{1/2}=\frac{\Id+hM}{1+h}.
\]
If $q>0$, then
\begin{equation}
\label{eq:KL-scalar-lower}
    q^{-1}-1+\log q
    \geq
    \frac{(q-1)^2}{2\max\{q,1\}^2}.
\end{equation}
Indeed, for $q\geq1$ the left side is
$\int_1^q(t-1)t^{-2}\,\dd t\geq(q-1)^2/(2q^2)$, while for $q\leq1$ it is
$\int_q^1(1-t)t^{-2}\,\dd t\geq(1-q)^2/2$.
The eigenvalues of $Q$ satisfy
\[
    q_i-1=-\frac{h}{1+h}s_i,
    \qquad
    \max\{q_i,1\}\leq\frac{1+hK}{1+h}.
\]
Applying \eqref{eq:KL-scalar-lower} eigenvalue by eigenvalue therefore gives
\[
    \KL_C(\rho_{a^+}\|\rho_a)
    \geq
    \frac{h^2}{4(1+hK)^2}\|S\|_{\Frob}^2.
\]
Finally,
\[
    \|S\|_{\Frob}^2=\Tr(CRCR)
    =\Tr\bigl(C[RCR]\bigr)
    \geq L\Tr(RCR),
\]
because $RCR\succeq0$ and $C\succeq L\Id$.  Combining the mean and covariance parts gives
\[
    \KL(\rho_{a^+}\|\rho_a)
    \geq
    \frac{Lh^2}{4(1+hK)^2}
    \left(\|g(a)\|^2+\Tr(RCR)\right),
\]
which is \eqref{eq:KL-forward-gradient-lower}.
\end{proof}

\begin{theorem}[Improved global hitting time, KL scheme]
\label{thm:KL-global}
Under \Cref{assump:logconcave-smooth}, with
$0<\Delta t\leq1/(\beta\lambda_{\max})$, the hitting time
$N_{\KL}^{\glob}(\delta):=\inf\{n\geq0:\Delta(a_n)\leq\delta\}$ of \eqref{eq:KL-det}
satisfies, for any $\delta\in(0,\Delta_0)$,
\[
    \boxed{
    N_{\KL}^{\glob}(\delta)
    =
    O\!\left(
        \frac1{\Delta t}\logp\frac1{\beta\lambda_0}
        +
        \frac\kappa{\Delta t}\logp\frac{\Delta_0}{\delta}
    \right).
    }
\]
\end{theorem}

\begin{proof}
Apply \Cref{prop:KL-onestep} with $L=L_n$, $\Lambda=\lambda_{\max}$, and
$K=\beta\lambda_{\max}$.  Combining
\eqref{eq:KL-forward-descent}, \eqref{eq:KL-forward-gradient-lower}, and
\Cref{lem:W2-PL} gives
\[
    \Delta(a_{n+1})
    \leq
    \left(
        1-
        \frac{\alpha L_n\Delta t}{2(1+\Delta t\beta\lambda_{\max})^2}
    \right)\Delta(a_n).
\]
After $N_{\cov,-}^{\KL}$ iterations, \Cref{lem:KL-cov} gives
$L_n\geq1/(2\beta)$.  Since
$\Delta t\beta\lambda_{\max}\leq1$, one has
$(1+\Delta t\beta\lambda_{\max})^2\leq4$, and therefore the per-step factor is at most
$1-\Delta t/(16\kappa)$.  The claimed hitting-time bound follows by adding the covariance
burn-in and using $-\log(1-x)\geq x$.
\end{proof}

\section{Sharpness of fixed-step discrete global localization}
\label{sec:disc-sharp}

The preceding upper bounds have the global-localization scale
\[
    \frac{\kappa}{\Delta t}
    \log\frac{\Delta_0}{\delta}.
\]
For the globally certified stepsize $\Delta t=\Theta(\kappa^{-1})$, this becomes
$\Theta(\kappa^2)$ iterations for a constant-factor reduction of the objective.
We now show that this polynomial dependence is sharp for both explicit Fisher--Rao
schemes.  The construction is one-dimensional, starts with covariance already on the
curvature scale, and can be chosen with a dimension-free Hessian-Lipschitz constant.
Thus the lower bound concerns the global localization stage itself, rather than the
covariance burn-in.

\subsection{Scalar schemes and scope of the result}

We normalize
\[
    \alpha=1,
    \qquad
    \beta=\kappa,
\]
and write $a=(m,c)\in\R\times(0,\infty)$.  For
$X\sim\N(m,c)$ define
\begin{equation}
\label{eq:disc-sharp-bA}
    b(m,c):=\E[V'(X)],
    \qquad
    A(m,c):=\E[V''(X)].
\end{equation}
The scalar Riemannian and KL/Bregman updates are, respectively,
\begin{equation}
\label{eq:disc-sharp-R-map}
    m^+=m-hc\,b(m,c),
    \qquad
    c^+=c\exp\bigl(h[1-cA(m,c)]\bigr),
\end{equation}
and
\begin{equation}
\label{eq:disc-sharp-KL-map}
    m^+=m-hc\,b(m,c),
    \qquad
    c^+=\frac{(1+h)c}{1+hcA(m,c)}.
\end{equation}
We consider a fixed globally safe scale
\begin{equation}
\label{eq:disc-sharp-step}
    h=\frac{\gamma}{\kappa},
    \qquad
    0<\gamma\leq\frac12,
\end{equation}
where $\gamma$ is independent of $\kappa$.  The restriction $\gamma\leq1/2$ places the
step inside the certified range of both \Cref{thm:Riem-global,thm:KL-global}; the same
construction works, with changed numerical constants, for every fixed $\gamma\in(0,2)$.

The lower bound below proves sharpness of the factor $\kappa/h$ for a constant-factor
energy reduction.  It does not claim that the entire product
$\kappa h^{-1}\log(\Delta_0/\delta)$ is sharp for every target accuracy, and it does not
apply to adaptive, relative-curvature, line-search, or implicit discretizations.

\subsection{A flat-top curvature bump}

Fix an even function $\varphi\in C_c^\infty(\R)$ satisfying
\begin{equation}
\label{eq:flat-bump}
    0\leq\varphi\leq1,
    \qquad
    \varphi(u)=1\quad\text{for }|u|\leq\frac12,
    \qquad
    \varphi(u)=0\quad\text{for }|u|\geq1.
\end{equation}
Set
\[
    I_\varphi:=\int_\R\varphi(u)\,\dd u,
    \qquad
    M_\varphi:=\|\varphi'\|_\infty.
\]
Fix a numerical constant $L_0>0$ and define
\begin{equation}
\label{eq:bump-width-mass}
    w_\kappa:=\frac{(\kappa-1)M_\varphi}{L_0},
    \qquad
    B_\kappa:=w_\kappa I_\varphi,
    \qquad
    H_\kappa:=(\kappa-1)B_\kappa.
\end{equation}
Thus
\begin{equation}
\label{eq:bump-scales}
    w_\kappa=\Theta(\kappa),
    \qquad
    H_\kappa=\Theta(\kappa^2),
\end{equation}
where the implicit constants depend only on $\varphi$ and $L_0$.  The reason for the
choice $w_\kappa=\Theta(\kappa)$ is that it makes the Hessian-Lipschitz constant uniform:
\begin{equation}
\label{eq:bump-LH}
    \frac{\kappa-1}{w_\kappa}\|\varphi'\|_\infty=L_0.
\end{equation}

\subsection{The bump train}

Choose numerical constants $T>0$ sufficiently small and $Y>0$ sufficiently large,
depending only on $\varphi$ and $L_0$, and set
\begin{equation}
\label{eq:bump-count-x0}
    N_\kappa:=\left\lfloor\frac{T\kappa^2}{\gamma}\right\rfloor,
    \qquad
    x_0:=\frac{Y\kappa^4}{\gamma}.
\end{equation}
Define the positive bump centers recursively by
\begin{equation}
\label{eq:center-recursion}
    x_{j+1}
    :=x_j-
    \frac{\gamma}{\kappa^2}
    \left[
        x_j+H_\kappa\left(N_\kappa-j+\frac12\right)
    \right],
    \qquad
    0\leq j<N_\kappa.
\end{equation}

\begin{lemma}[Geometry of the bump centers]
\label{lem:bump-centers}
For suitable universal choices of $T$ and $Y$, and all sufficiently large $\kappa$,
\begin{equation}
\label{eq:center-lower}
    x_{N_\kappa}\geq\frac34x_0,
\end{equation}
and
\begin{equation}
\label{eq:center-spacing}
    x_j-x_{j+1}\geq\frac{3Y}{4}\kappa^2,
    \qquad
    0\leq j<N_\kappa.
\end{equation}
In particular, the intervals $[x_j-w_\kappa,x_j+w_\kappa]$ are pairwise disjoint, and
they are disjoint from their reflections about the origin.
\end{lemma}

\begin{proof}
Summing \eqref{eq:center-recursion} and using $x_j\leq x_0$ gives
\begin{align*}
    x_0-x_{N_\kappa}
    &\leq
    \frac{\gamma N_\kappa}{\kappa^2}x_0
    +\frac{\gamma H_\kappa}{\kappa^2}
      \sum_{j=0}^{N_\kappa-1}
      \left(N_\kappa-j+\frac12\right)\\
    &\leq
    T x_0
    +C_\varphi\frac{T^2\kappa^4}{\gamma},
\end{align*}
where $C_\varphi$ depends only on $\varphi$ and $L_0$, because
$H_\kappa\leq C_\varphi\kappa^2$ and
$N_\kappa\leq T\kappa^2/\gamma$.  Since
$x_0=Y\kappa^4/\gamma$, first choose $T$ small and then $Y$ large so that the right-hand
side is at most $x_0/4$.  This proves \eqref{eq:center-lower}.

By \eqref{eq:center-recursion} and \eqref{eq:center-lower},
\[
    x_j-x_{j+1}
    \geq
    \frac{\gamma}{\kappa^2}x_{N_\kappa}
    \geq
    \frac{3\gamma}{4\kappa^2}\frac{Y\kappa^4}{\gamma}
    =\frac{3Y}{4}\kappa^2.
\]
Since $w_\kappa=O(\kappa)$, the positive supports are disjoint for all large $\kappa$.
Moreover $x_{N_\kappa}\asymp\kappa^4/\gamma\gg w_\kappa$, so the positive supports are
also disjoint from the reflected negative supports.
\end{proof}

Define the potential by
\begin{equation}
\label{eq:bump-potential}
    V_\kappa''(x)
    :=1+(\kappa-1)
    \sum_{j=0}^{N_\kappa}
    \left[
        \varphi\left(\frac{x-x_j}{w_\kappa}\right)
        +
        \varphi\left(\frac{x+x_j}{w_\kappa}\right)
    \right],
\end{equation}
with $V_\kappa'(0)=0$ and $V_\kappa(0)=0$.

\begin{lemma}[Regularity and exact condition number]
\label{lem:bump-potential-regularity}
For all sufficiently large $\kappa$, the function $V_\kappa$ is even and belongs to
$C^\infty(\R)$.  Moreover,
\begin{equation}
\label{eq:bump-curvature-bounds}
    1\leq V_\kappa''(x)\leq\kappa
    \qquad\text{for every }x\in\R,
\end{equation}
with both bounds attained, and
\begin{equation}
\label{eq:bump-Hess-Lip}
    \|V_\kappa'''\|_\infty\leq L_0.
\end{equation}
Consequently $e^{-V_\kappa}$ is integrable and the target has exact condition number
$\kappa$.  Its optimizer-whitened condition number also equals $\kappa$.
\end{lemma}

\begin{proof}
The reflected construction makes $V_\kappa''$ even; the conditions at the origin then make
$V_\kappa'$ odd and $V_\kappa$ even.  By \Cref{lem:bump-centers}, at most one summand in
\eqref{eq:bump-potential} is nonzero at any point.  Hence
$1\leq V_\kappa''\leq1+(\kappa-1)=\kappa$.  The value $1$ is attained outside the bump
supports and the value $\kappa$ is attained on every flat top.  Similarly,
\[
    |V_\kappa'''(x)|
    \leq
    \frac{\kappa-1}{w_\kappa}\|\varphi'\|_\infty
    =L_0.
\]
Since $V_\kappa''\geq1$ and $V_\kappa(0)=V_\kappa'(0)=0$,
$V_\kappa(x)\geq x^2/2$, which proves integrability.

In one dimension, optimizer whitening multiplies every Hessian value by the same positive
scalar $c_\star$.  It therefore leaves the ratio
$\sup V_\kappa''/\inf V_\kappa''$ unchanged, so $\kappa_\star=\kappa$.
\end{proof}

\subsection{Gaussian averages along the train}

Set
\begin{equation}
\label{eq:matched-covariance}
    c_\kappa:=\frac1\kappa,
\end{equation}
and define the ideal state $a_j^\circ:=(x_j,c_\kappa)$.  We first compute the averaged
curvature and score at these states.

\begin{lemma}[Averaged curvature and score at the centers]
\label{lem:bump-averages}
There exist constants $C,c>0$ and an integer $p\geq1$, independent of $\kappa$, such that
for every $0\leq j\leq N_\kappa$,
\begin{align}
\label{eq:center-A-estimate}
    A(x_j,c_\kappa)
    &=\kappa+r^A_{j,\kappa},
    &|r^A_{j,\kappa}|
    &\leq C\kappa^p e^{-c\kappa^3},\\
\label{eq:center-b-estimate}
    b(x_j,c_\kappa)
    &=x_j+H_\kappa\left(N_\kappa-j+\frac12\right)
      +r^b_{j,\kappa},
    &|r^b_{j,\kappa}|
    &\leq C\kappa^p e^{-c\kappa^3}.
\end{align}
The same estimates hold uniformly on the tubes
\begin{equation}
\label{eq:bump-tube}
    \mathcal T_{j,\kappa}
    :=\left\{
        (m,c):
        |m-x_j|\leq\frac{w_\kappa}{8},
        \quad
        |\kappa c-1|\leq\frac18
    \right\},
\end{equation}
in the following form: for $(m,c)\in\mathcal T_{j,\kappa}$,
\begin{align}
\label{eq:tube-A}
    |A(m,c)-\kappa|
    &\leq \varepsilon_\kappa,\\
\label{eq:tube-b}
    \left|
        b(m,c)
        -x_j-H_\kappa\left(N_\kappa-j+\frac12\right)
        -\kappa(m-x_j)
    \right|
    &\leq \varepsilon_\kappa,
\end{align}
where
\begin{equation}
\label{eq:epsilon-kappa}
    \varepsilon_\kappa:=C\kappa^p e^{-c\kappa^3}.
\end{equation}
\end{lemma}

\begin{proof}
Let $X=x_j+Z/\sqrt\kappa$, $Z\sim\N(0,1)$.  The current bump equals one whenever
$|X-x_j|\leq w_\kappa/2$.  Since $w_\kappa\asymp\kappa$,
\[
    \Prob\left(|X-x_j|>\frac{w_\kappa}{2}\right)
    \leq2\exp(-c\kappa^3).
\]
Every other bump is at distance at least $c\kappa^2$ from $x_j$, which is at least
$c\kappa^{5/2}$ Gaussian standard deviations.  Summing over
$N_\kappa+1=O(\kappa^2)$ bumps changes only the polynomial prefactor.  This proves
\eqref{eq:center-A-estimate}.

To compute the score, put
\[
    \Phi(s):=\int_{-\infty}^{s}\varphi(u)\,\dd u.
\]
For $x>0$, away from an exponentially negligible event on which the Gaussian sample crosses
the origin, integration of \eqref{eq:bump-potential} gives
\[
    V_\kappa'(x)
    =x+(\kappa-1)w_\kappa
      \sum_{k=0}^{N_\kappa}
      \Phi\left(\frac{x-x_k}{w_\kappa}\right).
\]
The bumps with $k>j$ lie strictly between the origin and $x_j$ and contribute their full
mass $I_\varphi$; those with $k<j$ contribute zero, up to the same Gaussian-tail error.
For the current bump, evenness of $\varphi$ gives
$\Phi(s)+\Phi(-s)=I_\varphi$.  Since $(X-x_j)/w_\kappa$ is symmetric,
\[
    \E\Phi\left(\frac{X-x_j}{w_\kappa}\right)=\frac{I_\varphi}{2}.
\]
Multiplying by $(\kappa-1)w_\kappa$ proves \eqref{eq:center-b-estimate}.

Now take $(m,c)\in\mathcal T_{j,\kappa}$.  The mean remains at least $3w_\kappa/8$
from the edge of the current flat top, and the standard deviation is at most
$\sqrt{9/(8\kappa)}$.  Therefore the probability of reaching a transition layer is still
bounded by $C e^{-c\kappa^3}$, uniformly on the tube.  This proves \eqref{eq:tube-A}.
Finally, $\partial_m b(m,c)=A(m,c)$, so integration in the mean variable gives the term
$\kappa(m-x_j)$ with an error bounded by $\varepsilon_\kappa$.  The covariance derivative
is
\[
    \partial_c b(m,c)=\frac12\E[V_\kappa'''(X)],
\]
and $V_\kappa'''$ is supported only on the transition layers, which are reached with
probability $O(e^{-c\kappa^3})$ throughout the tube.  Integrating in $c$ proves
\eqref{eq:tube-b} after increasing the polynomial prefactor in \eqref{eq:epsilon-kappa}.
\end{proof}

\subsection{Shadowing for both covariance maps}

Let $(m_j,c_j)$ denote the iterates of either \eqref{eq:disc-sharp-R-map} or
\eqref{eq:disc-sharp-KL-map}, initialized at
\begin{equation}
\label{eq:disc-sharp-init}
    m_0=x_0,
    \qquad
    c_0=c_\kappa=\frac1\kappa.
\end{equation}
Define the mean and relative-covariance errors
\begin{equation}
\label{eq:shadow-errors}
    e_j:=m_j-x_j,
    \qquad
    q_j:=\kappa c_j-1.
\end{equation}

\begin{lemma}[Uniform discrete shadowing]
\label{lem:discrete-shadowing}
For either covariance map, there exist constants $C,c>0$ and an integer $p\geq1$ such that,
for all sufficiently large $\kappa$ and every $0\leq j\leq N_\kappa$,
\begin{equation}
\label{eq:shadowing-estimate}
    |e_j|+|q_j|
    \leq C\kappa^p e^{-c\kappa^3}.
\end{equation}
In particular, every iterate belongs to $\mathcal T_{j,\kappa}$ and
\begin{equation}
\label{eq:shadowing-conclusion}
    m_j=x_j+o(1),
    \qquad
    c_j=\frac1\kappa(1+o(1)),
\end{equation}
uniformly over the whole time interval $0\leq j\leq N_\kappa$.
\end{lemma}

\begin{proof}
Write
\[
    B_j:=x_j+H_\kappa\left(N_\kappa-j+\frac12\right).
\]
By \eqref{eq:center-recursion}, $x_{j+1}=x_j-\gamma B_j/\kappa^2$.  On the tube,
\Cref{lem:bump-averages} gives
\[
    b(m_j,c_j)=B_j+\kappa e_j+r_j,
    \qquad
    |r_j|\leq\varepsilon_\kappa.
\]
Since $c_j=(1+q_j)/\kappa$, the mean recursion is exactly
\begin{equation}
\label{eq:mean-shadow-recursion}
    e_{j+1}
    =\left(1-\frac{\gamma(1+q_j)}{\kappa}\right)e_j
      -\frac{\gamma q_j}{\kappa^2}B_j
      -\frac{\gamma(1+q_j)}{\kappa^2}r_j.
\end{equation}
The center geometry gives $|B_j|\leq C\kappa^4/\gamma$.  Hence, as long as
$|q_j|\leq1/8$,
\begin{equation}
\label{eq:mean-shadow-bound}
    |e_{j+1}|
    \leq
    \left(1-\frac{7\gamma}{8\kappa}\right)|e_j|
    +C\kappa^2|q_j|
    +C\frac{\varepsilon_\kappa}{\kappa^2}.
\end{equation}

For the Riemannian covariance map, write
$A(m_j,c_j)=\kappa+\delta_j$ with $|\delta_j|\leq\varepsilon_\kappa$.  Then
\begin{equation}
\label{eq:q-R-exact}
    1+q_{j+1}
    =(1+q_j)
      \exp\left[
        \frac{\gamma}{\kappa}
        \left(
            1-(1+q_j)\left(1+\frac{\delta_j}{\kappa}\right)
        \right)
      \right].
\end{equation}
For $|q_j|\leq1/8$ and large $\kappa$, the derivative of the right-hand side with
respect to $q_j$ is between zero and $1-\gamma/(2\kappa)$.  Its variation with respect to
$\delta_j$ is at most $C\gamma|\delta_j|/\kappa^2$.  Since the map fixes $q=0$ when
$\delta=0$, the mean-value theorem yields
\begin{equation}
\label{eq:q-R-bound}
    |q_{j+1}|
    \leq
    \left(1-\frac{\gamma}{2\kappa}\right)|q_j|
    +C\frac{\varepsilon_\kappa}{\kappa^2}.
\end{equation}

For the KL/Bregman map,
\begin{equation}
\label{eq:q-KL-exact}
    1+q_{j+1}
    =
    \frac{(1+\gamma/\kappa)(1+q_j)}
         {1+(\gamma/\kappa)(1+q_j)(1+\delta_j/\kappa)}.
\end{equation}
The same mean-value argument gives, on $|q_j|\leq1/8$,
\begin{equation}
\label{eq:q-KL-bound}
    |q_{j+1}|
    \leq
    \left(1-\frac{\gamma}{2\kappa}\right)|q_j|
    +C\frac{\varepsilon_\kappa}{\kappa^2}.
\end{equation}

Starting from $q_0=0$, either \eqref{eq:q-R-bound} or \eqref{eq:q-KL-bound} implies
\begin{equation}
\label{eq:q-summed}
    |q_j|
    \leq
    C\frac{\varepsilon_\kappa}{\gamma\kappa}
    \qquad
    (0\leq j\leq N_\kappa).
\end{equation}
Substituting this into \eqref{eq:mean-shadow-bound}, using $e_0=0$, and summing the
geometric recursion gives
\begin{equation}
\label{eq:e-summed}
    |e_j|
    \leq
    C_\gamma\kappa^p e^{-c\kappa^3}
\end{equation}
for a possibly larger integer $p$.  Because exponential decay dominates every polynomial,
\eqref{eq:q-summed}--\eqref{eq:e-summed} are eventually much smaller than $1/8$ and
$w_\kappa/8$, respectively.  The tube assumptions used in the argument therefore close by
induction, proving \eqref{eq:shadowing-estimate}.
\end{proof}

\subsection{The objective gap remains macroscopic}

For $c>0$ define the Gaussian-smoothed potential
\begin{equation}
\label{eq:Fc}
    F_c(m):=\E_{X\sim\N(m,c)}[V_\kappa(X)].
\end{equation}
Up to a target-dependent additive constant,
\begin{equation}
\label{eq:scalar-objective}
    \calE_\kappa(m,c)=F_c(m)-\frac12\log c.
\end{equation}

\begin{lemma}[Constant-factor lower bound for the objective]
\label{lem:disc-objective-gap}
There exists a numerical constant $q\in(0,1)$, independent of $\kappa$, such that the
iterates of either scheme satisfy
\begin{equation}
\label{eq:disc-gap-lower-all}
    \Delta(m_j,c_j)\geq q\,\Delta(m_0,c_0),
    \qquad
    0\leq j\leq N_\kappa,
\end{equation}
for all sufficiently large $\kappa$.
\end{lemma}

\begin{proof}
Since $V_\kappa$ is even, $F_c$ is even and $F_c'(0)=0$.  Moreover,
\[
    F_c''(m)=A(m,c)\geq1.
\]
Hence
\begin{equation}
\label{eq:Fc-lower}
    F_c(m)-F_c(0)\geq\frac12m^2.
\end{equation}
At the matched covariance $c_\kappa=1/\kappa$, \Cref{lem:bump-centers} gives
\begin{equation}
\label{eq:final-mean-lower}
    F_{c_\kappa}(x_j)-F_{c_\kappa}(0)
    \geq\frac12x_j^2
    \geq\frac12x_{N_\kappa}^2
    \geq\frac{9}{32}x_0^2.
\end{equation}

For the converse bound at initialization, integration of
\eqref{eq:bump-potential} gives the global estimate
\begin{equation}
\label{eq:Vprime-global}
    |V_\kappa'(x)|
    \leq |x|+H_\kappa(N_\kappa+1).
\end{equation}
Thus, for $m\geq0$,
\[
    F_{c_\kappa}'(m)
    \leq
    m+\sqrt{\frac{2}{\pi\kappa}}+H_\kappa(N_\kappa+1).
\]
Integrating from zero to $x_0$ yields
\begin{align}
\label{eq:initial-mean-upper}
    F_{c_\kappa}(x_0)-F_{c_\kappa}(0)
    \leq
    \frac12x_0^2
    +\left[
        H_\kappa(N_\kappa+1)
        +\sqrt{\frac{2}{\pi\kappa}}
    \right]x_0
    \leq C_0x_0^2,
\end{align}
where $C_0$ is independent of $\kappa$; the last inequality follows from
$H_\kappa(N_\kappa+1)=O(\kappa^4/\gamma)$ and
$x_0=Y\kappa^4/\gamma$.
Combining \eqref{eq:final-mean-lower} and \eqref{eq:initial-mean-upper} gives
\begin{equation}
\label{eq:mean-ratio}
    F_{c_\kappa}(x_j)-F_{c_\kappa}(0)
    \geq q_0
    \bigl[F_{c_\kappa}(x_0)-F_{c_\kappa}(0)\bigr],
    \qquad
    q_0:=\frac{9}{32C_0}>0.
\end{equation}

For each fixed covariance, $m\mapsto\calE_\kappa(m,c)$ is strictly convex and even, so the
unique optimizing mean is $m_\star=0$.  Let $c_\star$ denote the optimizing covariance and set
\[
    D_\kappa
    :=\calE_\kappa(0,c_\kappa)-\calE_\kappa(0,c_\star)
    \geq0.
\]
At covariance $c_\kappa$,
\[
    \Delta(x_j,c_\kappa)
    =F_{c_\kappa}(x_j)-F_{c_\kappa}(0)+D_\kappa.
\]
After replacing $q_0$ by $\min\{q_0,1\}$, \eqref{eq:mean-ratio} and $D_\kappa\geq0$ imply
\begin{equation}
\label{eq:ideal-gap-ratio}
    \Delta(x_j,c_\kappa)\geq q_0\Delta(x_0,c_\kappa).
\end{equation}

It remains to replace the ideal states by the actual iterates.  On the tube,
$|\partial_m\calE_\kappa|=|b(m,c)|\leq C\kappa^4/\gamma$, while
\[
    \left|\partial_c\calE_\kappa(m,c)\right|
    =\frac12\left|A(m,c)-\frac1c\right|
    \leq C\bigl(\varepsilon_\kappa+\kappa|q|\bigr).
\]
Together with \Cref{lem:discrete-shadowing}, the mean-value theorem therefore gives
\begin{equation}
\label{eq:objective-shadow}
    \sup_{0\leq j\leq N_\kappa}
    \left|
        \calE_\kappa(m_j,c_j)-\calE_\kappa(x_j,c_\kappa)
    \right|
    =o(x_0^2).
\end{equation}
Since \eqref{eq:Fc-lower} gives
$\Delta(x_0,c_\kappa)\geq x_0^2/2$, the $o(x_0^2)$ error is
$o(\Delta(x_0,c_\kappa))$ and can be absorbed into
\eqref{eq:ideal-gap-ratio}, after decreasing $q_0$ by a fixed factor.  This proves
\eqref{eq:disc-gap-lower-all}.
\end{proof}

\begin{theorem}[Sharp fixed-step global localization for both discretizations]
\label{thm:disc-global-sharp}
Fix $\gamma\in(0,1/2]$ and $L_0>0$.  There exist constants
$T,q>0$ and $\kappa_0\geq1$, depending only on $\gamma$, $L_0$, and the fixed bump
$\varphi$, such that for every $\kappa\geq\kappa_0$ there is an even potential
$V_\kappa\in C^\infty(\R)$ satisfying
\begin{equation}
\label{eq:disc-sharp-final-assumptions}
    1\leq V_\kappa''\leq\kappa,
    \qquad
    \|V_\kappa'''\|_\infty\leq L_0,
\end{equation}
and an initialization
\begin{equation}
\label{eq:disc-sharp-final-init}
    m_0=\frac{Y\kappa^4}{\gamma},
    \qquad
    c_0=\frac1\kappa,
\end{equation}
such that, when either \eqref{eq:disc-sharp-R-map} or
\eqref{eq:disc-sharp-KL-map} is run with exact expectations and stepsize
$h=\gamma/\kappa$, one has
\begin{equation}
\label{eq:disc-sharp-final-gap}
    \Delta_n\geq q\Delta_0
    \qquad
    \text{for every }
    0\leq n\leq
    \left\lfloor\frac{T\kappa^2}{\gamma}\right\rfloor.
\end{equation}
Consequently, for constant-factor objective reduction,
\begin{equation}
\label{eq:disc-sharp-final-complexity}
    N_{\Riem}=\Omega\left(\frac{\kappa^2}{\gamma}\right)
    =\Omega\left(\frac{\kappa}{h}\right),
    \qquad
    N_{\KL}=\Omega\left(\frac{\kappa^2}{\gamma}\right)
    =\Omega\left(\frac{\kappa}{h}\right).
\end{equation}
The lower bound holds with $\kappa_\star=\kappa$, with the covariance initially in the
curvature-scale band, and with a dimension-free Hessian-Lipschitz constant.
\end{theorem}

\begin{proof}
The regularity and curvature claims follow from
\Cref{lem:bump-potential-regularity}; the trajectory statement is the combination of
\Cref{lem:discrete-shadowing,lem:disc-objective-gap}.  The equality
$\kappa_\star=\kappa$ is part of \Cref{lem:bump-potential-regularity}, and
$c_0=1/\kappa$ lies exactly at the lower edge of the curvature-scale band
$[1/\kappa,1]$.  Finally, $h=\gamma/\kappa$ converts
$\kappa^2/\gamma$ into $\kappa/h$.
\end{proof}

\begin{corollary}[Sharpness of the global-localization prefactor]
\label{cor:disc-global-sharp}
For the globally fixed-step Riemannian and KL/Bregman schemes, the factor
$\kappa/\Delta t$ in \Cref{thm:Riem-global,thm:KL-global} is sharp, up to constants
independent of $\kappa$ for each fixed $\gamma=\kappa\Delta t$, for a constant-factor
reduction of the energy.  In particular, with
$\Delta t=\Theta(\kappa^{-1})$, both methods have worst-case global-localization complexity
\[
    \Theta(\kappa^2)
\]
in the condition-number exponent.  This remains true under the additional assumption
$\|V'''\|_\infty=O(1)$.
\end{corollary}

\begin{remark}[Why the construction is slow]
Along the shadowed trajectory,
\[
    c_nA(m_n,c_n)=1+o(1),
\]
so both covariance maps are almost stationary.  At the same time $c_n\simeq1/\kappa$ and
$b(m_n,c_n)/m_n=\Theta(1)$.  Hence
\[
    h\frac{c_nb(m_n,c_n)}{m_n}
    =\Theta\left(\frac{\gamma}{\kappa^2}\right),
\]
and a constant-factor motion of the mean requires $\Theta(\kappa^2/\gamma)$ steps.  The
lower bound is therefore the product of the globally safe fixed-step scale
$h=\Theta(\kappa^{-1})$ and a Fisher--Rao global mean scale of order $\kappa^{-1}$.
\end{remark}

\begin{remark}[Limitations of the lower bound]
The theorem proves sharpness of the polynomial condition-number dependence for explicit methods
committed to one globally fixed step.  It does not rule out an $O(\kappa)$ global-localization
result for an adaptive relative-curvature step, Armijo or backtracking line search, or an implicit
or Rosenbrock-stabilized mean update.  Indeed, the construction has
$c_nA(m_n,c_n)\simeq1$, so a method whose stability test is based on the current relative
curvature could accept a step of order one and avoid this particular slowdown.
\end{remark}


\section{Local convergence: the Gaussian core and a Gaussian-sharp region}
\label{sec:local}

This section proves the local convergence used in the three-stage complexity statements.
Its organizing principle is a \emph{Gaussian-core decomposition}. The natural gradient
field of the standard Gaussian target is nonlinear in the covariance but can be dissipated
exactly; treating it as a Taylor remainder---as a full-field expansion around $a_\star$
must---discards a large, genuinely dissipative contribution. We therefore write
\[
    F=F_{\mathrm G}+H,
    \qquad
    F_{\mathrm G}(m,C):=(-Cm,\;C-C^2),
\]
retain $F_{\mathrm G}$ exactly, and Taylor-bound only the non-Gaussian perturbation $H$.
Because $H\equiv0$ whenever the target is Gaussian, the certified energy-sublevel region is
sharp for Gaussian targets: it attains the exact maximal threshold in every dimension
(\Cref{cor:gauss-sharp,prop:sharp}), with the exact rate $2\lambda_{\min}(C)$. For general
strongly log-concave targets the region degrades continuously through a non-Gaussian
modulus $K_{\mathrm{ng}}(\delta)$ that vanishes as the target approaches Gaussianity. The
switching region is an energy sublevel set, forward invariant under energy descent, and
entry is certified online by a computable Bures--Wasserstein residual test
(\Cref{subsec:entry}).

\subsection{Whitened coordinates and the linearized gap}
\label{subsec:whitened}

Recall the optimizer-whitened variables and curvature constants from
\eqref{eq:optimizer-whitening}--\eqref{eq:star-curvature}.  By
\Cref{lem:optimizer-whitening,lem:kappa-star-kappa}, the optimizer is $(0,\Id)$,
$\E_{\N(0,\Id)}[\nabla\widehat V(Z)]=0$,
$\E_{\N(0,\Id)}[\nabla^2\widehat V(Z)]=\Id$, and
$1\leq\kappastar\leq\kappa^2$.  The KL divergence and the dynamics are invariant under
$\theta\mapsto C_\star^{-1/2}(\theta-m_\star)$, so \emph{throughout this section we work in
whitened coordinates}: $a=(m,C)$ is the whitened parameter,
$a_\star=(0,\Id)$, $g(a)=-\E_{\N(m,C)}[\nabla\widehat V]$,
$A(a)=\E_{\N(m,C)}[\nabla^2\widehat V]$, $g(a_\star)=0$, $A(a_\star)=\Id$. Statements
transfer back through the fixed equilibrium norms
\begin{equation}
\label{eq:local-norms}
    \norm{a-a_\star}_\star^2:=\norm{\widehat m}^2+\tfrac12\norm{\widehat C-\Id}_{\Frob}^2,
    \qquad
    \norm{a-a_\star}_{\star,\Delta t}^2:=\norm{\widehat m}^2+\tfrac{1+\Delta t}2\norm{\widehat C-\Id}_{\Frob}^2,
\end{equation}
of which $\norm{\cdot}_\star=\norm{\cdot}_{a_\star}$ is the Fisher--Rao tangent norm
\eqref{eq:tangent-norm} frozen at $a_\star$. We write $\xi:=a-a_\star$, use $D(a):=\norm{\xi}_\star^2$,
and denote the field $F(a):=(Cg(a),\,C-CA(a)C)$, so \eqref{eq:NG-flow} reads $\dot a=F(a)$,
$F(a_\star)=0$. Set
\begin{equation}
\label{eq:Gamma}
    \Gamma:=\min\Bigl\{\beta_\star-1,\,N_\theta\bigl(3+\tfrac4{\sqrt\pi}(1+\log\kappa_\star)\bigr)\Bigr\},
    \qquad
    \gund:=\frac1{4+\Gamma} .
\end{equation}

\begin{definition}[Linearized score operators]
\label{def:linear-score-ops}
Let $Z\sim\N(0,\Id)$ and write $B(Z):=\nabla^2\widehat V(Z)$, so that
$\E B(Z)=\Id$.  For $u\in\R^{N_\theta}$ and $X\in\Sym(N_\theta)$ define
\[
\begin{aligned}
    \mathcal T[X]_k
    &:=\E\!\left[\Tr(XB(Z))Z_k\right],\\
    \mathcal T^*[u]_{ij}
    &:=\E\!\left[(u^\top Z)B_{ij}(Z)\right],\\
    \mathcal H[X]_{ij}
    &:=\E\!\left[B_{ij}(Z)\bigl(Z^\top XZ-\Tr X\bigr)\right].
\end{aligned}
\]
These definitions require only the bounded Hessian assumed above.  Gaussian score
identities give
\[
    u^\top\mathcal T[X]=\Tr(\mathcal T^*[u]X),
    \qquad
    \Tr(\mathcal H[X]Y)=\Tr(X\mathcal H[Y]).
\]
\end{definition}

\begin{lemma}[Diagonal-mode bounds]
\label{lem:linear-diagonal-bounds}
Fix an orthogonal basis in which $X=\diag(\lambda)$, and define
\[
    \widetilde T_{ki}:=\E[B_{ii}(Z)Z_k],
    \qquad
    G_{ij}:=\E[B_{ii}(Z)Z_j^2],
    \qquad
    D:=G-\mathbf 1\mathbf 1^\top .
\]
The matrix $D$ is symmetric as the restriction of the self-adjoint operator $\mathcal H$ to
the diagonal subspace.  Moreover,
\begin{equation}
\label{eq:T-Schur}
    \widetilde T^\top\widetilde T\preceq\Id+D,
\end{equation}
and
\begin{equation}
\label{eq:D-Gamma}
    \lambda_{\max}(D)\leq
    \Gamma
    :=\min\left\{
        \beta_\star-1,
        N_\theta\left(3+\frac4{\sqrt\pi}(1+\log\kappa_\star)\right)
    \right\}.
\end{equation}
\end{lemma}

\begin{proof}
For $\eta(Z):=u+\diag(\lambda)Z$, Gaussian integration by parts gives the exact identity
\begin{equation}
\label{eq:eta-B-identity}
    \E[\eta(Z)^\top B(Z)\eta(Z)]
    =\|u\|^2+2u^\top\widetilde T\lambda+\|\lambda\|^2+\lambda^\top D\lambda.
\end{equation}
Indeed, the constant term is $u^\top\E[B]u=\|u\|^2$, the cross term is
$2u^\top\widetilde T\lambda$, and two Gaussian integrations by parts give, componentwise,
\[
    \E[B_{ij}(Z)Z_iZ_j]=\delta_{ij}+G_{ij}-1.
\]
Summing this identity against $\lambda_i\lambda_j$ produces the last two terms in
\eqref{eq:eta-B-identity}.
Because $B(Z)\succeq0$, the block quadratic form on the right is nonnegative for all
$(u,\lambda)$.  Taking its Schur complement with respect to the identity mean block proves
\eqref{eq:T-Schur}.

We next bound $D$.  Taking $u=0$ in \eqref{eq:eta-B-identity} and using
$B(Z)\preceq\beta_\star\Id$ gives
\[
    \|\lambda\|^2+\lambda^\top D\lambda
    \leq\beta_\star\|\lambda\|^2,
\]
so $\lambda_{\max}(D)\leq\beta_\star-1$.

For the dimension-dependent bound, first note that
\begin{equation}
\label{eq:G-tail-bound}
    \alpha_\star\leq G_{ij}
    \leq 4+\frac4{\sqrt\pi}(1+\log\kappa_\star).
\end{equation}
The lower bound follows from $B_{ii}\geq\alpha_\star$.  For the upper bound, if
$\kappa_\star\leq e^{1/2}$ it is immediate from
$B_{ii}\leq\beta_\star\leq\kappa_\star$.  Otherwise set
$t_0:=\sqrt{2\log\kappa_\star}>1$.  Since $\E B_{ii}=1$,
\[
\begin{aligned}
    G_{ij}
    &\leq
    t_0^2\E B_{ii}
    +\beta_\star\E[Z_j^2\mathbf 1_{\{|Z_j|\geq t_0\}}]\\
    &\leq
    2\log\kappa_\star
    +2\kappa_\star\int_{t_0}^{\infty}t^2\frac{e^{-t^2/2}}{\sqrt{2\pi}}\,\dd t.
\end{aligned}
\]
Integration by parts and Mills' ratio yield
\[
    \int_{t_0}^{\infty}t^2\frac{e^{-t^2/2}}{\sqrt{2\pi}}\,\dd t
    \leq2t_0\frac{e^{-t_0^2/2}}{\sqrt{2\pi}},
\]
so the tail contribution is at most
$4\sqrt{\log\kappa_\star/\pi}$.  This implies \eqref{eq:G-tail-bound}.
Consequently
\[
    |D_{ij}|=|G_{ij}-1|
    \leq3+\frac4{\sqrt\pi}(1+\log\kappa_\star).
\]
Therefore, for every $\lambda$,
\[
    \lambda^\top D\lambda
    \leq
    \left(3+\frac4{\sqrt\pi}(1+\log\kappa_\star)\right)\|\lambda\|_1^2
    \leq
    N_\theta\left(3+\frac4{\sqrt\pi}(1+\log\kappa_\star)\right)\|\lambda\|^2.
\]
Combining the two bounds proves \eqref{eq:D-Gamma}.
\end{proof}

\begin{proposition}[Linearized gap]
\label{prop:linearized}
Let $J_\star:=DF(a_\star)$ and $\Lstar:=-J_\star$.  Then
\begin{equation}
\label{eq:Lstar-formula}
    \Lstar(u,X)
    =\left(
        u+\frac12\mathcal T[X],
        X+\mathcal T^*[u]+\frac12\mathcal H[X]
    \right).
\end{equation}
The operator $\Lstar$ is self-adjoint in $\inner{\cdot}{\cdot}_\star$, and if
\[
    \gstar:=\lambda_{\min}(\Lstar),
    \qquad
    \Lambda_\star:=\lambda_{\max}(\Lstar),
\]
then
\begin{equation}
\label{eq:linear-gap-bounds}
    \boxed{
    \frac1{4+\Gamma}\leq\gstar\leq\Lambda_\star\leq\frac{4+\Gamma}{2}.
    }
\end{equation}
For a Gaussian target, $\mathcal T=0$, $\mathcal H=0$, and
$\gstar=\Lambda_\star=1$.
\end{proposition}

\begin{proof}
Bonnet--Price differentiation at $a_\star=(0,\Id)$ gives
\[
\begin{aligned}
    \E_{\N(u,\Id+X)}[\nabla\widehat V]
    &=u+\frac12\mathcal T[X]+O(\|(u,X)\|_\star^2),\\
    \E_{\N(u,\Id+X)}[\nabla^2\widehat V]
    &=\Id+\mathcal T^*[u]+\frac12\mathcal H[X]+O(\|(u,X)\|_\star^2).
\end{aligned}
\]
Substituting these expansions into
$F(m,C)=(-C\E[\nabla\widehat V],\,C-C\E[\nabla^2\widehat V]C)$ proves
\eqref{eq:Lstar-formula}.  The adjoint identities in
\Cref{def:linear-score-ops} show that $\Lstar$ is self-adjoint.

Its quadratic form is
\begin{equation}
\label{eq:Q-linear}
\begin{aligned}
    Q(u,X)
    &:=\inner{\Lstar(u,X)}{(u,X)}_\star\\
    &=\|u\|^2+u^\top\mathcal T[X]
      +\frac12\|X\|_{\Frob}^2
      +\frac14\Tr(X\mathcal H[X]).
\end{aligned}
\end{equation}
Diagonalize $X$ by an orthogonal change of variables.  In the corresponding rotated
coordinates, \eqref{eq:Q-linear} becomes
\[
    Q(u,\diag\lambda)
    =\|u\|^2+u^\top\widetilde T\lambda
      +\frac12\|\lambda\|^2+\frac14\lambda^\top D\lambda.
\]
Let $\nu:=1/(4+\Gamma)$.  Completing the square in $u$, it is enough to verify
\[
    \frac{1-\nu}{2}\Id+\frac14D
      -\frac1{4(1-\nu)}\widetilde T^\top\widetilde T\succeq0.
\]
Using \eqref{eq:T-Schur} and the fact that the coefficient of $D$ after substitution is
negative, \eqref{eq:D-Gamma} reduces this to
\[
    \frac{1-\nu}{2}-\frac{1+\nu\Gamma}{4(1-\nu)}\geq0,
\]
or equivalently
\[
    1-\nu(4+\Gamma)+2\nu^2\geq0,
\]
which holds for $\nu=1/(4+\Gamma)$.  Thus
$Q(u,X)\geq\nu\|(u,X)\|_\star^2$.

For the upper bound, \Cref{lem:linear-diagonal-bounds} gives
\[
\begin{aligned}
    Q(u,\diag\lambda)
    &\leq
    \|u\|^2+\sqrt{1+\Gamma}\,\|u\|\,\|\lambda\|
      +\left(\frac12+\frac\Gamma4\right)\|\lambda\|^2\\
    &\leq
    \left(1+\frac\Gamma4+\frac14\sqrt{\Gamma^2+8\Gamma+8}\right)
    \left(\|u\|^2+\frac12\|\lambda\|^2\right)\\
    &\leq
    \frac{4+\Gamma}{2}\left(\|u\|^2+\frac12\|\lambda\|^2\right).
\end{aligned}
\]
The spectral theorem now proves \eqref{eq:linear-gap-bounds}.
\end{proof}

\begin{proposition}[Linearized contractions of the deterministic maps]
\label{prop:linearized-maps}
Let $T_{\Riem,h}$ and $T_{\KL,h}$ denote the deterministic maps
\eqref{eq:Riem-det} and \eqref{eq:KL-det} in whitened coordinates.  If
$0<h\leq2/(4+\Gamma)$, then
\begin{align}
\label{eq:Riem-linear-contraction}
    \|DT_{\Riem,h}(a_\star)\xi\|_\star
    &\leq
    \left(1-\frac h{4+\Gamma}\right)\|\xi\|_\star,\\
\label{eq:KL-linear-contraction}
    \|DT_{\KL,h}(a_\star)\xi\|_{\star,h}
    &\leq
    \left(1-\frac h{4+2h+\Gamma}\right)\|\xi\|_{\star,h}.
\end{align}
\end{proposition}

\begin{proof}
For the Riemannian map,
\[
    DT_{\Riem,h}(a_\star)=\Id+hJ_\star=\Id-h\Lstar.
\]
By \Cref{prop:linearized}, the spectrum of $\Lstar$ lies in
$[1/(4+\Gamma),(4+\Gamma)/2]$.  The stepsize condition implies
$h\Lambda_\star\leq1$, so every eigenvalue of $\Id-h\Lstar$ lies in
$[0,1-h/(4+\Gamma)]$, proving \eqref{eq:Riem-linear-contraction}.

For the KL map, direct differentiation of \eqref{eq:KL-det} at $(0,\Id)$ gives
\begin{equation}
\label{eq:KL-linear-map-formula}
    M_{\KL,h}(u,X)
    :=DT_{\KL,h}(a_\star)(u,X)
    =\left(
       (1-h)u-\frac h2\mathcal T[X],
       \frac{X-h\mathcal T^*[u]-\frac h2\mathcal H[X]}{1+h}
    \right).
\end{equation}
This operator is self-adjoint in
\[
    \inner{(u,X)}{(v,Y)}_{\star,h}
    :=u^\top v+\frac{1+h}{2}\Tr(XY),
\]
and satisfies the exact identity
\begin{equation}
\label{eq:KL-linear-Q-identity}
    \|\xi\|_{\star,h}^2
      -\inner{M_{\KL,h}\xi}{\xi}_{\star,h}
    =hQ(\xi),
\end{equation}
where $Q$ is \eqref{eq:Q-linear}.  We next prove the weighted coercivity explicitly.  Diagonalize $X$ as before and let
$\nu>0$.  The inequality
$Q(\xi)\geq\nu\|\xi\|_{\star,h}^2$ is equivalent to
\[
\begin{aligned}
    &(1-\nu)\|u\|^2+u^\top\widetilde T\lambda
    +\frac{1-\nu(1+h)}2\|\lambda\|^2
    +\frac14\lambda^\top D\lambda\geq0.
\end{aligned}
\]
Completing the square in $u$, it is sufficient that
\[
    \frac{1-\nu(1+h)}2\Id+\frac14D
    -\frac1{4(1-\nu)}\widetilde T^\top\widetilde T\succeq0.
\]
Using \eqref{eq:T-Schur} and then \eqref{eq:D-Gamma}, this follows from
\[
    \frac{1-\nu(1+h)}2-\frac{1+\nu\Gamma}{4(1-\nu)}\geq0,
\]
which is equivalent to
\[
    1-\nu(4+2h+\Gamma)+2(1+h)\nu^2\geq0.
\]
At $\nu=1/(4+2h+\Gamma)$ the left side equals $2(1+h)\nu^2>0$.  Hence
\[
    Q(\xi)\geq\frac1{4+2h+\Gamma}\|\xi\|_{\star,h}^2.
\]  Hence \eqref{eq:KL-linear-Q-identity} bounds the largest eigenvalue of
$M_{\KL,h}$ by $1-h/(4+2h+\Gamma)$.

Finally, \Cref{prop:linearized} also gives
$Q(\xi)\leq(4+\Gamma)\|\xi\|_{\star,h}^2/2$.  Therefore
$h\leq2/(4+\Gamma)$ implies
$\inner{M_{\KL,h}\xi}{\xi}_{\star,h}\geq0$, so the spectrum of the self-adjoint map is
nonnegative.  Its operator norm is thus its largest eigenvalue, proving
\eqref{eq:KL-linear-contraction}.
\end{proof}

For the remainder of the local theory, any certified lower bound
$0<\gamma_0\leq\gstar$ may be used.  The universal choice is
$\gamma_0=\gund=1/(4+\Gamma)$, while the Gaussian target has the exact value
$\gstar=1$.

\subsection{Energy geometry: entropy-enhanced coercivity and bands}

\begin{lemma}[Entropy-enhanced coercivity]
\label{lem:entropy-coerc}
Write $B:=C^{1/2}$ with eigenvalues $s_1,\dots,s_{N_\theta}$, and set
$\phi_{\alpha_\star}(s):=s-1-\log s+\tfrac{\alpha_\star}2(s-1)^2$. Then
\begin{equation}
\label{eq:entropy-coerc}
    \Delta(m,C)\geq\frac{\alpha_\star}2\norm{m}^2+\sum_{i=1}^{N_\theta}\phi_{\alpha_\star}(s_i).
\end{equation}
For a Gaussian target ($\alpha_\star=1$) this is an equality, and
$\phi_1(\sqrt c)=\tfrac12(c-1-\log c)$, so
\begin{equation}
\label{eq:gauss-energy}
    \Delta_{\mathrm G}(m,C)=\tfrac12\norm{m}^2+\sum_{i=1}^{N_\theta}\varphi(c_i),
    \qquad
    \varphi(c):=\tfrac12(c-1-\log c),
\end{equation}
with $c_i$ the eigenvalues of $C$.
\end{lemma}

\begin{proof}
With $X=m+BZ$, $\Delta(m,C)=\E[\widehat V(m+BZ)-\widehat V(Z)]-\log\det B$. By
$\alpha_\star$-strong convexity,
$\widehat V(m+BZ)-\widehat V(Z)\geq\nabla\widehat V(Z)^\top(m+(B-\Id)Z)+\tfrac{\alpha_\star}2\norm{m+(B-\Id)Z}^2$.
Taking expectations: $\E[\nabla\widehat V(Z)]^\top m=0$; by Gaussian integration by parts
$\E[\nabla\widehat V(Z)^\top(B-\Id)Z]=\Tr((B-\Id)\E[\nabla^2\widehat V(Z)])=\Tr(B-\Id)$;
and $\E\norm{m+(B-\Id)Z}^2=\norm{m}^2+\norm{B-\Id}_{\Frob}^2$. Hence
$\Delta\geq\tfrac{\alpha_\star}2\norm{m}^2+\Tr(B-\Id)-\log\det B+\tfrac{\alpha_\star}2\norm{B-\Id}_{\Frob}^2$,
which diagonalizes to \eqref{eq:entropy-coerc}. For a Gaussian target the strong-convexity
step is an equality; and
$\phi_1(s)=s-1-\log s+\tfrac12(s-1)^2=\tfrac12(s^2-1-\log s^2)=\tfrac12(c-1-\log c)$ with
$c=s^2$, giving \eqref{eq:gauss-energy}.
\end{proof}

\begin{corollary}[Covariance band and norm conversion]
\label{cor:band}
$\phi_{\alpha_\star}$ is strictly decreasing on $(0,1)$, increasing on $(1,\infty)$, with
minimum $0$ at $1$. For $\delta>0$ let $0<s_-(\delta)<1<s_+(\delta)$ solve
$\phi_{\alpha_\star}(s)=\delta$, and set $\ell_\delta:=s_-(\delta)^2$,
$U_\delta:=s_+(\delta)^2$. Then $\Delta(a)\leq\delta$ implies
$\ell_\delta\Id\preceq C\preceq U_\delta\Id$. Moreover, with
\[
    \vartheta(\delta):=\sup_{s\in[s_-(\delta),s_+(\delta)],\,s\neq1}\frac{(s^2-1)^2}{2\phi_{\alpha_\star}(s)}
    \leq\frac{(1+s_+(\delta))^2}{\alpha_\star+s_+(\delta)^{-2}},
    \qquad
    \mathfrak b(\delta)^2:=\max\Bigl\{\tfrac2{\alpha_\star},\vartheta(\delta)\Bigr\},
\]
one has $\Delta(a)\leq\delta\Rightarrow\norm{a-a_\star}_\star\leq\mathfrak b(\delta)\sqrt\delta$,
and likewise $\norm{a-a_\star}_{\star,\Delta t}\leq\mathfrak b_{\Delta t}(\delta)\sqrt\delta$
with $\mathfrak b_{\Delta t}(\delta)^2:=\max\{2/\alpha_\star,(1+\Delta t)\vartheta(\delta)\}$.
For a Gaussian target, $\ell_\delta$ and $U_\delta$ are the two roots of $\varphi(c)=\delta$.
\end{corollary}

\begin{proof}
The level set $\{\phi_{\alpha_\star}\leq\delta\}=[s_-,s_+]$, and each $\phi_{\alpha_\star}(s_i)\leq\delta$
by \eqref{eq:entropy-coerc}, so $s_i\in[s_-,s_+]$; squaring gives the band. For the
conversion, set $E_m:=\tfrac{\alpha_\star}2\norm{m}^2$, $E_C:=\sum_i\phi_{\alpha_\star}(s_i)$,
so $E_m+E_C\leq\delta$; then $\norm{m}^2=\tfrac2{\alpha_\star}E_m$ and
$\tfrac12\norm{C-\Id}_{\Frob}^2=\tfrac12\sum_i(s_i^2-1)^2\leq\vartheta(\delta)E_C$, whence
$\norm{a-a_\star}_\star^2\leq\max\{2/\alpha_\star,\vartheta(\delta)\}(E_m+E_C)\leq\mathfrak b(\delta)^2\delta$.
The bound on $\vartheta$ uses $\phi_{\alpha_\star}(s)\geq\tfrac{(s-1)^2}2(s_+^{-2}+\alpha_\star)$
on $[s_-,s_+]$.
\end{proof}

\begin{definition}[Regions, hull, moduli]
\label{def:region}
For $\delta>0$, $\Ureg_\delta:=\{a:\Delta(a)\leq\delta\}$ and
$\Hreg_\delta:=\{a_\star+s(a-a_\star):a\in\Ureg_\delta,s\in[0,1]\}$. Set
$\bar\Lambda_\delta:=\max\{U_\delta,1/\alpha_\star\}$ and the non-Gaussian moduli
\[
    \mathcal B_\delta:=\sup_{a\in\Hreg_\delta}\norm{A(a)-\Id}_{\op},
    \qquad
    \Sigma_\delta:=\sup_{a\in\Hreg_\delta}\Bigl(\E_{\rho_a}\norm{\nabla^2\widehat V(X)-A(a)}_{\Frob}^2\Bigr)^{1/2},
    \qquad
    \bar\Sigma_\delta:=\min\{\Sigma_\delta,\beta_\star-\alpha_\star\}.
\]
All three vanish identically for Gaussian targets, and each is nondecreasing in $\delta$.
\end{definition}

\begin{lemma}[Geometry and invariance]
\label{lem:region-geom}
Let $\delta>0$, $a=(m,C)\in\Ureg_\delta$. Then (i) $\ell_\delta\Id\preceq C\preceq U_\delta\Id$
and $\norm{a-a_\star}_\star\leq\mathfrak b(\delta)\sqrt\delta$; (ii) every $a'\in\Hreg_\delta$
also satisfies $\ell_\delta\Id\preceq C'\preceq U_\delta\Id$ and
$\norm{a'-a_\star}_\star\leq\mathfrak b(\delta)\sqrt\delta$; (iii) $\Ureg_\delta$ is compact;
(iv) $\Ureg_\delta$ is forward invariant for the flow, and for a scheme step with
$\calE(a_{n+1})\leq\calE(a_n)$.
\end{lemma}

\begin{proof}
(i) is \Cref{cor:band}. (ii): $C'=(1-s)\Id+sC$ has eigenvalues in $[\ell_\delta,U_\delta]$
(convex combinations of $1$ and $\lambda_i(C)\in[\ell_\delta,U_\delta]$, and
$\ell_\delta\leq1\leq U_\delta$), and $\norm{a'-a_\star}_\star=s\norm{\xi}_\star\leq\mathfrak b(\delta)\sqrt\delta$.
(iii): closed (continuity of $\Delta$) and bounded (part (i) and
$\norm{m}\leq\sqrt{2\delta/\alpha_\star}$). (iv): $\tfrac{\dd}{\dd t}\Delta=-\Gradnorm^2\leq0$
(\Cref{lem:metric-dissipation}); for a scheme step, $\calE(a_{n+1})\leq\calE(a_n)$ gives
$\Delta(a_{n+1})\leq\delta$.
\end{proof}

\subsection{The exact Gaussian core}

\begin{lemma}[Gaussian-core dissipation]
\label{lem:gauss-core}
Let $F_{\mathrm G}(m,C):=(-Cm,\,C-C^2)$, the natural gradient field of the standard
Gaussian target $\widehat V_{\mathrm G}(z)=\tfrac12\norm{z}^2$. Then:
\begin{enumerate}
\item $DF_{\mathrm G}(a_\star)=-\Id$ (the identity on $\R^{N_\theta}\times\Sym(N_\theta)$).
\item For every $a=(m,C)$, with $\xi=a-a_\star$,
\[
    \inner{\xi}{F_{\mathrm G}(a)}_\star
    =-m^\top Cm-\tfrac12\Tr\bigl(C(C-\Id)^2\bigr)
    =-\tfrac12\,q_{\mathrm G}(a)\,D(a),
\]
where
$q_{\mathrm G}(a):=\dfrac{2m^\top Cm+\Tr(C(C-\Id)^2)}{\norm{m}^2+\tfrac12\norm{C-\Id}_{\Frob}^2}$.
\item $q_{\mathrm G}(a)\geq2\lambda_{\min}(C)$, with equality when $m=0$ and $C-\Id$ is
supported on a minimal eigendirection of $C$. Consequently, on $\Ureg_\delta$,
$\inner{\xi}{F_{\mathrm G}(a)}_\star\leq-\ell_\delta D(a)$.
\end{enumerate}
\end{lemma}

\begin{proof}
(1) For $\zeta=(u,U)$ at $a_\star=(0,\Id)$: $D(-Cm)[\zeta]=-(Um+Cu)|_{a_\star}=-u$ and
$D(C-C^2)[\zeta]=U-(UC+CU)|_{a_\star}=-U$, so $DF_{\mathrm G}(a_\star)[\zeta]=-\zeta$.
(2) Write $X=C-\Id$. Then $C-C^2=-C(C-\Id)=-(X+X^2)$ (since $C(C-\Id)=(\Id+X)X=X+X^2$),
so with the inner product $\inner{(m,X)}{(m',X')}_\star=m^\top m'+\tfrac12\Tr(XX')$,
\[
    \inner{\xi}{F_{\mathrm G}(a)}_\star
    =-m^\top Cm+\tfrac12\Tr\bigl(X(-(X+X^2))\bigr)
    =-m^\top Cm-\tfrac12\Tr(X^2(\Id+X))
    =-m^\top Cm-\tfrac12\Tr(C X^2),
\]
using $\Id+X=C$. This equals $-\tfrac12 q_{\mathrm G}(a)D(a)$ by definition of $q_{\mathrm G}$.
(3) $2m^\top Cm\geq2\lambda_{\min}(C)\norm{m}^2$ and $\Tr(CX^2)\geq\lambda_{\min}(C)\norm{X}_{\Frob}^2$,
so the numerator of $q_{\mathrm G}$ is $\geq2\lambda_{\min}(C)(\norm{m}^2+\tfrac12\norm{X}_{\Frob}^2)=2\lambda_{\min}(C)D(a)$;
equality holds when $m=0$ and $X$ lies in the minimal eigenspace. On $\Ureg_\delta$,
$\lambda_{\min}(C)\geq\ell_\delta$ (\Cref{cor:band}), giving the last claim.
\end{proof}

\subsection{The non-Gaussian perturbation}

Write $H:=F-F_{\mathrm G}$. Explicitly, since $g_{\mathrm G}(a)=-m$ and $A_{\mathrm G}(a)=\Id$,
\begin{equation}
\label{eq:H-def}
    H(a)=\bigl(C(g(a)+m),\;-C(A(a)-\Id)C\bigr),
    \qquad H(a_\star)=0 .
\end{equation}
Both components are built from the genuinely non-Gaussian quantities $g(a)+m=-\E_{\rho_a}[\nabla\widehat V(X)-X]$
and $A(a)-\Id=\E_{\rho_a}[\nabla^2\widehat V(X)-\Id]$, each of which vanishes identically when
$\widehat V=\widehat V_{\mathrm G}$. The second-order Gaussian score calculus needed to bound
the second differential of $H$ is recorded next.

\begin{lemma}[Score identities and the sharp second-score constant]
\label{lem:second-score}
Fix $a=(m,C)$, $\zeta=(u,U)$, $v:=C^{-1/2}u$, $S:=C^{-1/2}UC^{-1/2}$,
$t:=\norm{\zeta}_a=(\norm{v}^2+\tfrac12\norm{S}_{\Frob}^2)^{1/2}$, and let $X=m+C^{1/2}Z$.
With $q_\zeta:=v^\top Z+\tfrac12(Z^\top SZ-\Tr S)$ and
$\chi_\zeta:=q_\zeta^2+\bigl(\tfrac12\Tr S^2-\norm{v}^2-2v^\top SZ-Z^\top S^2Z\bigr)$, for any
$\phi$ of polynomial growth and $\Phi(a):=\E_{\rho_a}[\phi(X)]$,
\[
    D\Phi(a)[\zeta]=\E[\phi\,q_\zeta],
    \qquad
    D^2\Phi(a)[\zeta,\zeta]=\E[\phi\,\chi_\zeta].
\]
Moreover $q_\zeta$ has Hermite chaoses $1,2$ and $\chi_\zeta$ has chaoses $2,3,4$ only, so
$\E[q_\zeta]=\E[\chi_\zeta]=0$, $\E[Z\chi_\zeta]=0$; and $\E[q_\zeta^2]=t^2$,
$\norm{\chi_\zeta}_{L^2}\leq\sqrt6\,t^2$ (sharp, attained by rank-one covariance directions).
\end{lemma}

\begin{proof}
\emph{Differentiation identities.} With $m_s=m+su$, $C_s=C+sU$ (positive definite for
small $|s|$), differentiating $\log\rho_{a_s}(x)=-\tfrac12\log\det C_s-\tfrac12(x-m_s)^\top C_s^{-1}(x-m_s)+\text{const}$
at fixed $x=m+C^{1/2}Z$ gives $\partial_s\log\rho_{a_s}|_0=q_\zeta$ and
$\partial_s^2\log\rho_{a_s}|_0=\tfrac12\Tr S^2-\norm{v}^2-2v^\top SZ-Z^\top S^2Z=:p_\zeta$
(the entropy term contributes $-\tfrac12\Tr S$ and $+\tfrac12\Tr S^2$, the quadratic term
$v^\top Z+\tfrac12Z^\top SZ$ and $-\norm{v}^2-2v^\top SZ-Z^\top S^2Z$). Since
$\partial_s\rho_{a_s}=\rho_{a_s}\partial_s\log\rho_{a_s}$ and
$\partial_s^2\rho_{a_s}=\rho_{a_s}((\partial_s\log\rho_{a_s})^2+\partial_s^2\log\rho_{a_s})$,
differentiation under the integral (justified by polynomial growth of $\phi$ and locally
bounded polynomial score factors) yields $D\Phi(a)[\zeta]=\E[\phi\,q_\zeta]$ and
$D^2\Phi(a)[\zeta,\zeta]=\E[\phi\,\chi_\zeta]$ with $\chi_\zeta=q_\zeta^2+p_\zeta$.

\emph{Chaos structure and $\E[q_\zeta^2]$.} Write $\ell:=v^\top Z$ (chaos $1$) and
$Q:=Z^\top SZ-\Tr S$ (chaos $2$), so $q_\zeta=\ell+\tfrac12Q$ has chaoses $1,2$ and, by
orthogonality and $\operatorname{Var}(Z^\top SZ)=2\norm{S}_{\Frob}^2$,
$\E[q_\zeta^2]=\norm{v}^2+\tfrac14\cdot2\norm{S}_{\Frob}^2=t^2$. For $\chi_\zeta$, the
product formula for the second chaos gives $Q^2=(Q^2)_0+(Q^2)_2+(Q^2)_4$ with
$(Q^2)_0=2\norm{S}_{\Frob}^2$ and $(Q^2)_2=4(Z^\top S^2Z-\Tr S^2)=:4R_2$; and
$\ell Q=2v^\top SZ+G_3$ with $G_3$ the pure chaos-$3$ part (its chaos-$1$ part is
$\sum_k\E[\ell Q Z_k]Z_k=2v^\top SZ$, from
$\E[Z_i(Z_jZ_l-\delta_{jl})Z_k]=\delta_{ij}\delta_{lk}+\delta_{il}\delta_{jk}$). Writing
$Z^\top S^2Z=R_2+\Tr S^2$ in $p_\zeta$ gives $p_\zeta=-2v^\top SZ-R_2-\tfrac12\Tr S^2-\norm{v}^2$.
Collecting $\chi_\zeta=\ell^2+\ell Q+\tfrac14Q^2+p_\zeta$ by chaos order: the chaos-$0$ part
is $\norm{v}^2+\tfrac12\norm{S}_{\Frob}^2-\tfrac12\Tr S^2-\norm{v}^2=0$ (as
$\norm{S}_{\Frob}^2=\Tr S^2$); the chaos-$1$ part is $2v^\top SZ-2v^\top SZ=0$; and, using
$\tfrac14(Q^2)_2=R_2$, the chaos-$2$ part is $\chi_2=\ell^2-\norm{v}^2$. Thus
$\chi_2=\ell^2-\norm{v}^2$, $\chi_3=G_3=\ell Q-2v^\top SZ$, $\chi_4=\tfrac14(Q^2)_4$, all
mutually orthogonal; in particular $\E[\chi_\zeta]=0$ and $\E[Z\chi_\zeta]=0$.

\emph{Sharp constant.} By orthogonality
$\norm{\chi_\zeta}_{L^2}^2=\E[\chi_2^2]+\E[\chi_3^2]+\E[\chi_4^2]$. With $a:=\norm{v}^2$,
$b:=\norm{S}_{\Frob}^2$, $c:=v^\top S^2v$, $d_4:=\Tr S^4$: $\E[\chi_2^2]=\operatorname{Var}((v^\top Z)^2)=2a^2$;
$\E[\chi_3^2]=\E[(\ell Q)^2]-4\norm{v^\top SZ}_{L^2}^2=2ab+4c$; and
$\E[\chi_4^2]=\tfrac1{16}\norm{(Q^2)_4}_{L^2}^2=\tfrac12b^2+d_4$. Hence
$\norm{\chi_\zeta}_{L^2}^2=2a^2+2ab+4c+\tfrac12b^2+d_4$. Now $c=v^\top S^2v\leq\norm{v}^2\norm{S}_{\op}^2\leq\norm{v}^2\norm{S}_{\Frob}^2=ab$
and $d_4=\sum_i\lambda_i(S)^4\leq(\sum_i\lambda_i(S)^2)^2=b^2$, so
$\norm{\chi_\zeta}_{L^2}^2\leq2a^2+6ab+\tfrac32b^2\leq6(a+\tfrac b2)^2=6t^4$. Equality holds
for $v=0$ (so $a=c=0$) and $S$ rank one (so $d_4=b^2$): then
$\norm{\chi_\zeta}_{L^2}^2=\tfrac12b^2+b^2=\tfrac32b^2=6(b/2)^2=6t^4$.
\end{proof}

\begin{lemma}[Non-Gaussian second-differential modulus]
\label{lem:Kng}
Define
\[
    K_{\mathrm{ng}}(\delta):=\sup_{a\in\Hreg_\delta}\sup_{\zeta\neq0}
    \frac{\norm{D^2H(a)[\zeta,\zeta]}_\star}{\norm{\zeta}_\star^2}.
\]
Then $K_{\mathrm{ng}}(\delta)\leq\bigl(K_{\mathrm{ng},m}(\delta)^2+\tfrac12K_{\mathrm{ng},C}(\delta)^2\bigr)^{1/2}$
with
\[
    K_{\mathrm{ng},m}(\delta):=\frac{U_\delta^{3/2}}{\ell_\delta^2}
    \Bigl(\sqrt3\,\bar\Sigma_\delta+\Sigma_\delta+\sqrt{\Sigma_\delta^2+2\mathcal B_\delta^2}\Bigr),
    \qquad
    K_{\mathrm{ng},C}(\delta):=\frac{U_\delta^2}{\ell_\delta^2}
    \bigl((4\sqrt2+\sqrt6)\Sigma_\delta+4\mathcal B_\delta\bigr).
\]
In particular $K_{\mathrm{ng}}(\delta)=0$ for every Gaussian target, and
$K_{\mathrm{ng}}(\delta)=O(\mathcal B_\delta+\Sigma_\delta)$ as the target approaches
Gaussianity.
\end{lemma}

\begin{proof}
Write $x:=\norm{v}$, $y:=\norm{S}_{\Frob}/\sqrt2$, $t^2=x^2+y^2$,
$\norm{U}_{\op}\leq\sqrt2U_\delta y$, $\norm{U}_{\Frob}\leq\sqrt2U_\delta y$,
$\norm{C}_{\op}\leq U_\delta$. We bound the two blocks of $D^2H$ and convert
$t\leq\ell_\delta^{-1}\norm{\zeta}_\star$.

\emph{Mean block.} $H_m(a)=C(g(a)+m)$, so
$D^2H_m[\zeta,\zeta]=2U\,D(g+m)[\zeta]+C\,D^2g[\zeta,\zeta]$ (as $D^2m=0$). Since
$g(a)+m=-\E[\nabla\widehat V(X)-X]$ and, by the Ornstein--Uhlenbeck inverse of $q_\zeta$
(with $\nabla\psi_1=v+\tfrac12SZ$),
$D(g+m)[\zeta]=-\E[(\nabla^2\widehat V(X)-\Id)C^{1/2}(v+\tfrac12SZ)]$; splitting
$\nabla^2\widehat V-\Id=(\nabla^2\widehat V-A(a))+(A(a)-\Id)$ and using
$\E[\nabla^2\widehat V-A(a)]=0$, $\E[SZ]=0$,
\[
    \norm{D(g+m)[\zeta]}
    \leq\sqrt{U_\delta}\Bigl(\mathcal B_\delta\norm{v}+\tfrac1{\sqrt2}\Sigma_\delta y\Bigr).
\]
For the second derivative, $\chi_\zeta$ has chaoses $\geq2$, so with
$\psi_2=(-\mathcal L)^{-1}\chi_\zeta$ ($\E\norm{\nabla\psi_2}^2\leq\tfrac12\norm{\chi_\zeta}_{L^2}^2\leq3t^4$,
$\E\nabla\psi_2=0$),
$D^2g[\zeta,\zeta]=-\E[(\nabla^2\widehat V(X)-A(a))C^{1/2}\nabla\psi_2]$, hence
$\norm{D^2g[\zeta,\zeta]}\leq\sqrt{3U_\delta}\,\bar\Sigma_\delta t^2$, where
$\bar\Sigma_\delta=\min\{\Sigma_\delta,\beta_\star-\alpha_\star\}$ arises because one may
bound $\E\norm{(\nabla^2\widehat V-A)C^{1/2}\nabla\psi_2}$ either by
$\norm{\nabla^2\widehat V-A}_{\op}\leq\beta_\star-\alpha_\star$ pointwise (an operator-norm
bound, dimension-free) or by the Frobenius dispersion
$(\E\norm{\nabla^2\widehat V-A}_{\Frob}^2)^{1/2}\leq\Sigma_\delta$ via Cauchy--Schwarz,
taking whichever is smaller; no bound of the form $\norm{\nabla^2\widehat V-A}_{\Frob}\leq\beta_\star-\alpha_\star$
is used. Therefore
\[
    \norm{D^2H_m[\zeta,\zeta]}
    \leq U_\delta^{3/2}\bigl(2\sqrt2\mathcal B_\delta xy+2\Sigma_\delta y^2+\sqrt3\bar\Sigma_\delta t^2\bigr),
\]
whose quadratic form in $(x,y)$ has largest eigenvalue
$\sqrt3\bar\Sigma_\delta+\Sigma_\delta+\sqrt{\Sigma_\delta^2+2\mathcal B_\delta^2}$, giving
$K_{\mathrm{ng},m}$ after $t\leq\ell_\delta^{-1}\norm{\zeta}_\star$.

\emph{Covariance block.} $H_C(a)=-C(A(a)-\Id)C=-CWC$ with $W:=A(a)-\Id$; the product rule gives
$D^2(CWC)[\zeta,\zeta]=2U(DW)C+2C(DW)U+2UWU+C(D^2W)C$. With
$\norm{DW[\zeta]}_{\Frob}=\norm{DA[\zeta]}_{\Frob}\leq\Sigma_\delta t$,
$\norm{W}_{\op}\leq\mathcal B_\delta$, and $\norm{D^2W[\zeta,\zeta]}_{\Frob}=\norm{D^2A[\zeta,\zeta]}_{\Frob}\leq\sqrt6\Sigma_\delta t^2$,
\[
    \norm{D^2H_C[\zeta,\zeta]}_{\Frob}
    \leq4\sqrt2U_\delta^2\Sigma_\delta ty+4\mathcal B_\delta U_\delta^2y^2+\sqrt6U_\delta^2\Sigma_\delta t^2
    \leq U_\delta^2\bigl((4\sqrt2+\sqrt6)\Sigma_\delta+4\mathcal B_\delta\bigr)t^2,
\]
using $ty\leq t^2$, $y^2\leq t^2$; converting $t$ gives $K_{\mathrm{ng},C}$. The pair-norm
bound follows from $\norm{D^2H[\zeta,\zeta]}_\star^2=\norm{(D^2H)_m}^2+\tfrac12\norm{(D^2H)_C}_{\Frob}^2$.
Every constant is proportional to $\mathcal B_\delta$, $\Sigma_\delta$, or $\bar\Sigma_\delta$,
all of which vanish for Gaussian targets.
\end{proof}

\begin{corollary}[Non-Gaussian Taylor remainder]
\label{cor:taylor-ng}
For $a=a_\star+\xi\in\Ureg_\delta$, write $F(a)=J_\star\xi+N_{\mathrm G}(\xi)+R_{\mathrm{ng}}(\xi)$,
where $N_{\mathrm G}(\xi):=F_{\mathrm G}(a)-J_{\mathrm G}\xi$ (the exact Gaussian nonlinearity,
$J_{\mathrm G}=-\Id$) and $R_{\mathrm{ng}}(\xi):=H(a)-DH(a_\star)\xi$. Then
$\norm{R_{\mathrm{ng}}(\xi)}_\star\leq\tfrac12K_{\mathrm{ng}}(\delta)\norm{\xi}_\star^2$.
\end{corollary}

\begin{proof}
$R_{\mathrm{ng}}(\xi)=\int_0^1(1-s)D^2H(a_\star+s\xi)[\xi,\xi]\dd s$ (as $H(a_\star)=0$);
each segment point is in $\Hreg_\delta$, so \Cref{lem:Kng} applies pointwise.
\end{proof}

\subsection{Main continuous-time local theorem}

\begin{theorem}[Gaussian-core local convergence of the flow]
\label{thm:core-flow}
Let $0<\gamma_0\leq\gstar$ and $\delta>0$, and define
\begin{equation}
\label{eq:rho-core}
    \rho_{\mathrm{core}}(\delta):=2\gamma_0-2+2\ell_\delta-K_{\mathrm{ng}}(\delta)\,\mathfrak b(\delta)\sqrt\delta .
\end{equation}
If $\rho_{\mathrm{core}}(\delta)>0$ and $\Delta(a_0)\leq\delta$, then the flow
\eqref{eq:NG-flow} remains in $\Ureg_\delta$ and
\[
    D(a_t)=\norm{a_t-a_\star}_\star^2\leq e^{-\rho_{\mathrm{core}}(\delta)\,t}\,D(a_0),\qquad t\geq0 .
\]
\end{theorem}

\begin{proof}
$\Ureg_\delta$ is forward invariant and compact (\Cref{lem:region-geom}), so the flow
exists for all $t\geq0$ in $\Ureg_\delta$; write $\xi_t=a_t-a_\star$. Since $\norm{\cdot}_\star$
is a fixed inner-product norm,
\[
    \tfrac12\tfrac{\dd}{\dd t}D(a_t)
    =\inner{\xi_t}{F(a_t)}_\star
    =\inner{\xi_t}{F_{\mathrm G}(a_t)}_\star+\inner{\xi_t}{H(a_t)}_\star .
\]
By \Cref{lem:gauss-core}(3), $\inner{\xi_t}{F_{\mathrm G}(a_t)}_\star\leq-\ell_\delta D(a_t)$.
Next, $H(a_t)=B_\star\xi_t+R_{\mathrm{ng}}(\xi_t)$ with $B_\star=DH(a_\star)=J_\star-J_{\mathrm G}=J_\star+\Id$,
which is self-adjoint (both $J_\star$ and $J_{\mathrm G}=-\Id$ are, by
\Cref{prop:linearized,lem:gauss-core}). Hence
\[
    \inner{\xi_t}{B_\star\xi_t}_\star
    =\inner{\xi_t}{J_\star\xi_t}_\star+\norm{\xi_t}_\star^2
    \leq-\gamma_0 D(a_t)+D(a_t)=(1-\gamma_0)D(a_t),
\]
using $\inner{\xi}{J_\star\xi}_\star\leq-\gstar\norm{\xi}_\star^2\leq-\gamma_0\norm{\xi}_\star^2$.
Finally, by \Cref{cor:taylor-ng} and $\norm{\xi_t}_\star\leq\mathfrak b(\delta)\sqrt\delta$
(\Cref{lem:region-geom}),
\[
    \inner{\xi_t}{R_{\mathrm{ng}}(\xi_t)}_\star
    \leq\norm{\xi_t}_\star\cdot\tfrac12K_{\mathrm{ng}}(\delta)\norm{\xi_t}_\star^2
    \leq\tfrac12K_{\mathrm{ng}}(\delta)\mathfrak b(\delta)\sqrt\delta\,D(a_t).
\]
Adding the three bounds,
$\tfrac12\tfrac{\dd}{\dd t}D\leq[-\ell_\delta+(1-\gamma_0)+\tfrac12K_{\mathrm{ng}}\mathfrak b\sqrt\delta]D$,
i.e.\ $\tfrac{\dd}{\dd t}D\leq-\rho_{\mathrm{core}}(\delta)D$; Gr\"onwall gives the claim.
\end{proof}

\begin{corollary}[Universal certified level]
\label{cor:core-universal}
With $\gamma_0=\gund$, $\rho_{\mathrm{core}}(\delta)\to2\gund$ as $\delta\downarrow0$
(since $\ell_\delta\uparrow1$ and $\mathfrak b(\delta)\sqrt\delta\downarrow0$), so for every
$\rho<2\gund$ there is $\delta_{\mathrm c}(\rho)>0$ with $\rho_{\mathrm{core}}(\delta)\geq\rho$
on $(0,\delta_{\mathrm c}(\rho)]$. Each such level is admissible in \Cref{thm:core-flow}.
\end{corollary}

\begin{corollary}[Exact sharpness for Gaussian targets]
\label{cor:gauss-sharp}
Let the target be Gaussian. Then $\gstar=1$, $K_{\mathrm{ng}}\equiv0$, and \Cref{thm:core-flow}
holds with the \emph{exact} rate
\[
    \rho_{\mathrm{core}}(\delta)=2\ell_\delta .
\]
Consequently, for every rate $\rho\in(0,2)$, the largest energy sublevel on which rate
$\rho$ is certified is $\{\Delta\leq\Delta_{\mathrm G}^\sharp(\rho)\}$ with
\begin{equation}
\label{eq:gauss-sharp-thresh}
    \Delta_{\mathrm G}^\sharp(\rho)=\varphi\!\Bigl(\tfrac\rho2\Bigr)=\tfrac12\Bigl(\tfrac\rho2-1-\log\tfrac\rho2\Bigr),
\end{equation}
independent of dimension, and this threshold is attained (\Cref{prop:sharp}).
\end{corollary}

\begin{proof}
For a Gaussian target $F=F_{\mathrm G}$, $H\equiv0$, so $\gstar=1$ (\Cref{prop:linearized})
and $K_{\mathrm{ng}}\equiv0$ (\Cref{lem:Kng}); \eqref{eq:rho-core} gives $\rho_{\mathrm{core}}=2\cdot1-2+2\ell_\delta=2\ell_\delta$.
Rate $\rho$ is certified iff $2\ell_\delta\geq\rho$, i.e.\ $\ell_\delta\geq\rho/2$; since
$\ell_\delta$ is the lower root of $\varphi(c)=\delta$ (\Cref{cor:band}) and $\varphi$ is
decreasing on $(0,1)$, this is $\delta\leq\varphi(\rho/2)$.
\end{proof}

\subsection{Sharpness benchmark}

\begin{definition}[Instantaneous rate and maximal threshold]
\label{def:sharp}
For $a\neq a_\star$, the instantaneous squared-norm rate and the maximal rate-$\rho$
threshold are
\[
    q_{\mathrm c}(a):=\frac{-2\inner{a-a_\star}{F(a)}_\star}{\norm{a-a_\star}_\star^2},
    \qquad
    \sharpthr(\rho):=\inf\{\Delta(a):a\neq a_\star,\;q_{\mathrm c}(a)<\rho\}
\]
for $\rho>0$ (with $\inf\emptyset=+\infty$).
\end{definition}

\begin{proposition}[The benchmark is exact and dominates the certified level]
\label{prop:sharp}
Let $\rho>0$. (1) If $\Delta(a_0)<\sharpthr(\rho)$ then $D(a_t)\leq e^{-\rho t}D(a_0)$ for all
$t\geq0$. (2) For every $\delta>\sharpthr(\rho)$ there is $a_0$ with $\Delta(a_0)\leq\delta$
and instantaneous rate $<\rho$; so $\sharpthr(\rho)$ is the largest energy threshold
guaranteeing rate $\rho$. (3) For a Gaussian target, $\sharpthr(\rho)=\Delta_{\mathrm G}^\sharp(\rho)=\varphi(\rho/2)$
for $\rho\in(0,2)$, in every dimension, and the certified level of \Cref{cor:gauss-sharp}
attains it.
\end{proposition}

\begin{proof}
(1) By energy descent $\Delta(a_t)\leq\Delta(a_0)<\sharpthr(\rho)$, so $q_{\mathrm c}(a_t)\geq\rho$
(else $\Delta(a_t)\geq\sharpthr(\rho)$); thus $\tfrac{\dd}{\dd t}D=-q_{\mathrm c}(a_t)D\leq-\rho D$.
(2) By definition of the infimum there is $a_0\neq a_\star$ with $q_{\mathrm c}(a_0)<\rho$ and
$\Delta(a_0)<\delta$. (3) For a Gaussian target $q_{\mathrm c}(a)=q_{\mathrm G}(a)$, and by
\Cref{lem:gauss-core}(3) $q_{\mathrm G}(a)\geq2\lambda_{\min}(C)$ with equality when $m=0$
and $C-\Id$ is in the minimal eigendirection. If $q_{\mathrm c}(a)<\rho$ then
$\lambda_{\min}(C)<\rho/2$, so some eigenvalue $c_j<\rho/2$ and, by \eqref{eq:gauss-energy},
$\Delta_{\mathrm G}(a)\geq\varphi(c_j)>\varphi(\rho/2)$ (as $\varphi$ decreases on $(0,1)$);
hence $\sharpthr(\rho)\geq\varphi(\rho/2)$. Conversely, $m=0$,
$C=\diag(c,1,\dots,1)$ with $c\uparrow\rho/2$ has $q_{\mathrm G}=2c<\rho$ and
$\Delta_{\mathrm G}=\varphi(c)\to\varphi(\rho/2)$, giving $\sharpthr(\rho)\leq\varphi(\rho/2)$.
Equality with \Cref{cor:gauss-sharp} is \eqref{eq:gauss-sharp-thresh}.
\end{proof}

\begin{remark}[Rate versus region]
\label{rem:rate-region}
For Gaussian targets both the region and the rate are exact: the certified threshold equals
$\sharpthr(\rho)$ and the certified rate equals the instantaneous rate at the active
constraint. The universal linearized bound $\gund=1/(4+\Gamma)$ is used only for general
targets and only as a lower bound on $\gstar$; when the exact gap $\gstar$ is available
(e.g.\ $\gstar=1$ for Gaussian, or computed from the linearization) it should be used in
\eqref{eq:rho-core}. As $\delta\downarrow0$, $\rho_{\mathrm{core}}(\delta)\to2\gstar$, the
exact asymptotic squared-norm rate.
\end{remark}

\subsection{Deterministic KL/Bregman scheme: exact Gaussian interval}
\label{subsec:kl-local}

The deterministic KL update \eqref{eq:KL-det} in whitened coordinates reads
$m_{n+1}=m_n+\Delta t\,C_ng(a_n)$, $C_{n+1}=(1+\Delta t)(C_n^{-1}+\Delta t\,A(a_n))^{-1}$,
and we measure it in $D_{\Delta t}(a):=\norm{a-a_\star}_{\star,\Delta t}^2=\norm{m}^2+\tfrac{1+\Delta t}2\norm{C-\Id}_{\Frob}^2$.

\begin{proposition}[Exact Gaussian KL map and sharp interval]
\label{prop:kl-gauss}
Let the target be Gaussian, so $g(a)=-m$, $A(a)=\Id$, and \eqref{eq:KL-det} becomes
\[
    m_{n+1}=(\Id-\Delta t\,C_n)m_n,
    \qquad
    C_{n+1}=(1+\Delta t)\,C_n(\Id+\Delta t\,C_n)^{-1}.
\]
In the eigenbasis of $C_n$, each covariance eigenvalue $c$ and the aligned mean component
transform by the scalar factors
\[
    r_m(c)=(1-\Delta t\,c)^2,
    \qquad
    r_C(c)=\frac1{(1+\Delta t\,c)^2},
    \qquad
    c^+=\frac{(1+\Delta t)c}{1+\Delta t\,c}.
\]
Fix $0<\Delta t<2$ and $q\in(0,1)$ with $q\geq q_0(\Delta t):=\max\{(1+\Delta t)^{-2},(1-\Delta t)^2\}$,
and define the interval
\begin{equation}
\label{eq:kl-interval}
    I_{\Delta t}^{\mathrm{KL}}(q):=\Bigl[\tfrac{q^{-1/2}-1}{\Delta t},\;\tfrac{1+\sqrt q}{\Delta t}\Bigr]\ni1 .
\end{equation}
Then: (i) $r_m(c)\leq q$ and $r_C(c)\leq q$ for all $c\in I_{\Delta t}^{\mathrm{KL}}(q)$;
(ii) the scalar map $c\mapsto c^+$ maps $I_{\Delta t}^{\mathrm{KL}}(q)$ into itself; hence
(iii) if $\operatorname{spec}(C_n)\subseteq I_{\Delta t}^{\mathrm{KL}}(q)$ then
$\operatorname{spec}(C_{n+k})\subseteq I_{\Delta t}^{\mathrm{KL}}(q)$ and
\[
    D_{\Delta t}(a_{n+k})\leq q^k\,D_{\Delta t}(a_n),\qquad k\geq0 .
\]
For the admissible stepsizes $\Delta t\leq\sqrt2$ (in particular $\Delta t\leq\tfrac2{4+\Gamma}$),
$q_0(\Delta t)=(1+\Delta t)^{-2}$. The energy sublevel
$\Delta_{\mathrm{G,KL}}^{\mathrm{band}}(q):=\min\{\varphi(\ell),\varphi(u)\}$, with
$[\ell,u]=I_{\Delta t}^{\mathrm{KL}}(q)$, is the largest sublevel \emph{contained in} this
invariant interval and is thus a sufficient---but generally conservative---entry certificate;
the exact contraction threshold is \Cref{prop:kl-sharp}.
\end{proposition}

\begin{proof}
The mean update $m_{n+1}=(\Id-\Delta t\,C_n)m_n$ and $C_{n+1}=(1+\Delta t)C_n(\Id+\Delta t\,C_n)^{-1}$
share the eigenbasis of $C_n$; in it, $c^+=\frac{(1+\Delta t)c}{1+\Delta t\,c}$ and the
aligned mean scales by $1-\Delta t\,c$, while $c^+-1=\frac{c-1}{1+\Delta t\,c}$ gives the
covariance factor $r_C(c)=(1+\Delta t\,c)^{-2}$. Thus
$D_{\Delta t}(a_{n+1})=\sum_i[r_m(c_i)(m_i)^2+\tfrac{1+\Delta t}2r_C(c_i)(c_i-1)^2]$.
(i) $r_m(c)\leq q\Leftrightarrow|1-\Delta t\,c|\leq\sqrt q\Leftrightarrow c\in[\tfrac{1-\sqrt q}{\Delta t},\tfrac{1+\sqrt q}{\Delta t}]$,
and $r_C(c)\leq q\Leftrightarrow1+\Delta t\,c\geq q^{-1/2}\Leftrightarrow c\geq\tfrac{q^{-1/2}-1}{\Delta t}$.
Since $q^{-1/2}-1\geq1-\sqrt q$ (equivalent to $(q^{-1/4}-q^{1/4})^2\geq0$), the covariance
lower bound dominates, giving \eqref{eq:kl-interval}; membership $1\in I_{\Delta t}^{\mathrm{KL}}(q)$
requires both $\tfrac{q^{-1/2}-1}{\Delta t}\leq1$ (i.e.\ $q\geq(1+\Delta t)^{-2}$) and
$1\leq\tfrac{1+\sqrt q}{\Delta t}$ (i.e.\ $\sqrt q\geq\Delta t-1$, automatic for
$\Delta t\leq1$ and equal to $q\geq(1-\Delta t)^2$ for $1<\Delta t<2$), so
$q\geq q_0(\Delta t)$. (ii) $c^+$ lies strictly between $c$ and $1$ (because
$c^+-1=\frac{c-1}{1+\Delta t\,c}$ has the sign of $c-1$ with smaller magnitude, and
$c^+-c=\frac{\Delta t\,c(1-c)}{1+\Delta t\,c}$ has the opposite sign), so an interval
containing $1$ is preserved. (iii) While $\operatorname{spec}(C_n)\subseteq I_{\Delta t}^{\mathrm{KL}}(q)$,
each scalar factor is $\leq q$, so $D_{\Delta t}$ contracts by $q$; forward invariance of the
interval propagates this for all $k$. Finally $\Delta_{\mathrm{G,KL}}^{\mathrm{band}}(q)$ is
the largest energy sublevel whose covariance band $[\ell_\delta,U_\delta]$ (\Cref{cor:band})
lies inside $[\ell,u]$, i.e.\ the largest $\delta$ with $\varphi(\ell)\geq\delta$ and
$\varphi(u)\geq\delta$; this is a sufficient certificate for the band condition of (i)--(iii),
not a claim of maximality over all contracting states.
\end{proof}

\begin{proposition}[Exact discrete contraction threshold]
\label{prop:kl-sharp}
For the Gaussian KL map $T_{\Delta t}$ and $a\neq a_\star$, let
$q_{\Delta t}(a):=D_{\Delta t}(T_{\Delta t}a)/D_{\Delta t}(a)$, and for $q\in(q_0(\Delta t),1)$
define the exact contraction threshold
\[
    \Delta_{\mathrm{G,KL}}^\sharp(q):=\inf\{\Delta_{\mathrm G}(a):a\neq a_\star,\;q_{\Delta t}(a)>q\}.
\]
Then, with $w:=\tfrac{1+\Delta t}2$, $r_m(c)=(1-\Delta t\,c)^2$, $r_C(c)=(1+\Delta t\,c)^{-2}$,
\begin{equation}
\label{eq:kl-sharp}
    \Delta_{\mathrm{G,KL}}^\sharp(q)=\inf_{c>0}\Bigl[\varphi(c)+\tfrac12z_{\min}(c;q)\Bigr],
    \quad
    z_{\min}(c;q)=
    \begin{cases}
    0, & r_C(c)>q,\\[2pt]
    w(c-1)^2\dfrac{q-r_C(c)}{r_m(c)-q}, & r_C(c)\leq q<r_m(c),\\[8pt]
    +\infty, & r_m(c)\leq q\ \text{and}\ r_C(c)\leq q .
    \end{cases}
\end{equation}
It satisfies $\Delta_{\mathrm{G,KL}}^{\mathrm{band}}(q)\leq\Delta_{\mathrm{G,KL}}^\sharp(q)$,
and $\Delta(a)<\Delta_{\mathrm{G,KL}}^\sharp(q)$ guarantees $q_{\Delta t}(a)\leq q$; it is the
largest open energy sublevel on which one KL step contracts $D_{\Delta t}$ by $q$ (the
infimum being over the strict-violation set $\{q_{\Delta t}(a)>q\}$). Equality
$\Delta_{\mathrm{G,KL}}^{\mathrm{band}}=\Delta_{\mathrm{G,KL}}^\sharp$ holds whenever
$\varphi(\ell)\leq\varphi(u)$ (the lower endpoint is the active constraint). The infimum in
\eqref{eq:kl-sharp} is one-dimensional and numerically computable.
\end{proposition}

\begin{proof}
In the eigenbasis of $C$, aligning the mean with covariance eigendirections (which only
lowers the energy needed to violate contraction), $D_{\Delta t}$, the map, and
$\Delta_{\mathrm G}=\tfrac12\norm m^2+\sum_i\varphi(c_i)$ all separate across coordinates.
Writing $D_{\Delta t}(T_{\Delta t}a)-qD_{\Delta t}(a)=\sum_i\bigl[(r_m(c_i)-q)m_i^2+w(r_C(c_i)-q)(c_i-1)^2\bigr]$,
the condition $q_{\Delta t}(a)>q$ is $\sum_i\psi_i>0$ with $\psi_i:=(r_m(c_i)-q)m_i^2+w(r_C(c_i)-q)(c_i-1)^2$.
Coordinates with $\psi_i\leq0$ only add energy without helping, so the cheapest violating
state perturbs a single coordinate: minimize $\varphi(c)+\tfrac12z$ over $c>0$, $z=m^2\geq0$
subject to $(r_m(c)-q)z+w(r_C(c)-q)(c-1)^2>0$. For fixed $c$: if $r_C(c)>q$ then $z=0$
already violates ($\psi>0$); if $r_C(c)\leq q<r_m(c)$ the constraint is
$z>w(c-1)^2\frac{q-r_C(c)}{r_m(c)-q}=:z_{\min}(c;q)$, attained in the limit; if both factors
are $\leq q$ then $\psi\leq0$ for all $z\geq0$, so $z_{\min}=+\infty$. Taking the infimum
over $c$ gives \eqref{eq:kl-sharp}. For the comparison
$\Delta_{\mathrm{G,KL}}^{\mathrm{band}}\leq\Delta_{\mathrm{G,KL}}^\sharp$: any state with
$q_{\Delta t}(a)>q$ has $\sum_i\psi_i>0$, so some coordinate has $\psi>0$, which forces
$c\notin[\ell,u]$ (both factors are $\leq q$ inside). If $c<\ell<1$ then
$\varphi(c)>\varphi(\ell)$; if $c>u>1$ then $\varphi(c)>\varphi(u)$; either way
$\varphi(c)+\tfrac12z\geq\varphi(c)>\min\{\varphi(\ell),\varphi(u)\}$, and taking the infimum
gives $\Delta_{\mathrm{G,KL}}^\sharp\geq\Delta_{\mathrm{G,KL}}^{\mathrm{band}}$. Conversely,
$m=0$ with $c\uparrow\ell$ from below has $r_C(c)>q$ (so $q_{\Delta t}>q$) and
$\Delta_{\mathrm G}(0,c)\to\varphi(\ell)$, giving $\Delta_{\mathrm{G,KL}}^\sharp\leq\varphi(\ell)$;
hence when $\varphi(\ell)\leq\varphi(u)$ the lower endpoint attains the infimum and
$\Delta_{\mathrm{G,KL}}^\sharp=\varphi(\ell)=\Delta_{\mathrm{G,KL}}^{\mathrm{band}}$. At the
upper endpoint, by contrast, $z_{\min}(c)\to+\infty$ as $c\downarrow u$ (there
$r_m(c)\downarrow q$ while $q-r_C(c)$ stays positive), so the upper band value $\varphi(u)$
can be strictly conservative when $\varphi(u)<\varphi(\ell)$.
\end{proof}

\begin{proposition}[General-target KL contraction via the Gaussian-core defect]
\label{prop:kl-general}
For a general target write $T_{\Delta t}^{\mathrm{KL}}$ for the map \eqref{eq:KL-det} and
$T_{\Delta t,\mathrm G}^{\mathrm{KL}}$ for its Gaussian core (\Cref{prop:kl-gauss}), and set
\[
    q_{\mathrm G}^{\mathrm{KL}}(\delta):=\max_{c\in[\ell_\delta,U_\delta]}\Bigl\{(1-\Delta t\,c)^2,\;(1+\Delta t\,c)^{-2}\Bigr\},
    \quad
    \mathfrak d_{\Delta t}^{\mathrm{KL}}(\delta):=\sup_{a\in\Ureg_\delta\setminus\{a_\star\}}
    \frac{\bigl[D_{\Delta t}(T_{\Delta t}^{\mathrm{KL}}a)-D_{\Delta t}(T_{\Delta t,\mathrm G}^{\mathrm{KL}}a)\bigr]_+}{D_{\Delta t}(a)} .
\]
Then for $a_n\in\Ureg_\delta$,
$D_{\Delta t}(a_{n+1})\leq\bigl(q_{\mathrm G}^{\mathrm{KL}}(\delta)+\mathfrak d_{\Delta t}^{\mathrm{KL}}(\delta)\bigr)D_{\Delta t}(a_n)$.
The defect is given exactly by the non-Gaussian discrepancies
$g(a)+m$ and $A(a)-\Id$: with $M_0:=C^{-1}+\Delta t\,\Id$, $M:=C^{-1}+\Delta t\,A(a)$, the map
differences are $e_m=\Delta t\,C(g(a)+m)$ and
$e_C=-(1+\Delta t)\Delta t\,M_0^{-1}(A(a)-\Id)M^{-1}$, and
\[
    D_{\Delta t}(T^{\mathrm{KL}}a)-D_{\Delta t}(T_{\mathrm G}^{\mathrm{KL}}a)
    =2\langle m_{\mathrm G}^+,e_m\rangle+\norm{e_m}^2
    +(1+\Delta t)\langle C_{\mathrm G}^+-\Id,e_C\rangle_{\Frob}+\tfrac{1+\Delta t}2\norm{e_C}_{\Frob}^2 .
\]
In particular $\mathfrak d_{\Delta t}^{\mathrm{KL}}(\delta)=0$ for every Gaussian target, so
\Cref{prop:kl-gauss} is the exact specialization. The defect is finite: both maps fix
$a_\star$ and $D_{\Delta t}(T^{\mathrm{KL}}a)-D_{\Delta t}(T_{\mathrm G}^{\mathrm{KL}}a)$ is
$O(\norm{a-a_\star}_\star^2)$ near $a_\star$ (each of $g(a)+m$ and $A(a)-\Id$ is
$O(\norm{a-a_\star}_\star)$ and vanishes at $a_\star$), so the ratio in
$\mathfrak d_{\Delta t}^{\mathrm{KL}}$ extends continuously at $a_\star$ and is bounded on the
compact $\Ureg_\delta$. Finally, every step initiated in $\Ureg_\delta$ contracts
$D_{\Delta t}$ by the displayed factor; if in addition
$\Delta t\leq1/(\beta_\star\bar\Lambda_\delta)$ (so that the KL one-step descent keeps
$\calE$ nonincreasing and hence $\Ureg_\delta$ forward invariant,
\Cref{lem:region-geom}(iv)) and
$q_{\mathrm G}^{\mathrm{KL}}(\delta)+\mathfrak d_{\Delta t}^{\mathrm{KL}}(\delta)<1$, the
contraction iterates geometrically on $\Ureg_\delta$.
\end{proposition}

\begin{proof}
$D_{\Delta t}(a_{n+1})=D_{\Delta t}(T^{\mathrm{KL}}a_n)\leq D_{\Delta t}(T_{\mathrm G}^{\mathrm{KL}}a_n)+\mathfrak d_{\Delta t}^{\mathrm{KL}}(\delta)D_{\Delta t}(a_n)$
by definition of the defect, and $D_{\Delta t}(T_{\mathrm G}^{\mathrm{KL}}a_n)\leq q_{\mathrm G}^{\mathrm{KL}}(\delta)D_{\Delta t}(a_n)$
by \Cref{prop:kl-gauss}(i) and the band $\operatorname{spec}(C_n)\subseteq[\ell_\delta,U_\delta]$
(\Cref{cor:band}). The mean difference is immediate from \eqref{eq:KL-det}. For the
covariance, the resolvent identity $M^{-1}-M_0^{-1}=-M_0^{-1}(M-M_0)M^{-1}=-\Delta t\,M_0^{-1}(A(a)-\Id)M^{-1}$
and $C^+=(1+\Delta t)M^{-1}$, $C_{\mathrm G}^+=(1+\Delta t)M_0^{-1}$ give $e_C$; the energy
difference is the polarization identity for $D_{\Delta t}$. For a Gaussian target
$g(a)+m=0$ and $A(a)-\Id=0$, so $e_m=e_C=0$ and $\mathfrak d_{\Delta t}^{\mathrm{KL}}\equiv0$.
\end{proof}

\begin{remark}[Precision coordinates]
\label{rem:precision}
The covariance update \eqref{eq:KL-det} is, in precision coordinates $P=C^{-1}$, the exact
scaled-Euler step $P_{n+1}=P_n+\tfrac{\Delta t}{1+\Delta t}(A(a_n)-P_n)$; for a Gaussian
target ($A\equiv\Id$) this is $P_{n+1}-\Id=\tfrac1{1+\Delta t}(P_n-\Id)$, a contraction by
$(1+\Delta t)^{-1}$ per step in every dimension and for every stepsize, with no nonlinear
defect. This is the precision-coordinate face of \Cref{prop:kl-gauss}.
\end{remark}

\subsection{Deterministic Riemannian scheme}
\label{subsec:riem-local}

\begin{proposition}[Exact Gaussian Riemannian map]
\label{prop:riem-gauss}
The Riemannian retraction update in whitened coordinates,
$m_{n+1}=m_n+\Delta t\,C_ng(a_n)$, $C_{n+1}=C_n^{1/2}\exp(\Delta t\,C_n^{-1/2}F_C(a_n)C_n^{-1/2})C_n^{1/2}$,
reduces for a Gaussian target to $m_{n+1}=(\Id-\Delta t\,C_n)m_n$ and, in the eigenbasis of
$C_n$, $c^+=c\,e^{\Delta t(1-c)}$. The scalar factors are $r_m^{\mathrm R}(c)=(1-\Delta t\,c)^2$
and
\[
    r_C^{\mathrm R}(c)=\Bigl(\frac{c\,e^{\Delta t(1-c)}-1}{c-1}\Bigr)^2,
    \qquad r_C^{\mathrm R}(1)=(1-\Delta t)^2 .
\]
For $0<\Delta t<2$ and $q\in(0,1)$ with $q\geq(1-\Delta t)^2$ (so that
$r_m^{\mathrm R}(1)=r_C^{\mathrm R}(1)=(1-\Delta t)^2\leq q$, i.e.\ $1$ belongs to the
defining set), let $I_{\Delta t}^{\mathrm R}(q)=[\ell,u]$ be the connected component
containing $1$ of $\{c>0:r_m^{\mathrm R}(c)\leq q,\ r_C^{\mathrm R}(c)\leq q\}$. If
additionally
\[
    \Delta t\leq\min\Bigl\{1,\;\tfrac{\log u}{u-1}\Bigr\}
\]
(the first part prevents a point below $1$ from overshooting above $1$, the second prevents
points in $(1,u]$ from crossing below $1$, so that $c\mapsto c\,e^{\Delta t(1-c)}$ keeps
$[\ell,u]$ forward invariant) and $\operatorname{spec}(C_n)\subseteq I_{\Delta t}^{\mathrm R}(q)$,
then $D(a_{n+k})\leq q^kD(a_n)$ for all $k\geq0$, and the general-target defect
$\mathfrak d_{\Delta t}^{\mathrm R}(\delta)$, defined as in \Cref{prop:kl-general}, vanishes
for Gaussian targets.
\end{proposition}

\begin{proof}
For a Gaussian target $F_C(a)=C-C^2=-C(C-\Id)$, so $C^{-1/2}F_CC^{-1/2}=\Id-C$ and the
retraction gives $C^{+}=C^{1/2}e^{\Delta t(\Id-C)}C^{1/2}=C\,e^{\Delta t(\Id-C)}$ (as $C$ and
$\Id-C$ commute); in the eigenbasis, $c^+=c\,e^{\Delta t(1-c)}$. Then
$c^+-1=c\,e^{\Delta t(1-c)}-1$, giving $r_C^{\mathrm R}$; the limit $c\to1$ is
$r_C^{\mathrm R}(1)=(1-\Delta t)^2$ by expansion of $c\,e^{\Delta t(1-c)}$ at $c=1$. The
mean factor is $r_m^{\mathrm R}(c)=(1-\Delta t\,c)^2$ as in the KL case. On
$I_{\Delta t}^{\mathrm R}(q)$ both factors are $\leq q$; the no-overshoot condition ensures
$c\mapsto c\,e^{\Delta t(1-c)}$ keeps the interval forward invariant, so the per-step
contraction propagates. Vanishing of the defect for Gaussian targets is as in
\Cref{prop:kl-general}.
\end{proof}

The Riemannian interval endpoints solve scalar transcendental equations; the KL interval
\eqref{eq:kl-interval} is closed-form, so we take \Cref{prop:kl-gauss} as the primary
discrete local theorem and \Cref{prop:riem-gauss} as its parallel.

\subsection{General-target discrete local rate}
\label{subsec:disc-general}

For general targets we record a linear-rate local theorem for both schemes; the Gaussian
specialization is sharper (the exact intervals above), but the following supplies the clean
per-step rate used in the complexity statements of \Cref{sec:three-stage}. Let
$\gamma_{\Riem}:=\gund$ and $\gamma_{\KL}(\Delta t):=\tfrac1{4+2\Delta t+\Gamma}$ denote the
certified linearized per-step gaps of the two scheme maps at $a_\star$ in the norms
$\norm{\cdot}_\star$ and $\norm{\cdot}_{\star,\Delta t}$ respectively, as proved in
\Cref{prop:linearized-maps}.  We use the shorthand
\[
    \mathfrak b_{\star,\Riem}(\delta):=\mathfrak b(\delta),
    \qquad
    \mathfrak b_{\star,\KL}(\delta):=\mathfrak b_{\Delta t}(\delta).
\]
For a scheme map $T$ write
$\Kdisc(\delta):=\Delta t^{-1}\sup_{a\in\Hreg_\delta}\norm{D^2T(a)}$, finite by
compactness of $\Hreg_\delta$ (\Cref{lem:region-geom}) and smoothness of $T$ on it, where
$\norm{D^2T(a)}:=\sup_{\xi,\zeta\neq0}\norm{D^2T(a)[\xi,\zeta]}_{\star,\bullet}/(\norm{\xi}_{\star,\bullet}\norm{\zeta}_{\star,\bullet})$
is the bilinear operator norm and the normalization by $\Delta t$ reflects
$\norm{D^2T}=O(\Delta t)$ (the scheme maps are $\Id+O(\Delta t)$).

\begin{theorem}[General-target discrete local convergence]
\label{thm:disc-local}
Fix a scheme $\bullet\in\{\Riem,\KL\}$ with gap $\gamma_\bullet$ and norm
$\norm{\cdot}_{\star,\bullet}$. Suppose
\[
    0<\Delta t\leq\min\Bigl\{\tfrac2{4+\Gamma},\,\tfrac1{\beta_\star\bar\Lambda_\delta}\Bigr\},
    \qquad
    \Kdisc(\delta)\,\mathfrak b_{\star,\bullet}(\delta)\sqrt\delta\leq\gamma_\bullet .
\]
If $\Delta(a_0)\leq\delta$, then the iterates remain in $\Ureg_\delta$ and
\[
    \norm{a_n-a_\star}_{\star,\bullet}\leq\bigl(1-\tfrac12\Delta t\,\gamma_\bullet\bigr)^n\norm{a_0-a_\star}_{\star,\bullet}.
\]
In particular the Riemannian rate is $(1-\tfrac{\Delta t}{2(4+\Gamma)})^n$ and the KL rate is
$(1-\tfrac{\Delta t}{2(4+2\Delta t+\Gamma)})^n$. Since $\Kdisc(\delta)\to\Kdisc(0)<\infty$ and
$\mathfrak b_{\star,\bullet}(\delta)\sqrt\delta\to0$ as $\delta\downarrow0$, the radius
condition holds for all sufficiently small $\delta$.
\end{theorem}

\begin{proof}
Invariance: \Cref{lem:Riem-descent,prop:KL-onestep} give
$\calE(a_{n+1})\leq\calE(a_n)$ under
$\Delta t\leq1/(\beta_\star\bar\Lambda_\delta)$, so $a_{n+1}\in\Ureg_\delta$ by
\Cref{lem:region-geom}(iv), and $\norm{\xi_n}_{\star,\bullet}\leq\mathfrak b_{\star,\bullet}(\delta)\sqrt\delta$.
Decompose $a_{n+1}-a_\star=M_\bullet\xi_n+\tfrac12 D^2T(\tilde a_n)[\xi_n,\xi_n]$ by
Taylor's theorem with $M_\bullet=DT(a_\star)$ and $\tilde a_n\in\Hreg_\delta$; the linear part
contracts by $1-\Delta t\gamma_\bullet$ by \Cref{prop:linearized-maps}.
The Taylor remainder satisfies
$\norm{\tfrac12 D^2T(\tilde a_n)[\xi_n,\xi_n]}_{\star,\bullet}\leq\tfrac12\norm{D^2T(\tilde a_n)}\norm{\xi_n}_{\star,\bullet}^2
=\tfrac12\Delta t\Kdisc(\delta)\norm{\xi_n}_{\star,\bullet}^2$
(the $\Delta t$ coming from $\norm{D^2T}=\Delta t\Kdisc$, not from Taylor's formula), which is
$\leq\tfrac12\Delta t\gamma_\bullet\norm{\xi_n}_{\star,\bullet}$ by the radius condition and
$\norm{\xi_n}_{\star,\bullet}\leq\mathfrak b_{\star,\bullet}(\delta)\sqrt\delta$. Adding gives the
per-step factor $1-\Delta t\gamma_\bullet+\tfrac12\Delta t\gamma_\bullet=1-\tfrac12\Delta t\gamma_\bullet$;
iterating completes the proof.
\end{proof}

\subsection{Online and a posteriori entry criteria}
\label{subsec:entry}

The energy test $\Delta(a_n)\leq\delta$ is not directly computable, since
$\Delta(a)=\calE(a)-\calE(a_\star)$ involves the unknown optimal value $\calE(a_\star)$. The
following sufficient tests use only the current iterate and known target quantities.

The local analysis above is carried out in equilibrium-whitened coordinates, but the
algorithm runs in the original coordinates and $C_\star$ is unknown online. The entry tests
must therefore be stated in original coordinates, where all quantities are available.

\begin{proposition}[Computable residual gate, original coordinates]
\label{prop:entry-residual}
Let $\alpha>0$ be the strong-convexity constant of the target ($\nabla^2V\succeq\alpha\Id$),
and at $a=(m,C)$ in original coordinates set $g(a)=-\E_{\rho_a}[\nabla V(X)]$,
$A(a)=\E_{\rho_a}[\nabla^2V(X)]$, and the Bures--Wasserstein gradient residual
\[
    \mathcal R_{\mathrm{BW}}(a)^2:=\norm{g(a)}^2+\Tr\!\bigl[(C^{-1}-A(a))\,C\,(C^{-1}-A(a))\bigr].
\]
Then $\mathcal R_{\mathrm{BW}}(a)^2\geq2\alpha\,\Delta(a)$, so for any a priori certified local
level $\delta_{\mathrm{loc}}$,
\[
    \mathcal R_{\mathrm{BW}}(a_n)^2\leq2\alpha\,\delta_{\mathrm{loc}}
    \quad\Longrightarrow\quad
    \Delta(a_n)\leq\delta_{\mathrm{loc}} .
\]
The left-hand side is evaluated from $g(a_n)$, $A(a_n)$, $C_n$, and $\alpha$, none of which
requires $C_\star$, $m_\star$, or $\calE(a_\star)$. When a covariance floor $C_n\succeq\ell_n\Id$
is available, the Fisher--Rao gradient norm gives an alternative intrinsic test,
$\Gradnorm(a_n)^2\leq\alpha\ell_n\delta_{\mathrm{loc}}\Rightarrow\Delta(a_n)\leq\delta_{\mathrm{loc}}$,
convenient when $\Gradnorm(a_n)$ is already evaluated by the algorithm; by
\Cref{lem:FR-W2} it is more restrictive than the Bures--Wasserstein gate
($\{\Gradnorm^2\leq\alpha\ell_n\delta_{\mathrm{loc}}\}\subseteq\{\mathcal R_{\mathrm{BW}}^2\leq2\alpha\delta_{\mathrm{loc}}\}$),
so the Bures--Wasserstein residual remains the primary online certificate. Once either test
triggers, energy descent keeps $\Delta(a_k)\leq\delta_{\mathrm{loc}}$ for all $k\geq n$, so the
local theorems apply permanently.
\end{proposition}

\begin{proof}
$\mathcal R_{\mathrm{BW}}(a)^2=\norm{\grad^{W_2}\calE(a)}_{W_2}^2$ is the squared
Bures--Wasserstein gradient norm of $\calE$ at $\rho_a$, and the Bures--Wasserstein PL
inequality $\norm{\grad^{W_2}\calE(a)}_{W_2}^2\geq2\alpha\,\Delta(a)$ holds by
\Cref{lem:W2-PL}; the Fisher--Rao variant with the floor $\ell_n$ is \Cref{lem:FR-W2}. The
implications are immediate, and forward invariance of $\Ureg_{\delta_{\mathrm{loc}}}$ is
\Cref{lem:region-geom}(iv).
\end{proof}

\begin{remark}[Computability of the threshold]
\label{rem:threshold-computable}
The gate detects entry to $\Ureg_{\delta_{\mathrm{loc}}}$ without knowing $\calE(a_\star)$;
the level $\delta_{\mathrm{loc}}$ itself is a target-adaptive quantity
(through $\alpha_\star,\beta_\star,\Sigma_\delta,\mathcal B_\delta,K_{\mathrm{ng}}$). A fully
precomputable conservative choice follows from the a priori bounds
$\alpha_\star\geq\kappa^{-1}$, $\beta_\star\leq\kappa$,
$\Sigma_\delta\leq\sqrt{N_\theta}(\beta_\star-\alpha_\star)$,
$\mathcal B_\delta\leq\max\{\beta_\star-1,1-\alpha_\star\}$, which reduce the admissibility
condition of \Cref{thm:core-flow} to a scalar inequality in $\alpha,\beta,N_\theta$ alone,
yielding a computable lower bound on the admissible $\delta_{\mathrm{loc}}$. In the
deterministic setting $g(a_n)$ and $A(a_n)$ are exact expectations and the gate is an exact
test; under Monte Carlo evaluation it becomes an estimated gate and should be applied with a
confidence margin, so ``computable online'' here means exact-or-certified evaluation rather
than a single-sample criterion.
\end{remark}

\begin{proposition}[Exact Gaussian spectral gates]
\label{prop:entry-gauss}
Let the target be Gaussian with precision $H$ (so $C_\star=H^{-1}$), and put
$S_n:=H^{1/2}C_nH^{1/2}$ (computable from $C_n$ and $H$, without the target mean). Then:
\begin{enumerate}
\item \emph{(Flow, $0<\rho\leq2$.)} $\lambda_{\min}(S_n)\geq\rho/2$ certifies instantaneous
rate $\geq\rho$ for the flow, and is forward invariant: each covariance eigenvalue obeys
$\dot s=s(1-s)$ and moves monotonically toward $1$, so for $\rho\leq2$ (hence $\rho/2\leq1$)
an eigenvalue $\geq\rho/2$ stays $\geq\rho/2$. (For $\rho>2$ the threshold exceeds $1$ and is
not forward invariant, since eigenvalues above $1$ decrease toward $1$.)
\item \emph{(KL scheme.)} $\operatorname{spec}(S_n)\subseteq I_{\Delta t}^{\mathrm{KL}}(q)$
certifies per-step contraction $D_{\Delta t}(a_{n+1})\leq q\,D_{\Delta t}(a_n)$ and is forward
invariant (\Cref{prop:kl-gauss}).
\end{enumerate}
These gates can be substantially less restrictive than the residual gate of
\Cref{prop:entry-residual} because they impose no condition on the mean and depend only on
the target-relative covariance spectrum.
\end{proposition}

\begin{proof}
$S_n$ is the whitened covariance $\widehat C_n$, with the same eigenvalues as
$C_n^{1/2}HC_n^{1/2}$. (1) For the flow $q_{\mathrm c}=q_{\mathrm G}\geq2\lambda_{\min}(\widehat C)=2\lambda_{\min}(S_n)$
by \Cref{lem:gauss-core}(3); the stated forward invariance for $\rho\leq2$ follows from the
scalar dynamics $\dot s=s(1-s)$. (2) is \Cref{prop:kl-gauss}(iii) in whitened coordinates.
\end{proof}
\section{Deterministic three-stage complexity}
\label{sec:three-stage}

\begin{theorem}[Deterministic three-stage complexity]
\label{thm:det-three-stage}
Assume \Cref{assump:logconcave-smooth} and
$\lambda_{0,\min}\Id\preceq C_0\preceq\lambda_{0,\max}\Id$. Let $\delta_{\mathrm c}>0$ be an
energy level admissible for \Cref{thm:core-flow} with $\rho_{\mathrm{core}}(\delta_{\mathrm c})\geq\gund$,
and $\delta_{\Riem},\delta_{\KL}>0$ levels admissible for \Cref{thm:disc-local} for the two
schemes (with the stepsizes below). Such levels exist: by \Cref{cor:core-universal}
$\rho_{\mathrm{core}}(\delta)\to2\gund>\gund$ as $\delta\downarrow0$, and the discrete radius
conditions hold for all small $\delta$ (\Cref{thm:disc-local}). For sufficiently small
$\epsilon>0$ (specifically $\epsilon<\mathfrak b(\delta_\bullet)^2\delta_\bullet$ for the
relevant level):
\begin{enumerate}
\item Continuous-time flow \eqref{eq:NG-flow}:
\[
    T_{\mathrm c}(\epsilon)
    =
    O\!\left(
        \logp\frac1{\betastar\lambda_{0,\star}}
        +\kappastar\logp\frac{\Delta_0}{\delta_{\mathrm c}}
        +(4+\Gamma)\log\frac1\epsilon
    \right)
\]
suffices for $\norm{a_T-a_\star}_\star^2\leq\epsilon$.  The first two terms may be replaced by
the smaller of this affine-invariant expression and
$\logp(1/(\beta\lambda_0))+\kappa\logp(\Delta_0/\delta_{\mathrm c})$.
\item Riemannian scheme \eqref{eq:Riem-det}, with
$0<\Delta t\leq\min\{1/(2\beta\lambda_{\max}),\,2/(4+\Gamma),\,1/(\beta_\star\bar\Lambda_{\delta_{\Riem}})\}$:
\[
    N_{\Riem}(\epsilon)
    =
    O\!\left(
        \frac1{\Delta t}\logp\frac1{\beta\lambda_0}
        +
        \frac\kappa{\Delta t}\logp\frac{\Delta_0}{\delta_{\Riem}}
        +
        \frac{4+\Gamma}{\Delta t}\log\frac1\epsilon
    \right)
\]
suffices for $\norm{a_N-a_\star}_\star^2\leq\epsilon$.
\item KL scheme \eqref{eq:KL-det}, with
$0<\Delta t\leq\min\{1/(\beta\lambda_{\max}),\,2/(4+\Gamma),\,1/(\beta_\star\bar\Lambda_{\delta_{\KL}})\}$:
\[
    N_{\KL}(\epsilon)
    =
    O\!\left(
        \frac1{\Delta t}\logp\frac1{\beta\lambda_0}
        +
        \frac\kappa{\Delta t}\logp\frac{\Delta_0}{\delta_{\KL}}
        +
        \frac{4+2\Delta t+\Gamma}{\Delta t}\log\frac1\epsilon
    \right)
\]
suffices for $\norm{a_N-a_\star}_{\star,\Delta t}^2\leq\epsilon$.
\end{enumerate}
\end{theorem}

\begin{proof}
\emph{Continuous time.} By \Cref{thm:cont-global} with $\delta=\delta_{\mathrm c}$, after time
$O(\logp\tfrac1{\betastar\lambda_{0,\star}}+\kappastar\logp\tfrac{\Delta_0}{\delta_{\mathrm c}})$ one has
$\Delta(a_t)\leq\delta_{\mathrm c}$, i.e.\ $a_t\in\Ureg_{\delta_{\mathrm c}}$ by
\Cref{def:region}, and $\norm{a_t-a_\star}_\star\leq\mathfrak b(\delta_{\mathrm c})\sqrt{\delta_{\mathrm c}}$ by
\Cref{lem:region-geom}(2). \Cref{thm:core-flow}, restarted at this time, gives
\[
    \norm{a_{t+s}-a_\star}_\star^2
    \leq
    e^{-s/(4+\Gamma)}\,\mathfrak b(\delta_{\mathrm c})^2\delta_{\mathrm c},
\]
so $s=(4+\Gamma)\log\bigl(\mathfrak b(\delta_{\mathrm c})^2\delta_{\mathrm c}/\epsilon\bigr)$ further time
suffices, which is $O((4+\Gamma)\log\tfrac1\epsilon)$ for small $\epsilon$ (since
$\delta_{\mathrm c}$ is independent of $\epsilon$).

\emph{Riemannian.} By \Cref{thm:Riem-global} with $\delta=\delta_{\Riem}$, after
$O(\tfrac1{\Delta t}\logp\tfrac1{\beta\lambda_0}+\tfrac\kappa{\Delta t}\logp\tfrac{\Delta_0}{\delta_{\Riem}})$
iterations one has $\Delta(a_n)\leq\delta_{\Riem}$, i.e.\
$a_n\in\Ureg_{\delta_{\Riem}}$, and
$\norm{a_n-a_\star}_\star\leq \mathfrak b(\delta_{\Riem})\sqrt{\delta_{\Riem}}$. The stepsize satisfies the
hypotheses of \Cref{thm:disc-local} for $\delta=\delta_{\Riem}$, so
\[
    \norm{a_{n+k}-a_\star}_\star
    \leq
    \Bigl(1-\frac{\Delta t}{2(4+\Gamma)}\Bigr)^{k}\mathfrak b(\delta_{\Riem})\sqrt{\delta_{\Riem}},
\]
and $k=O(\tfrac{4+\Gamma}{\Delta t}\log\tfrac1\epsilon)$ further iterations suffice for
$\norm{a_{n+k}-a_\star}_\star^2\leq\epsilon$, using $-\log(1-x)\geq x$.

\emph{KL.} By \Cref{thm:KL-global} with $\delta=\delta_{\KL}$, after
$O(\tfrac1{\Delta t}\logp\tfrac1{\beta\lambda_0}+\tfrac\kappa{\Delta t}\logp\tfrac{\Delta_0}{\delta_{\KL}})$
iterations one has $\Delta(a_n)\leq\delta_{\KL}$, i.e.\
$a_n\in\Ureg_{\delta_{\KL}}$, and
$\norm{a_n-a_\star}_{\star,\Delta t}\leq \mathfrak b_{\Delta t}(\delta_{\KL})\sqrt{\delta_{\KL}}$ by
\Cref{lem:region-geom}(2). The stepsize satisfies the hypotheses of
\Cref{thm:disc-local} for $\delta=\delta_{\KL}$, so
\[
    \norm{a_{n+k}-a_\star}_{\star,\Delta t}
    \leq
    \Bigl(1-\frac{\Delta t}{2(4+2\Delta t+\Gamma)}\Bigr)^{k}
    \mathfrak b_{\Delta t}(\delta_{\KL})\sqrt{\delta_{\KL}},
\]
and $k=O(\tfrac{4+2\Delta t+\Gamma}{\Delta t}\log\tfrac1\epsilon)$ further iterations
suffice, again using $-\log(1-x)\geq x$.
\end{proof}

\section{Transport bootstrap: Wasserstein burn-in followed by Fisher--Rao tail}
\label{sec:transport-bootstrap}

The deterministic covariance bootstrap in the preceding sections shows that pure
Fisher--Rao/natural-gradient dynamics has only logarithmic dependence on the initial lower
covariance. This is already a substantial improvement over the frozen-bound analysis. However,
when $C_0$ is severely underdispersed, even the logarithmic term
$\logp(1/(\beta\lambda_0))$ is a genuine Fisher--Rao warm-up cost. This section records a
clean way to remove that term: use Wasserstein transport only to lift the covariance to the
curvature scale, and then switch permanently to Fisher--Rao.

The main point is not to run a Wasserstein--Fisher--Rao splitting at every iteration. Instead, we
separate the roles of the two geometries:
\[
    \boxed{
    \begin{aligned}
    &\text{Wasserstein transport lifts an underdispersed covariance;}\\
    &\text{Fisher--Rao performs the global-to-local and local descent.}
    \end{aligned}
    }
\]
This avoids an adaptive transport schedule and avoids the cost of re-evaluating a Fisher--Rao
step after a Wasserstein half-step at every iteration.

\subsection{Covariance-scale comparison}

Let
\[
    A(a):=\E_{\rho_a}[\nabla^2V(X)]
    =-\E_{\rho_a}[\nabla_\theta^2\log\rho_{\post}(X)],
    \qquad
    \alpha\Id\preceq A(a)\preceq\beta\Id.
\]
For the Fisher--Rao covariance equation,
\[
    \dot C=C-CAC,
\]
if $\mu_i(t)$ is an eigenvalue of $C_t$ with normalized eigenvector $u_i(t)$, then
\[
    \dot\mu_i(t)
    =\mu_i(t)-\mu_i(t)^2h_i(t)
    =\mu_i(t)(1-\mu_i(t)h_i(t)),
    \qquad
    h_i(t):=u_i(t)^\top A(a_t)u_i(t)\in[\alpha,\beta].
\]
Thus, in the underdispersed regime $\mu_i\ll1/\beta$,
\[
    \dot\mu_i(t)\approx \mu_i(t),
\]
so the Fisher--Rao covariance growth is multiplicative. Starting from
$\mu_i(0)=\lambda_0\ll1/\beta$, reaching the scale $1/\beta$ necessarily takes time of order
$\log(1/(\beta\lambda_0))$.

By contrast, the unit-strength Gaussian Bures--Wasserstein flow is
\begin{equation}
\label{eq:W-flow}
    \dot m_t=g(a_t),
    \qquad
    \dot C_t=2\Id-C_tA(a_t)-A(a_t)C_t.
\end{equation}
For the same eigenpair,
\[
    \dot\mu_i(t)=2(1-\mu_i(t)h_i(t)).
\]
Thus, when $\mu_i\ll1/\beta$,
\[
    \dot\mu_i(t)\approx2,
\]
so Wasserstein covariance growth is additive and scale-independent. This is the precise
reason transport is better suited to the covariance warm-up.

The same comparison also explains why transport should not be kept forever. In a flat direction
with curvature $h_i=\alpha$, growing from $1/\beta$ to the target scale $1/\alpha$ takes
Wasserstein time of order
\[
    \frac1\alpha-\frac1\beta=O\!\left(\frac1\alpha\right),
\]
whereas Fisher--Rao, whose growth is approximately multiplicative until the target scale, takes
only
\[
    O\!\left(\log\frac\beta\alpha\right)=O(\log\kappa).
\]
Thus the clean covariance-scale picture is
\[
    \boxed{
    \text{Wasserstein is best for }0\to O(1/\beta),
    \qquad
    \text{Fisher--Rao is best for shape calibration }1/\beta\to1/\alpha.
    }
\]
This comparison is a covariance-scale dominance statement, not a pointwise claim that one
energy-dissipation rate dominates the other at every state. It is moreover
\emph{coordinate-dependent}: unlike the Fisher--Rao natural-gradient dynamics, the
Bures--Wasserstein flow \eqref{eq:W-flow} is not affine invariant, so the additive scale
$1/\beta$ and the $O(1/\beta)$ warm-start time below are stated in the given
(non-whitened) coordinates and would change under an affine reparametrization of the
target.

\subsection{Continuous-time Wasserstein bootstrap}

\begin{lemma}[Wasserstein covariance bootstrap]
\label{lem:W-cov-bootstrap}
Assume \Cref{assump:logconcave-smooth} and
$\lambda_{0,\min}\Id\preceq C_0\preceq\lambda_{0,\max}\Id$. Let
$\lambda_0:=\lambda_{0,\min}$ and
$\lambda_{\max}:=\max\{\lambda_{0,\max},1/\alpha\}$. Along the Gaussian Wasserstein flow
\eqref{eq:W-flow},
\[
    C_t\succeq \ell_W(t)\Id,
    \qquad
    \ell_W(t):=\frac1\beta+\left(\lambda_0-\frac1\beta\right)e^{-2\beta t},
\]
while $C_t\preceq\lambda_{\max}\Id$. Consequently, for any fixed $c\in(0,1)$, if
\[
    T_W^{\mathrm c}(c)
    :=
    \begin{cases}
    0,
    & \lambda_0\ge c/\beta,\\[3pt]
    \dfrac1{2\beta}\log\dfrac{1-\beta\lambda_0}{1-c},
    & \lambda_0<c/\beta,
    \end{cases}
\]
then
\[
    C_{T_W^{\mathrm c}(c)}\succeq \frac c\beta\Id.
\]
In particular, for fixed $c$, the Wasserstein covariance warm-up time is bounded by a constant
multiple of $1/\beta$ and has no logarithmic dependence on $\lambda_0$.
\end{lemma}

\begin{proof}
Let $\mu_i(t)$ be an eigenvalue of $C_t$ with normalized eigenvector $u_i(t)$. Along
\eqref{eq:W-flow},
\[
    \dot\mu_i(t)
    =u_i(t)^\top(2\Id-C_tA(a_t)-A(a_t)C_t)u_i(t)
    =2\bigl(1-\mu_i(t)h_i(t)\bigr),
\]
where $h_i(t):=u_i(t)^\top A(a_t)u_i(t)\in[\alpha,\beta]$. Hence
\[
    \dot\mu_i(t)\ge 2(1-\beta\mu_i(t)).
\]
Let $\ell_W$ solve
\[
    \dot\ell_W(t)=2(1-\beta\ell_W(t)),
    \qquad
    \ell_W(0)=\lambda_0.
\]
The solution is
\[
    \ell_W(t)=\frac1\beta+\left(\lambda_0-\frac1\beta\right)e^{-2\beta t}.
\]
The scalar comparison argument gives $\mu_i(t)\ge\ell_W(t)$ for every $i$, and hence
$C_t\succeq\ell_W(t)\Id$.

For the upper bound, if $\mu_i(t)\ge1/\alpha$, then
\[
    \dot\mu_i(t)=2(1-\mu_i(t)h_i(t))\le 2(1-\alpha\mu_i(t))\le0.
\]
Therefore no eigenvalue can cross the barrier $1/\alpha$ from below, and any eigenvalue that
starts above $1/\alpha$ is pushed downward. Since $C_0\preceq\lambda_{0,\max}\Id$, this gives
$C_t\preceq\lambda_{\max}\Id$.

Finally, if $\lambda_0<c/\beta$, the condition $\ell_W(t)\ge c/\beta$ is equivalent to
\[
    e^{-2\beta t}\le \frac{1-c}{1-\beta\lambda_0},
\]
which gives the stated expression for $T_W^{\mathrm c}(c)$. If $\lambda_0\ge c/\beta$, no
Wasserstein warm-up time is needed.
\end{proof}

\begin{theorem}[Continuous transport-bootstrap natural gradient]
\label{thm:cont-W-to-FR}
Assume \Cref{assump:logconcave-smooth}, and let $\delta_{\mathrm c}>0$ be a level admissible
for \Cref{thm:core-flow} with $\rho_{\mathrm{core}}(\delta_{\mathrm c})\geq\gund$ (such levels
exist by \Cref{cor:core-universal}). Fix $c\in(0,1)$. Run the
Gaussian Wasserstein flow \eqref{eq:W-flow} until time $T_W^{\mathrm c}(c)$ from
\Cref{lem:W-cov-bootstrap}, and then switch permanently to the Fisher--Rao flow
\eqref{eq:NG-flow}. Then, for sufficiently small $\epsilon>0$,
\[
    T_{W\to\FR}^{\mathrm c}(\epsilon;c)
    \le
    T_W^{\mathrm c}(c)
    +\min\left\{
        \frac{\kappa}{c}\logp\frac{\Delta_0}{\delta_{\mathrm c}},
        \logp\frac{\kappa}{c\betastar}
        +\kappastar\logp\frac{\Delta_0}{\delta_{\mathrm c}}
    \right\}
    +(4+\Gamma)\log\frac{\mathfrak b(\delta_{\mathrm c})^2\delta_{\mathrm c}}{\epsilon}
\]
suffices for $\norm{a_T-a_\star}_\star^2\leq\epsilon$.
In particular, for any fixed $c\in(0,1)$,
\[
    \boxed{
    T_{W\to\FR}^{\mathrm c}(\epsilon;c)
    =
    O\!\left(
        \frac1\beta
        +\min\left\{
            \kappa\logp\frac{\Delta_0}{\delta_{\mathrm c}},
            \logp\frac{\kappa}{\betastar}
            +\kappastar\logp\frac{\Delta_0}{\delta_{\mathrm c}}
        \right\}
        +(4+\Gamma)\log\frac1\epsilon
    \right).
    }
\]
Thus the pure-Fisher--Rao covariance-burn-in term
$\logp(1/(\beta\lambda_0))$ is replaced by a Wasserstein warm-start time of order $1/\beta$.
\end{theorem}

\begin{proof}
During the Wasserstein stage, $\calE$ is decreasing because \eqref{eq:W-flow} is the
Bures--Wasserstein gradient flow of $\calE$ on the Gaussian manifold
\citep{lambert2022variational}. Hence, at the switching state $a_{\mathrm b}$,
\[
    \Delta(a_{\mathrm b})\le\Delta_0,
    \qquad
    C_{\mathrm b}\succeq \frac c\beta\Id.
\]
After the switch, the covariance floor persists: $c/\beta\leq1/\beta$, and along the
Fisher--Rao flow the lower envelope of \Cref{lem:cont-cov} started from the floor
$c/\beta$ is nondecreasing (the scalar envelope
$\ell(t)=\lambda e^t/(1+\beta\lambda(e^t-1))$ with $\lambda=c/\beta$ satisfies
$\ell(0)=c/\beta$ and $\ell'(t)=\lambda e^t(1-\beta\lambda)/(1+\beta\lambda(e^t-1))^2\geq0$),
so $C_t\succeq\tfrac c\beta\Id$ for the whole Fisher--Rao stage. Hence the Fisher--Rao/Wasserstein
comparison in \Cref{lem:FR-W2} gives, along the Fisher--Rao stage,
\[
    \Gradnorm(a_t)^2
    \ge
    \frac{c}{2\beta}\norm{\grad^{W_2}\calE(a_t)}_{W_2}^2
    \ge
    \frac{c}{2\beta}\cdot2\alpha\,\Delta(a_t)
    =
    \frac{c}{\kappa}\Delta(a_t).
\]
Since $\tfrac{\dd}{\dd t}\Delta(a_t)=-\Gradnorm(a_t)^2$ along Fisher--Rao
(\Cref{lem:metric-dissipation}),
\[
    \frac{\dd}{\dd t}\Delta(a_t)
    \le
    -\frac{c}{\kappa}\Delta(a_t).
\]
Therefore, after an additional time
\[
    \frac{\kappa}{c}\logp\frac{\Delta_0}{\delta_{\mathrm c}},
\]
we have $\Delta(a_t)\le\delta_{\mathrm c}$.  Alternatively, the switching covariance satisfies
\[
    \widehat C_{\mathrm b}
    =C_\star^{-1/2}C_{\mathrm b}C_\star^{-1/2}
    \succeq\frac c\beta C_\star^{-1}
    \succeq\frac c\kappa\Id,
\]
so \Cref{thm:cont-global}, restarted at $a_{\mathrm b}$, gives the affine-invariant tail bound
\[
    \logp\frac{\kappa}{c\betastar}
    +\kappastar\logp\frac{\Delta_0}{\delta_{\mathrm c}}.
\]
Taking the better of the two bounds yields $\Delta(a_t)\leq\delta_{\mathrm c}$, i.e.\
$a_t\in\Ureg_{\delta_{\mathrm c}}$ by
\Cref{def:region}, with $\norm{a_t-a_\star}_\star\leq\mathfrak b(\delta_{\mathrm c})\sqrt{\delta_{\mathrm c}}$ by
\Cref{lem:region-geom}(2). \Cref{thm:core-flow}, restarted at this time, then gives
\[
    \norm{a_{t+s}-a_\star}_\star^2
    \le
    e^{-s/(4+\Gamma)}\,\mathfrak b(\delta_{\mathrm c})^2\delta_{\mathrm c}.
\]
Taking $s=(4+\Gamma)\log(\mathfrak b(\delta_{\mathrm c})^2\delta_{\mathrm c}/\epsilon)$ proves the explicit bound. The big-$O$
form follows from \Cref{lem:W-cov-bootstrap}, which gives
$T_W^{\mathrm c}(c)\le (2\beta)^{-1}\log(1/(1-c))$ when $c$ is fixed.
\end{proof}

\subsection{One-step discrete transport bootstrap}

The discrete version is even cleaner. A single Bures--Wasserstein forward--backward step creates
a covariance floor of size equal to the transport stepsize.

Given $a=(m,C)$ and a transport stepsize $\eta>0$, define
\begin{equation}
\label{eq:W-FB-step}
    g:=g(a),
    \qquad
    H:=\E_{\rho_a}[\nabla_\theta^2\log\rho_{\post}(X)]=-A(a),
    \qquad
    M:=\Id+\eta H=\Id-\eta A(a),
\end{equation}
\[
    m^+:=m+\eta g,
    \qquad
    \widetilde C:=MCM,
\]
\begin{equation}
\label{eq:W-FB-cov}
    C^+
    :=
    \frac12
    \left(
        \widetilde C+2\eta\Id+\bigl[\widetilde C(\widetilde C+4\eta\Id)\bigr]^{1/2}
    \right).
\end{equation}
The matrix square root is well-defined because $\widetilde C$ commutes with
$\widetilde C+4\eta\Id$.

\begin{lemma}[One-step Wasserstein covariance floor]
\label{lem:one-step-W-bootstrap}
Assume \Cref{assump:logconcave-smooth}, $C\preceq\lambda_{\max}\Id$ with
$\lambda_{\max}\ge1/\alpha$, and $0<\eta\le1/\beta$. Then the Wasserstein step
\eqref{eq:W-FB-step}--\eqref{eq:W-FB-cov} satisfies
\[
    C^+\succeq\eta\Id,
    \qquad
    C^+\preceq\lambda_{\max}\Id.
\]
Moreover, the standard Bures--Wasserstein forward--backward descent inequality
(the JKO/forward--backward descent guarantee of \citet{diao2023forward}, valid under the
stepsize restriction $\eta\le1/\beta$ used here) gives
\[
    \calE(m^+,C^+)\le\calE(m,C).
\]
\end{lemma}

\begin{proof}
Let
\[
    \psi_\eta(x)
    :=
    \frac12\left(x+2\eta+\sqrt{x(x+4\eta)}\right),
    \qquad x\ge0.
\]
The covariance update \eqref{eq:W-FB-cov} applies $\psi_\eta$ to the eigenvalues of
$\widetilde C$. Since $\psi_\eta$ is increasing and $\psi_\eta(0)=\eta$, every eigenvalue of
$C^+$ is at least $\eta$. Thus $C^+\succeq\eta\Id$.

For the upper bound, since $0<\eta\le1/\beta$ and $\alpha\Id\preceq A(a)\preceq\beta\Id$,
\[
    0\preceq M=\Id-\eta A(a)\preceq(1-\eta\alpha)\Id.
\]
Therefore
\[
    \widetilde C=MCM\preceq (1-\eta\alpha)^2\lambda_{\max}\Id.
\]
Because $\lambda_{\max}\ge1/\alpha$, we have $\eta\alpha\ge\eta/\lambda_{\max}$, hence
\[
    (1-\eta\alpha)^2\lambda_{\max}
    \le
    \left(1-\frac\eta{\lambda_{\max}}\right)^2\lambda_{\max}.
\]
The scalar identity
\[
    \psi_\eta\!\left(\left(1-\frac\eta L\right)^2L\right)=L,
    \qquad 0\le\eta\le L,
\]
with $L=\lambda_{\max}$ gives $C^+\preceq\lambda_{\max}\Id$.

Finally, \eqref{eq:W-FB-step}--\eqref{eq:W-FB-cov} is exactly the Bures--Wasserstein
forward--backward step for the splitting $\calE=\calE_1+\calE_2$, with the entropy part treated
backward. Under $\eta\le1/\beta$, the standard forward--backward descent inequality yields
$\calE(m^+,C^+)\le\calE(m,C)$.
\end{proof}

\begin{theorem}[Discrete transport-bootstrap natural gradient]
\label{thm:disc-W-to-FR}
Assume \Cref{assump:logconcave-smooth} and
$\lambda_{0,\min}\Id\preceq C_0\preceq\lambda_{0,\max}\Id$, and let the radii and
local levels $\delta_{\Riem},\delta_{\KL}$ be as in \Cref{thm:det-three-stage}. Fix $c\in(0,1]$ and set
\[
    \eta:=\frac c\beta.
\]
If $\lambda_0<c/\beta$, first perform one Wasserstein forward--backward step
\eqref{eq:W-FB-step}--\eqref{eq:W-FB-cov}; otherwise skip this step. Then run either the
Riemannian scheme \eqref{eq:Riem-det} or the KL scheme \eqref{eq:KL-det} with Fisher--Rao
stepsize $\Delta t$.

For the Riemannian scheme, if
\[
    0<\Delta t\le\min\left\{\frac1{2\beta\lambda_{\max}},\,\frac2{4+\Gamma},\,
    \frac1{\beta_\star\bar\Lambda_{\delta_{\Riem}}}\right\},
\]
then, for sufficiently small $\epsilon>0$,
\[
    \boxed{
    N_{W\to\Riem}(\epsilon;c)
    =
    \mathbf 1_{\{\lambda_0<c/\beta\}}
    +
    O\!\left(
        \frac{\kappa}{c\Delta t}\logp\frac{\Delta_0}{\delta_{\Riem}}
        +
        \frac{4+\Gamma}{\Delta t}\log\frac1\epsilon
    \right)
    }
\]
suffices for $\norm{a_N-a_\star}_\star^2\leq\epsilon$. For the KL scheme, if
\[
    0<\Delta t\le\min\left\{\frac1{\beta\lambda_{\max}},\,\frac2{4+\Gamma},\,
    \frac1{\beta_\star\bar\Lambda_{\delta_{\KL}}}\right\},
\]
then
\[
    \boxed{
    N_{W\to\KL}(\epsilon;c)
    =
    \mathbf 1_{\{\lambda_0<c/\beta\}}
    +
    O\!\left(
        \frac{\kappa}{c\Delta t}\logp\frac{\Delta_0}{\delta_{\KL}}
        +
        \frac{4+2\Delta t+\Gamma}{\Delta t}\log\frac1\epsilon
    \right)
    }
\]
suffices for $\norm{a_N-a_\star}_{\star,\Delta t}^2\leq\epsilon$.
Thus, for fixed $c$, the discrete covariance-burn-in term
$\Delta t^{-1}\logp(1/(\beta\lambda_0))$ is replaced by one Wasserstein bootstrap step.
\end{theorem}

\begin{proof}
If the Wasserstein step is skipped, then $C_0\succeq(c/\beta)\Id$ already. If the step is taken,
\Cref{lem:one-step-W-bootstrap} gives
\[
    C_{\mathrm b}\succeq\frac c\beta\Id,
    \qquad
    C_{\mathrm b}\preceq\lambda_{\max}\Id,
    \qquad
    \Delta(a_{\mathrm b})\le\Delta_0.
\]
Here $a_{\mathrm b}$ denotes the post-bootstrap state. Hence the subsequent Fisher--Rao stage
starts with the fixed covariance floor $c/\beta$ and the same upper bound $\lambda_{\max}$;
moreover, the covariance-floor recursions of \Cref{lem:Riem-cov,lem:KL-cov}, started
from $L_0=\min\{c/\beta,1/(2\beta)\}\geq c/(2\beta)$, are nondecreasing in $n$ (their
$y$-variables are nonincreasing), so $C_n\succeq\tfrac{c}{2\beta}\Id$ for the whole
Fisher--Rao stage; to keep the constants below simple we use the floor $c/(2\beta)$ in
the form of the one-step contractions.

For the Riemannian scheme, \Cref{lem:Riem-descent} with
$K=\beta\lambda_{\max}$ and $\Delta t\leq1/(2K)$, combined with
\Cref{lem:FR-W2} at the floor $L_n\geq c/(2\beta)$, gives
\[
    \Delta_{n+1}
    \leq
    \left(1-\frac{\alpha c\Delta t}{4\beta}\right)\Delta_n
    =
    \left(1-\frac{c\Delta t}{4\kappa}\right)\Delta_n.
\] Thus reaching $\delta_{\Riem}$ from a gap at most
$\Delta_0$ requires $O((\kappa/(c\Delta t))\logp(\Delta_0/\delta_{\Riem}))$ Riemannian
iterations. Once $\Delta_n\le\delta_{\Riem}$, \Cref{def:region} places the iterate in
$\Ureg_{\delta_{\Riem}}$, and \Cref{thm:disc-local} gives an additional
$O(((4+\Gamma)/\Delta t)\log(1/\epsilon))$ iterations, exactly as in the proof of
\Cref{thm:det-three-stage}.

For the KL scheme, \Cref{prop:KL-onestep} at the floor
$L_n\geq c/(2\beta)$ gives
\[
    \Delta_{n+1}
    \leq
    \left(
        1-
        \frac{\alpha(c/(2\beta))\Delta t}
             {2(1+\Delta t\beta\lambda_{\max})^2}
    \right)\Delta_n
    \leq
    \left(1-\frac{c\Delta t}{16\kappa}\right)\Delta_n,
\]
because $\Delta t\beta\lambda_{\max}\leq1$. So the
global localization phase takes $O((\kappa/(c\Delta t))\logp(\Delta_0/\delta_{\KL}))$
iterations, and once $\Delta_n\le\delta_{\KL}$, \Cref{thm:disc-local} gives the stated
local term, exactly as in the proof of \Cref{thm:det-three-stage}.
\end{proof}

\begin{remark}[What the transport bootstrap does and does not solve]
\label{rem:transport-bootstrap-caveat}
The transport bootstrap removes the small-initial-covariance burn-in from the deterministic
global theory. It does not, by itself, improve the fixed-step global-localization factor
$\kappa/\Delta t$: the lower construction of \Cref{thm:disc-global-sharp} already starts at
$c_0=1/\kappa$, after the covariance has reached the curvature scale. It also does not improve the
explicit local discretization factor $(4+\Gamma)/\Delta t$ for the Riemannian scheme or
$(4+2\Delta t+\Gamma)/\Delta t$ for the KL scheme. Thus covariance warm-up, fixed-step global
localization, and local stiffness are three distinct issues. To improve the latter two factors one
needs additional adaptive relative-curvature control, a line search, or a locally implicit,
proximal, or exponential-stabilized Fisher--Rao discretization.
\end{remark}

\begin{remark}[Practical switching rule]
The theorem suggests a particularly simple deterministic implementation. Pick a constant
$c\in(0,1]$, for example $c=1/2$, perform one Wasserstein step with
$\eta=c/\beta$ if $\beta\lambda_{\min}(C_0)<c$, and then switch permanently to the preferred
Fisher--Rao discretization. Equivalently, one may monitor
\[
    r_n:=\beta\lambda_{\min}(C_n)
\]
and switch once $r_n\ge c$. With the forward--backward Wasserstein bootstrap, this is
guaranteed to happen after one step by \Cref{lem:one-step-W-bootstrap} (which yields
$C^+\succeq(c/\beta)\Id$). The switch into the local
phase cannot be based on the energy gap $\Delta(a_n)=\calE(a_n)-\calE(a_\star)$ directly,
since $\calE(a_\star)$ is unknown; instead one uses the computable residual-based
sufficient tests of \Cref{prop:entry-residual}, e.g.\ switching once
$\Gradnorm(a_n)^2\le\alpha\ell\,\delta_\bullet$ with $\ell$ the current covariance floor.
\end{remark}

\begin{thebibliography}{9}

\bibitem[Lambert et al.(2022)]{lambert2022variational}
Marc Lambert, Sinho Chewi, Francis Bach, Silv\`ere Bonnabel, and Philippe Rigollet.
\newblock Variational inference via Wasserstein gradient flows.
\newblock In \emph{Advances in Neural Information Processing Systems}, volume 35, 2022.

\bibitem[Diao et al.(2023)]{diao2023forward}
Michael Ziyang Diao, Krishna Balasubramanian, Sinho Chewi, and Adil Salim.
\newblock Forward-backward Gaussian variational inference via JKO in the Bures--Wasserstein space.
\newblock In \emph{Proceedings of the 40th International Conference on Machine Learning}, 2023.

\end{thebibliography}

\end{document}
